% concatenated from PwCJS_K_Arxiv.tex
\documentclass[reqno]{amsart}
\usepackage[top=2cm,bottom=2cm,left=3cm,right=3cm]{geometry}
%\usepackage[left=3cm,right=3cm,top=3cm]{geometry}
%\parindent 0pt	
\usepackage{amssymb,amsthm,amsmath,amsxtra,dsfont,appendix,stix}
\usepackage{amsaddr}
%\usepackage[nottoc]{tocbibind}   % add bibliography to table of contents
%\usepackage{graphicx,tikz,esint,pgfplots}
%\usepackage[notcite, notref]{showkeys}
\usepackage{xcolor}
\definecolor{darkred}{RGB}{139,0,0}
\usepackage{hyperref} 
\hypersetup{
	colorlinks=true, 	
	citecolor=blue, 
	linkcolor=darkred,
	urlcolor=darkred
}
%[pagebackref=true]





\usepackage[nameinlink]{cleveref}
\usepackage[disable]{todonotes}

\usepackage{xparse,mathtools}
\usepackage{enumitem}
\usepackage{cancel}


\usepackage{mathrsfs}
%\usepackage{MnSymbol,wasysym}

\usepackage{tcolorbox}


\newcommand{\blue}{\color{blue}}
\newcommand{\red}{\color{red}}
\newcommand{\cyan}{\color{cyan}}



\newcommand{\1}{\mathds{1}}
\renewcommand{\i}{\mathbf{i}}
\renewcommand{\a}{\mathrm{a}}
\renewcommand{\d}{\mathrm{d}}
\newcommand{\A}{\mathcal{A}}
\newcommand{\B}{\mathcal{B}}
\newcommand{\C}{\mathds{C}}
\newcommand{\E}{\mathds{E}}

\newcommand{\F}{\mathcal{F}}
\newcommand{\G}{\mathrm{G}}
\newcommand{\W}{\mathrm{W}}
\renewcommand{\H}{\mathrm{H}}
\renewcommand{\O}{\mathcal{O}}
\renewcommand{\L}{\mathfrak{L}}
\newcommand{\M}{\mathrm{M}}
\newcommand{\N}{\mathbb{N}}
\newcommand{\No}{\mathcal{N}}
\newcommand{\Q}{\mathcal{Q}}
\renewcommand{\P}{\mathds{P}}
\newcommand{\X}{\mathrm{X}}
\newcommand{\R}{\mathds{R}}
\newcommand{\T}{\mathds{T}}
\newcommand{\U}{\mathscr{U}}
\newcommand{\Ai}{\operatorname{Ai}}
\newcommand{\SAi}{\operatorname{SAi}}
\newcommand{\V}{\mathscr{V}}
\newcommand{\REF}{{\color{red}REF}}


\newcommand{\re}{\text{\upshape Re\,}}
\newcommand{\im}{\text{\upshape Im\,}}
\newcommand{\res}{\text{\upshape Res\,}}


\newcommand{\sgn}{\operatorname{sgn}}
\newcommand{\Z}{\mathbb{Z}}
\newcommand{\D}{\mathds{D}}

\newcommand{\dif}{\mathrm{d}}
\newcommand{\bpsi}{\boldsymbol{\psi}}
\newcommand{\hpsi}{\hat{\psi}}
\newcommand{\tphi}{\boldsymbol{\chi}}
\newcommand{\Exp}{\mathbb{E}}

% \usepackage{titlesec}
% \titleformat{\subsection}[runin]
% {\normalfont\large\bfseries}{\thesubsection}{1em}{}
% \titleformat{\subsubsection}[runin]
% {\normalfont\bfseries}{}{0em}{} %\thesubsubsection

% % Define missing control sequence
% \makeatletter
% \def\@secnumpunct{.}
% \makeatother

\numberwithin{equation}{section}

\newtheorem{theorem}{Theorem}[section]
\newtheorem{conjecture}[theorem]{Conjecture}
\newtheorem{proposition}[theorem]{Proposition}
\newtheorem{corollary}[theorem]{Corollary}
\newtheorem{lemma}[theorem]{Lemma}
\newtheorem{remark}[theorem]{Remark}

\theoremstyle{definition} 
\newtheorem{definition}{Definition}[section]
\newtheorem{assumption}[definition]{Assumptions}
\newtheorem{figuretext}{Figure}
\usepackage{subcaption}
\usepackage[numbers,sort]{natbib}

\renewcommand{\Re}{\operatorname{Re}}
\renewcommand{\Im}{\operatorname{Im}}

\renewcommand\o{
\mathchoice
{{\scriptstyle\mathcal{O}}}% \displaystyle
{{\scriptstyle\mathcal{O}}}% \textstyle
{{\scriptscriptstyle\mathcal{O}}}% \scriptstyle
{\scalebox{.7}{$\scriptscriptstyle\mathcal{O}$}}%\scriptscriptstyle
}


%%%%%%%%%%%%%%%%%%%%%%%%%%%%%%%%%%%%%%%%%%
%%% Probability Macros
%%%%%%%%%%%%%%%%%%%%%%%%%%%%%%%%%%%%%%%%%%


\newcommand{\filt}{\mathscr{F}}


\newcommand{\Var}{\operatorname{Var}}
\newcommand{\Cov}{\operatorname{Cov}}
\newcommand{\dtv}{d_{TV}}
\newcommand{\weakto}{\Rightarrow}
\newcommand{\lawof}{\mathscr{L}}
\newcommand{\Bernoulli}{\operatorname{Bernoulli}}
\newcommand{\Binomial}{\operatorname{Binom}}
\newcommand{\Beta}{\operatorname{Beta}}
\newcommand{\Poisson}{\operatorname{Poisson}}
\newcommand{\Normal}{\operatorname{Normal}}
\newcommand{\Exponential}{\operatorname{Exp}}
\newcommand{\Unif}{\operatorname{Uni}}
\newcommand{\As}{\operatorname{a.s.}}
\newcommand{\lawequals}{\overset{\mathrm{law}}{=}}
\DeclareDocumentCommand{\Pto} {o} {
\IfNoValueTF {#1}
{\overset{\P}{\longrightarrow}}
{ \xrightarrow[ #1 \to \infty]{\P }}
}
\DeclareDocumentCommand{\Asto} {o} {
\IfNoValueTF {#1}
{\overset{\operatorname{a.s.}}{\longrightarrow}}
{
\xrightarrow[ #1 \to \infty]{\operatorname{a.s.} }
% \underset{#1 \to \infty}{\overset{\operatorname{a.s.}}{\longrightarrow}}
}
}
\DeclareDocumentCommand{\Mgfto} {o} {
\IfNoValueTF {#1}
{\overset{\operatorname{mgf}}{\longrightarrow}}
{ \xrightarrow[ #1 \to \infty]{\operatorname{mgf} }}
}

\DeclareDocumentCommand{\Wkto} {o} {
\IfNoValueTF {#1}
{\overset{\rm law}{\longrightarrow}}
{ \xrightarrow[ #1 \to \infty]{\operatorname{law}}}
}

\DeclareDocumentCommand \LPto { O{1} }
{\overset{\operatorname{\LP^{#1}}}{\longrightarrow}}



\title[Partition functions of two-dimensional Coulomb gases with circular root- and jump-type singularities]{Partition functions of two-dimensional Coulomb gases with circular root- and jump-type singularities}

\author[K. Noda]{Kohei Noda}
\address{{\small Institut de Recherche en Math\'ematique et Physique, Universit\'e catholique de Louvain, Louvain-la-Neuve, B--1348, Belgium}}
\email{kohei.noda@uclouvain.be}

%\author[E. Paquette]{Elliot Paquette}
%\address{Department of Mathematics, McGill University, Montreal, QC, Canada}
%\email{elliot.paquette@mcgill.ca}




\begin{document}

\maketitle



\begin{abstract}
In this paper, we study the random polynomial $p_n(\rho):=\prod_{j=1}^n (|z_j|-\rho)$, where the points $\{z_j\}_{j=1}^n$ are the eigenvalue moduli of random normal matrices with a radially symmetric potential. %, i.e., from a two-dimensional determinantal point process generalizing the Ginibre ensemble. 
We establish precise large $n$ asymptotic expansions for the moment generating function 
\[
\mathbb{E}\!\left[e^{\tfrac{u}{\pi}\im\log p_n(\rho)}\, e^{a\,\re\log p_n(\rho)}\right],  
\qquad u\in\R, \; a>-1,
\]
where $\rho>0$ lies in the bulk of the spectral droplet. The asymptotic expansion is expressed in terms of parabolic cylinder functions, which confirms a conjecture of Byun and Charlier. This also provides the first free energy expansion of two-dimensional Coulomb gases with general circular root- and jump-type singularities. 
While the $a=0$ case has already been widely studied in the literature due to its relation to counting statistics, we also obtain new results for this special case.   
\end{abstract}
\noindent
{\small{\sc AMS Subject Classification (2020)}: 41A60, 60B20, 60G55.}

\noindent
{\small{\sc Keywords}: Random normal matrices, counting statistics, partition functions, asymptotics}







\section{Introduction and statement of results}\label{section: introduction}
For $\boldsymbol{z}_n=(z_1,\dots,z_n)\in\C^n$, consider the $n$-fold integral 
%\todo[inline]{remove $\alpha$}
\begin{align}
\label{def of det partition function}
Z_{n,u,a}[Q]&:=\int_{\C^n}\prod_{1\leq j<k\leq n}|z_{k} -z_{j}|^{2}\prod_{j=1}^n |z_{j}|^{2\alpha}e^{-nQ(z_j)}\omega(z_j)\,\frac{d^2z_j}{\pi}, 
\end{align}
where $d^2z$ is the Lebesgue measure on $\C$, and $Q:\C\to\R$ is called the external potential.
Our assumptions on $Q$ are stated in Assumptions~\ref{Assumption_Q} below.
Here, $\omega(z)\equiv\omega(z;u,a)$ possesses a root singularity and a jump along the circle centered at 0 of radius $\rho>0$; more precisely, it is defined by 
\begin{equation}
\label{def of omega}
\omega(z):=|x-\rho|^{a}\begin{cases}
    e^{u}, & \mbox{if }x<\rho, \\
    1, & \mbox{if } x\geq \rho,
\end{cases}   
\quad x=|z|,\quad a>-1,\quad u\in\R. 
\end{equation}
Integrals of the form \eqref{def of det partition function} are typically called partition functions in the literature. They find applications in random matrix theory and statistical physics, and have therefore been widely studied in \cite{AFLS25,ACC2023c,Byun25,BFL25,BKS2023,BKSY2025,BSY2025,LS2017,WebbWong}. $Z_{n,\alpha,0,0}[Q]$ is also the normalization constant of the following probability measure:
\begin{align}\label{def of det point process}
d\mathbb{P}_{n}(\boldsymbol{z}_n):=\frac{1}{Z_{n,\alpha,0,0}[Q]} \prod_{1 \leq j < k \leq n} |z_{k} -z_{j}|^{2} \prod_{j=1}^{n}|z_{j}|^{2\alpha}e^{-n Q(z_{j})}\,\frac{d^2z_j}{\pi},\qquad \alpha>-1,
\end{align}
which represents the joint probability distribution of a random normal matrix \cite{BFreview,Mehta,Forrester}.
In particular, the choice $Q(z)=|z|^2$ corresponds to the complex Ginibre ensemble \cite{Ginibre}.
The measure \eqref{def of det point process} is a {\it determinantal point process} \cite{Forrester}. For a review of recent developments on non-Hermitian random matrices, see \cite{BFreview}.

The ratio $\frac{Z_{n,u,a}[Q]}{Z_{n,0,0}[Q]}$ gives the joint moment generating function of $(\re\log p_n(\rho),\im\log p_{n}(\rho))$, where $p_n$ is the random polynomial given by
\begin{equation}
p_n(x):=\prod_{j=1}^{n} \bigl(|z_j|-x\bigr).
\end{equation}
Note that the roots of $p_n$ are the moduli $\{|z_j|\}_{j=1}^n$, where $\{z_j\}_{j=1}^{n}$ are distributed according to \eqref{def of det point process}. In other words, if we denote by $\mathbb{E}$ the expectation with respect to \eqref{def of det point process}, then
\begin{equation}
\label{def of mathcal En}
 \mathcal{E}_{n,u,a} := \mathbb{E}\Bigl[e^{\frac{u}{\pi}\im \log p_n(\rho)}e^{a\re \log p_n(\rho)}\Bigr] = \frac{Z_{n,u,a}[Q]}{Z_{n,0,0}[Q]},
\end{equation}
with $u \in \R$, $a > -1$, $\rho >0$, and $\log p_n(\rho) = \log|p_n(\rho)| + \pi i \,\rm{N}_{\rho}$, where ${\rm N}_{\rho}:=\#\{z_j : |z_j| < \rho\}$. 
%Here, $\# A$ denotes the number of points in the finite set $A \subset \C$. 
The partition function $Z_{n,u,a}[Q]$ for $a\neq 0$ was first introduced in the work \cite{BC2022} of Byun and Charlier, where the potential $Q(z)$ in \eqref{def of det point process} was chosen to be the Mittag–Leffler potential $Q(z) = |z|^{2b}$ with $b > 0$. 
In \cite[Theorem 1.1]{BC2022}, they established the large $n$ asymptotic expansion of $\mathcal{E}_{n,u,a}$ up to and including the term of order 1, with a precise estimate for the error term. However, only the case $a\in\N$ was considered. 
An interesting phenomenon observed in their result is the appearance of associated Hermite polynomials in the coefficients of the terms of order $\sqrt{n}$ and $1$. The large $n$ asymptotics of $\mathcal{E}_{n,u,a}$ in the general case $a>-1$ was left open in \cite{BC2022}, and was conjectured to involve parabolic cylinder functions in place of the associated Hermite polynomials, see \cite[Remark 1.3]{BC2022}. 

\textit{In this paper, we obtain precise large $n$ asymptotics of $\mathcal{E}_{n,u,a}$ for general $u\in\R$, $a>-1$ and rotation-invariant $Q$, in the regime where $\rho$ lies in the bulk.} In particular, we confirm the conjecture from \cite{BC2022} that the asymptotics involve parabolic cylinder functions. This type of asymptotic behavior is completely new in random matrix theory, to the best of our knowledge. The case $a=0$ of our main results is already of interest, as it generalizes previous results on counting statistics of random normal matrix eigenvalues for general rotation-invariant potentials, as explained below.

For $a=0$, \eqref{def of mathcal En} reduces to the moment generating function of the disk counting statistics of random normal matrix eigenvalues, i.e.,
\begin{equation}
\label{def of disk counting statistics}
    \mathcal{E}_{n,u,0}=\mathbb{E}\bigl[e^{u\,{\rm N}_{\rho}}\bigr].
\end{equation}
Counting statistics of random normal matrix eigenvalues have attracted considerable attention in recent years, see e.g. \cite{ABE2023,ABES2023,LeeRiser2016,LMS2018,L et al 2019,FenzlLambert,SDMS2020,SDMS2021,ADM24,C2021 FH,ACCL1,ACCL2,CL2023}. 
In the work \cite{C2021 FH}, Charlier established the precise large $n$ asymptotic expansion of $\mathcal{E}_{n,u,0}$ for the Mittag–Leffler potential $Q(z)=|z|^{2b}$ with $b>0$. The asymptotics of the \textit{joint} moment generating functions, in the critical regime where all disk boundaries are merging at speed $n^{-1/2}$, were then obtained in the follow-up work \cite{CL2023}. In \cite{ACCL1,ACCL2}, Ameur, Charlier, Cronvall, and Lenells then treated the more difficult hard-edge regime where all disk boundaries are merging at speed $n^{-1}$ near a hard wall. In \cite{ACCL2}, the asymptotics contain an oscillatory term due to the fact that the particles accumulate on several components. The work \cite{ACM2024} also treats a hard-edge case, but in a simpler situation where there is no bulk. Leading order asymptotics of $\mathcal{E}_{n,u,0}$ were then obtained in \cite{ABES2023} for general rotation invariant potentials, and leading order asymptotics of $\mathrm{Var}[{\rm N}_{\rho}]$ in \cite{MMO25} for general potentials and domains. Our main result, in the special case $a=0$, improves on \cite{ABES2023} by providing the next two terms in the large $n$ asymptotics of $\mathcal{E}_{n,u,0}$. Since counting statistics has attracted considerable attention in recent years, for the convenience of the reader this particular case is stated separately in Theorem~\ref{theorem:asymptotic expansion of counting statistics} below. In Corollary~\ref{corollary:cumulant of counting statistics}, we also provide precise large $n$ asymptotics of all cumulants of $\mathrm{N}_{\rho}$ (not just the variance).


We will make the following assumptions on $Q$:
% from \cite{BKS2023,ABE2023,ABES2023,Seo} throughout this paper:
\begin{assumption}\label{Assumption_Q}
We suppose that the potential is rotation invariant, i.e., $Q(z)=q(|z|)$, and  satisfies the following conditions. 
\begin{enumerate}
\item[\textup{(1)}]\label{Assumption 1}
$\liminf_{|z|\to\infty}\frac{Q(z)}{2\log|z|}>1$, which guarantees that $Z_{n,u,a}[Q]<+\infty$. 
\item[\textup{(2)}] \label{Assumption 2}
$Q$ is $C^{6}$-smooth in a neighborhood of the droplet, subharmonic in $\C$, and strictly subharmonic (i.e., $\Delta Q(z)>0$) in a neighborhood of the droplet (the droplet is defined in \eqref{def of droplet} below).
\item[\textup{(3)}] $q'(0)>0$. 
\end{enumerate}
\end{assumption}
\begin{remark}
\label{remark:MF Potential}
Assumptions~\ref{Assumption_Q} cover a wide class of rotation invariant potentials.
Note however that it does not cover $Q(z)=|z|^{2b}$ for $b\neq 1$ as in this case $\lim_{r\to0}\Delta Q(r)$ either vanishes \textup{(}$b>1$\textup{)} or blows up \textup{(}$b\in(0,1)$\textup{)}.
\end{remark}
Under Assumptions~\ref{Assumption_Q}, the empirical measure $\frac1n\sum_{j=1}^n\delta_{z_j}$ of \eqref{def of det point process} converges weakly to the measure $\sigma_Q$ given by 
\begin{equation}
\label{def of sigma Q}
d\sigma_Q := \Delta Q \cdot \mathbf{1}_{S}\,\frac{d^2z}{\pi},
\end{equation}
where $S \equiv S_{Q}$ is a compact subset of $\C$ called the {\it droplet}.
Under parts (1) and (2) of Assumptions~\ref{Assumption_Q}, $S$ is of the form 
\begin{equation}
\label{def of droplet}
S = \mathbb{A}_{r_0,r_1}:=\{z \in \C : r_0 \leq |z| \leq r_1\},
\end{equation}
where $r_1$ is the smallest solution of $r q'(r)=2$; see, e.g., \cite{BKS2023,SaTo}. 
Part (3) of Assumptions~\ref{Assumption_Q} implies that $r_0=0$, i.e., $S=\mathbb{D}_{r_1}:=\{z\in\C:|z|\leq r_1\}$. (The analysis done in this paper can be adapted to the case $r_0>0$, which is in fact simpler.) 
%Moreover, although Assumption~\ref{Assumption_Q} does not fully encompass the setting of \cite{BC2022}, we believe that our results below is sufficiently general; see also Remarks~\ref{remark:counting statistics} and~\ref{remark for theorem cal EN 1}.   

\medskip
We recall that the complementary error function is defined by
\begin{equation}
\mathrm{erfc}(t):=\frac{2}{\sqrt{\pi}}\int_t^{+\infty}e^{-x^2}\,dx, \qquad t\in\R,     
\end{equation}
see e.g., \cite[Eq. (7.2.1)]{NIST}. 
Following \cite[Eq. (1.6)]{C2021 FH}, we introduce 
\begin{equation}
\label{def of mathcalFts}
\mathcal{F}(t,s):=\log\Bigl(1+\frac{s-1}{2}\mathrm{erfc}(t)\Bigr),    \qquad
t\in  \R, \quad s \in \C\backslash (-\infty, 0],
\end{equation}
where the principal branch is chosen for the $\log$. 


Our first main result, which corresponds to the case $a=0$, generalizes the result in \cite[Proposition~2.13]{ABES2023} by going beyond the leading term.  
\begin{theorem}[Counting statistics]
\label{theorem:asymptotic expansion of counting statistics}
Let $\rho\in(0,r_1)$, $\alpha>-1$, $u\in\R$, and $a=0$. 
Under Assumptions~\ref{Assumption_Q}, there exists $\delta>0$ such that, as $n\to+\infty$, we have   
\begin{equation}
\label{def of calEnu0 asymptotics}
\log\mathcal{E}_{n,u,0}
=C_1(u)\,n+C_2(u)\,\sqrt{n}+C_3(u)+\mathcal{O}\Bigl(\frac{(\log n)^3}{n^{\frac{1}{12}}}\Bigr),
\end{equation}  
uniformly for $u\in\{z\in\C:|z-x|\leq\delta\}$, where 
\begin{align}
\label{def of C1u disk count}
C_1(u)&:=
u\int_{\mathbb{D}_{\rho}}d\sigma_Q(z), 
\\
\label{def of C2u disk count}
C_2(u)&:=
\rho\sqrt{2\Delta Q(\rho)}\int_{0}^{+\infty}
\Bigl(
\mathcal{F}(x,e^{u})
+
\mathcal{F}(x,e^{-u})
\Bigr)\,dx,
\\
\begin{split}
\label{def of C3u disk count}
C_3(u)&:=
-
\Bigl(\alpha+\frac{1}{2}\Bigr)u
+\frac{1}{6}\Bigl(2+\frac{\rho\,\partial_r\Delta Q(\rho)}{\Delta Q(\rho)}\Bigr)u
\\
&\quad 
+\frac{1}{3}\Bigl(2+\frac{\rho\partial_r\Delta Q(\rho)}{\Delta Q(\rho)}\Bigr)
\int_{0}^{+\infty}
x\bigl(
\mathcal{F}(x,e^{u})
-
\mathcal{F}(x,e^{-u})
\bigr)\,dx.  
\end{split}
\end{align}
\end{theorem}



\begin{remark}[Consistency with Theorem 1.1 in \cite{C2021 FH}]
\label{remark:counting statistics}
As mentioned in Remark~\ref{remark:MF Potential}, Assumptions~\ref{Assumption_Q} do not cover the case $Q(z)=|z|^{2b}$ for $b\neq1$. However, surprisingly, 
substituting $Q(z)=|z|^{2b}$ into Theorem~\ref{theorem:asymptotic expansion of counting statistics} recovers \cite[Theorem~1.1]{BC2022} for any $b>0$, see Appendix~\ref{section: appendix counting}. 
Moreover, $C_3(u)$ is given in a simpler form than $C_3$ in \cite[Theorem 1.1]{C2021 FH}.
\end{remark}


Recall that the cumulants $\{\kappa_j\}_{j\in\mathbb{N}_{>0}}$ of the random variable ${\rm N}_{\rho}$ (see \eqref{def of mathcal En} below) are defined through the expansion 
\[
\log \mathbb{E}\bigl[e^{u\,{\rm N}_{\rho}}\bigr]
=
\kappa_1u+\frac{\kappa_2u^2}{2!}+\frac{\kappa_3u^3}{3!}+\cdots,\qquad u\to 0,
\]
or equivalently by
\begin{equation}
\label{def of cumulant}
\kappa_j=\partial_u^j\log \mathbb{E}\bigl[e^{u\,{\rm N}_{\rho}}\bigr]\Bigr|_{u=0},\qquad j\in\N.
\end{equation} 
\begin{comment}
More generally, the joint cumulants of $(\re\log p_n(\rho),\im\log p_{n}(\rho))$ are defined by 
\begin{equation}
\kappa_{j_u,j_a}:=\partial_u^{j_u}\partial_a^{j_a}\log\mathcal{E}_{n,u,a},\qquad
j_u,j_a\in\N_{>0}. 
\end{equation}
\end{comment}
In particular, since $\mathbb{E}[e^{u\,{\rm N}_{\rho}}]$ is analytic for $u\in\C$ and positive for any $u\in\R$, and Theorem~\ref{theorem:asymptotic expansion of counting statistics} is valid uniformly for $u\in\{z\in\C:|z-x|\leq \delta\}$ for some $\delta>0$,  Cauchy's formula implies that for any $u\in\R$, $j\in\N$, we have 
\begin{align*}
&\quad
\partial_u^{j}\Bigl\{\log \mathbb{E}\bigl[e^{u\,{\rm N}_{\rho}}\bigr]-\Bigl(C_1(u)\,n+C_2(u)\,\sqrt{n}+C_3(u)\Bigr)\Bigr\}
\\
&=
\frac{j!}{2\pi i}\oint_{|\zeta-u|=\frac{\delta}{2}}
\frac{\log \mathbb{E}\bigl[e^{\zeta\,{\rm N}_{\rho}}\bigr]-\bigl(C_1(\zeta)\,n+C_2(\zeta)\,\sqrt{n}+C_3(\zeta)\bigr)}{(\zeta-u)^{j+1}}\,d\zeta
=\mathcal{O}\Bigl(\frac{(\log n)^3}{n^{\frac{1}{12}}}\Bigr),\qquad
n\to+\infty.
\end{align*}
The above equation shows that \eqref{def of calEnu0 asymptotics} can be differentiated with respect to $u$ any fixed number of times without increasing the error term.  
Therefore, Theorem~\ref{theorem:asymptotic expansion of counting statistics} combined with \cite[Proof of Corollary 1.6]{C2021 FH} provides the following.

\begin{corollary}
\label{corollary:cumulant of counting statistics}  
Let $\rho\in(0,r_1)$, $\alpha>-1$, $u\in\R$, and $a=0$. 
Under Assumptions~\ref{Assumption_Q}, as $n\to+\infty$, we have 
\begin{equation}
\kappa_j=\begin{cases}
C_1'(0)\,n+C_3'(0)+\mathcal{O}\bigl(\frac{(\log n)^3}{n^{\frac{1}{12}}}\bigr),    & \mbox{if } j=1, \\
\partial_u^jC_3(0)+\mathcal{O}\bigl(\frac{(\log n)^3}{n^{\frac{1}{12}}}\bigr),    & \mbox{if } \text{$j\neq 1$ is odd}, \\
\partial_u^jC_2(0)\,\sqrt{n}+\mathcal{O}\bigl(\frac{(\log n)^3}{n^{\frac{1}{12}}}\bigr),     & \mbox{if } \text{$j$ is even}.
\end{cases}    
\end{equation}
\end{corollary}


The CLT from \cite[Proposition 2.13]{ABES2023} can also be deduced from Theorem~\ref{theorem:asymptotic expansion of counting statistics}.

\medskip
We now turn our attention to \eqref{def of mathcal En} in the general case $a>-1$.
For any $x\in\R$, $u\in\C$, and $a>-1$, define 
\begin{equation}
\label{def of cal H au}
\mathcal{H}_{a,u}(x):=\frac{\Gamma(a+1)}{\sqrt{2\pi}}e^{-\frac{1}{4}x^2}g_{a,u}(x),
    \qquad g_{a,u}(x):=
e^{u}D_{-a-1}(x)+D_{-a-1}(-x),
\end{equation}
where $D_{-\nu}(x)$ is the parabolic cylinder function \cite[Eq. (12.5.1)]{NIST}, which is defined by  
\begin{equation}
\label{def of parabolic cylinder function}
    D_{-\nu}(z)=U(\nu-\tfrac{1}{2},z),\qquad \re\nu>0. 
\end{equation}
Here,  
\begin{equation}
    U(\nu,z):=\frac{e^{-\frac{1}{4}z^2}}{\Gamma(\nu+\frac{1}{2})}
    \int_{0}^{+\infty}e^{-zt}t^{\nu-\frac{1}{2}}e^{-\frac{1}{2}t^2}\,dt,\qquad \re \nu>-\frac{1}{2}. 
\end{equation}
For the relationship between \eqref{def of parabolic cylinder function} and the associated Hermite polynomials, see Appendix~\ref{section:appendix parabolic cylinder function}. 
Note that for any $a>-1$, the integrand of $D_{-a-1}(x)$ is positive for all $x\in\R$.  
Consequently, $\mathcal{H}_{a,u}(x)$ is positive for all $x\in\R$ whenever $a>-1$ and $u\in\R$. 
Therefore, it follows that the logarithm $\log \mathcal{H}_{a,u}(x)$ is well defined for all $x\in\R$ and $u\in\R$.

The following is the main result of this paper. 
\begin{theorem}[Counting statistics and root-type statistics]
\label{theorem:calEn expansion}
Let $\rho\in(0,r_1)$, $\alpha>-1$, $u\in\R$, and $a>-1$. 
Under Assumptions~\ref{Assumption_Q}, there exists $\delta>0$ such that as $n\to+\infty$, we have 
\begin{equation}
\label{def of expansion of calEnua}
\log\mathcal{E}_{n,u,a}=C_1(u,a)\,n+C_2(u,a)\,\sqrt{n}+C_3(u,a)+\mathcal{O}\Bigl(\frac{(\log n)^3}{n^{\frac{1}{12}}}\Bigr),
\end{equation}
uniformly for $u\in\{z\in\C:|z-x|\leq \delta\}$ and $a$ in compact subsets of $(-1,+\infty)$, where 
\begin{align}
\label{def of mathcal En constant C1}
C_1(u,a)&:=
\int_{\C}\log\omega(z)\,d\sigma_{Q}(z)
=
\int_{\mathbb{D}_{\rho}}\bigl(u+a\log(\rho-|z|)\bigr)d\sigma_Q(z)
+\int_{\mathbb{A}_{\rho,r_1}}a\log (|z|-\rho)d\sigma_Q(z), 
\\
\label{def of mathcal En constant C2}
C_2(u,a)&:=
\rho\sqrt{\Delta Q(\rho)}
\int_{-\infty}^{+\infty}
\bigg(\log\bigl[\mathcal{H}_{a,u}(x)\bigr]
-a\log|x|-u\mathbf{1}_{(-\infty,0)}(x)
\bigg)
\,dx,
\\
C_3(u,a)&:=
-\frac{a}{2}\log\Bigl(\frac{r_1}{\rho}-1\Bigr)
-\frac{a(a-1)}{4}
\frac{r_1}{r_1-\rho}
-
\frac{a}{4}
\int_{0}^{r_1}
\frac{1}{x-\rho}
\Bigl(
\frac{x\partial_x\Delta Q(x)}{\Delta Q(x)}
-
\frac{\rho\partial_r\Delta Q(\rho)}{\Delta Q(\rho)}
\Bigr)
\,dx
\nonumber
\\
&\quad
+
\frac{a}{4}\Bigl(
4\alpha+a+2-\frac{\rho\partial_r\Delta Q(\rho)}{\Delta Q(\rho)}
\Bigr)\log\Bigl(\frac{r_1}{\rho}-1\Bigr)
\label{def of mathcal En constant C3}
\\
&\quad
-
\Bigl(\alpha+\frac{1}{2}\Bigr)u
-
\frac{a}{12}\Bigl(1-\frac{\rho\partial_r\Delta Q(\rho)}{\Delta Q(\rho)}\Bigr)u
+\frac{1}{6}\Bigl(2+\frac{\rho\,\partial_r\Delta Q(\rho)}{\Delta Q(\rho)}\Bigr)u
\nonumber
\\
&\quad 
+\frac{1}{6}\Bigl(2+\frac{\rho\partial_r\Delta Q(\rho)}{\Delta Q(\rho)}\Bigr)
\int_{-\infty}^{+\infty}
\bigg[
x\Bigl(\log\bigl[\mathcal{H}_{a,u}(x)\bigr] 
- u\mathbf{1}_{(-\infty,0)}(x)
\Big)
-ax\log|x| 
-\frac{a(a-1)x}{2(x^2+1)}
\bigg]\,dx. 
\nonumber
\end{align}
\end{theorem}
\begin{remark}\label{remark for theorem cal EN 1}
Again, as explained in Remark~\ref{remark:MF Potential}, Assumptions~\ref{Assumption_Q} do not cover the case $Q(z)=|z|^{2b}$ if $b\neq1$.  
However, in a similar way as in Remark~\ref{remark:counting statistics}, substituting $Q(z)=|z|^{2b}$ into Theorem~\ref{theorem:calEn expansion} recovers \cite[Theorem~1.1]{BC2022} for any $b>0$, see Appendix~\ref{section:appendix parabolic cylinder function}.
\end{remark}
\begin{figure}[t]
	\begin{subfigure}{0.44\textwidth}
    	\begin{center}	
    		\includegraphics[width=\textwidth]{figure/C3aComparison1040160.png}
            \subcaption{$a\mapsto C_3$}
    	\end{center}
    \end{subfigure}	 \quad 
	\begin{subfigure}{0.44\textwidth}
	\begin{center}	
		\includegraphics[width=\textwidth]{figure/C3rhoComparison.png}
            \subcaption{$\rho\mapsto C_3$}
	\end{center}
\end{subfigure}	
	\caption{    
The plot (a) shows $a\mapsto C_3$ (black line) and its comparison $a\mapsto \log\mathcal{E}_{n,u,a}-(C_1\,n+C_2\,\sqrt{n})$, where $Q(z)=0.2|z|^2+0.2345|z|^3, \alpha=0.667, u = 1.56$, $\rho=0.71r_1$, $n=10$ (red, dotted line), $n = 40$ (blue, dot-dashed line), and $n = 160$ (purple, dashed line)
The plot (b) shows $\rho\mapsto C_3$ (black line) and its comparison $\rho\mapsto \log\mathcal{E}_{n,u,a}-(C_1\,n+C_2\,\sqrt{n})$, where $Q,\alpha,u$ are same as before, $a=1.25$, and $n=100$ (red, dotted line), $n = 300$ (blue, dot-dashed line), and $n = 600$ (purple, dashed line). %Since $C_3$ is logarithmically divergent as $\rho\to0_{+}$, the behavior around the origin for the plot (b) is unstable for small size $n$.    
} \label{Fig_Comparison}
\end{figure} 
\begin{comment}
One of typical applications of \eqref{def of mathcal En} is to establish the joint cumulants of $(\frac{1}{\pi}\im\log p_n(\rho),\re\log p_n(\rho))$ for a fixed $\rho\in(0,r_1)$. 
Thus, the joint cumulants of $(\frac{1}{\pi}\im\log p_n(\rho),\re\log p_n(\rho))$ for a fixed $\rho\in(0,r_1)$ are defined by 
\begin{equation}
\kappa_{j_{u},j_{a}}
:=\partial_{u}^{j_{u}}\partial_{a}^{j_{a}}\log \mathcal{E}_n\bigr|_{u=0,a=0}, \qquad j_{u},j_{a}\in\mathbb{N}. 
\end{equation}
We write 
\[
\sigma_{Q}[\mathbb{D}_{\rho}]:=\int_{\mathbb{D}_{\rho}}d\sigma_Q(z). 
\]
Then, we can deduce from Theorem~\ref{theorem:calEn expansion} the following results. 
\begin{corollary}
\begin{description}
\item[(Asymptotics for the cumulants of each statistics)]
In particular, we obtain 
\begin{align}
\label{def of asymptotics of expectation of counting}
\mathbb{E}\bigl[N(\rho)\bigr]&\sim \sigma_Q[\mathbb{D}_{\rho}]\,n, 
\\
\label{def of asymptotics of variance of counting}
\mathrm{Var}\bigl[N(\rho)\bigr]&\sim \frac{\rho\sqrt{\Delta Q(\rho)}}{\sqrt{\pi}}\,\sqrt{n},
\\
\mathbb{E}\bigl[\log|p_n(\rho)|\bigr]&\sim\int_{S}\log[||z|-\rho|]\,d\sigma_Q(z)\,n, 
\\
\label{def of asymptotics of variance of real part of characteristic polynomial}
\mathrm{Var}\bigl[\log|p_n(\rho)|\bigr]&\sim
\rho\sqrt{\Delta Q(\rho)}
\int_{-\infty}^{+\infty}
\partial_a^2\log\Bigl[\frac{\Gamma(a+1)}{\sqrt{2\pi}}e^{-\frac{1}{4}x^2}g_0(x)\Bigr]
\Bigr|_{a=0,u=0}
\,dx
\end{align}
as $n\to+\infty$. 
\item[(Asymptotics for the joint cumulants)] 
\item[(Gaussian fluctuations)]
\end{description}    
\end{corollary}
\begin{remark}
We emphasize that \eqref{def of asymptotics of expectation of counting} is a consequence of the law of large number, and \eqref{def of asymptotics of variance of counting} was already obtained in \cite[Theorem 1.2]{ABE2023} and \cite[Proposition 2.13]{ABES2023}. 
\eqref{def of asymptotics of variance of real part of characteristic polynomial} seems to be new to our knowledge since $\log |p_n(\rho)|$ is discontinuous function of modulus of eigenvalues distributed to \eqref{def of det point process}.  
\end{remark}
\begin{remark}
In the special case $a=0$, by \eqref{def of D0 Dminus1}, \eqref{def of cal H au} can be simplified to 
\begin{equation} 
\label{def of cal H au a=0}
\mathcal{H}_{0,u}(x)=
1+\frac{e^u-1}{2}\mathrm{erfc}\Bigl(\frac{x}{\sqrt{2}}\Bigr)
=e^u\Bigl(1+\frac{e^{-u}-1}{2}\mathrm{erfc}\Bigl(-\frac{x}{\sqrt{2}}\Bigr) \Bigr),
\end{equation}
where we have used the identity that $1-\frac{1}{2}\mathrm{erfc}(-\tfrac{x}{\sqrt{2}})=\frac{1}{2}\mathrm{erfc}(\tfrac{x}{\sqrt{2}})$. 
Then, by \eqref{def of cal H au a=0}, \eqref{def of mathcal En constant C2} for $a=0$ satisfies 
\begin{equation}
\label{def of C2 a0}
C_2(u,0)
=
\rho\sqrt{2\Delta Q(\rho)}\int_{0}^{+\infty}
\Bigl(
\mathcal{F}(x,e^{u})
+
\mathcal{F}(x,e^{-u})
\Bigr)\,dx,
\end{equation}
It is straightforward to see that $C_3(u,0)$ is reduced to \eqref{def of mathcal En constant C3} by \eqref{def of cal H au a=0}. 
\end{remark}
\end{comment}
As mentioned earlier, Theorem~\ref{theorem:calEn expansion} confirms the conjecture stated in \cite[Remark~1.3]{BC2022}. 
Moreover, we obtained $C_3(u,a)$ in a simpler form than $C_3$ in \cite[Theorem~1.1]{BC2022}.
However, we currently do not have a conformal theoretic or geometric interpretation for \eqref{def of mathcal En constant C1}, \eqref{def of mathcal En constant C2}, and \eqref{def of mathcal En constant C3} as in \cite{ZW2006,BKS2023}.
We also believe that the error estimate in \eqref{def of expansion of calEnua} is not optimal, see \cite[Theorem~1.1 and Remark~1.2]{BC2022}.

\medskip

We finally deduce the large $n$ asymptotics of \eqref{def of det partition function} by combining Theorem~\ref{theorem:calEn expansion} with \cite[Theorem~1.4]{ACC2023c}.

We define
\begin{equation}
I_{Q}[\mu]:=\int_{\C^2}\log \frac{1}{|z-w|}\,d\mu(z)d\mu(w)+\int_{\C}Q(z)\,d\mu(z),   
\end{equation}
which is referred to as the weighted logarithmic energy, defined over all compactly supported Borel probability measures $\mu$.
In particular, if $Q$ is lower semi-continuous and finite on a set of positive capacity, Frostman’s theorem guarantees the existence of a unique equilibrium measure $\sigma_Q$ minimizing the weighted logarithmic energy, see \cite{SaTo}.
In particular, for the potential $Q$ satisfying Assumptions~\ref{Assumption_Q}, we have
\[
I_{Q}[\sigma_Q]
=q(r_1)-\log r_1-\frac{1}{4}\int_{0}^{r_1}rq'(r)^2\,dr.
\]
We define 
\begin{align*}
E_{Q}[\sigma_{Q}]&:=\int_{\C}\log \Delta Q\,d\sigma_{Q}, 
\\
F_{Q}[\sigma_{Q}]&:=\frac{1}{12}\log\frac{1}{r_1^2\Delta Q(r_1)}-\frac{1}{16}\frac{r_1\partial_r\Delta Q(r_1)}{\Delta Q(r_1)}+\frac{1}{24}\int_0^{r_1}\Bigl(\frac{\partial_r\Delta Q(r)}{\Delta Q(r)}\Bigr)^2r\,dr, 
\end{align*}
and for $\ell_{\alpha}(z):=2\alpha\log |z|$, 
\begin{align*}
\mathsf{e}_{\ell_{\alpha}}
&:=
\frac{1}{2}\int_{S}\ell_{\alpha}(z)\Delta \log \Delta Q(z)\,dA(z)+\frac{1}{8\pi}\int_{\partial S}\partial_{\mathrm{n}}\ell_{\alpha}(z)|dz|-\frac{1}{8\pi}\int_{\partial S}\ell_{\alpha}(z)\frac{\partial_{\mathrm{n}}\Delta Q(z)}{\Delta Q(z)}\,|dz|, 
\end{align*}
where “$\partial_{\mathrm{n}}$” designates differentiation in the normal direction to $\partial S$ pointing out from the droplet $S$. 
Here, $E_{Q}[\sigma_Q]$ represents the negative entropy of the equilibrium measure, while $F_{Q}[\sigma_Q]$ can be interpreted in terms of $\zeta$-regularized determinants associated with certain pseudo-differential operators. For additional background and details, we refer the reader to \cite{BKS2023,ACC2023c} and the references therein. 

We conclude with the following result.
\begin{theorem}[Partition function with circular- and root-type singularities]
\label{corollary:partition function}
Let $\rho\in(0,r_1)$, $\alpha>-1$, $u\in\R$, and $a>-1$. 
Under Assumptions~\ref{Assumption_Q}, there exists $\delta>0$ as $n\to+\infty$, we have   
\[
Z_{n,a,u}[Q]
=
\widetilde{C}_1n^2+\widetilde{C}_2n\log n+\widetilde{C}_3n+\widetilde{C}_4\sqrt{n}+\widetilde{C}_5\log n+\widetilde{C}_6+\mathcal{O}\Bigl(\frac{(\log n)^3}{n^{\frac{1}{12}}}\Bigr), 
\]
uniformly for $u\in\{z\in\C:|z-x|\leq\delta\}$ and $a$ in compact subsets of $(-1,+\infty)$, where 
\begin{align*}
\widetilde{C}_1&:=-I_Q[\sigma_Q],
\qquad
\widetilde{C}_2:=\frac{1}{2}, 
\\
\widetilde{C}_3&:=\frac{\log 2\pi}{2}-1-\frac{E_Q[\mu_Q]}{2}+\int_{\C}\ell_{\alpha}(z)\,d\sigma_Q(z)+C_1(u,a),
\qquad
\widetilde{C}_4:=C_2(u,a),
\qquad
\widetilde{C}_5:=\frac{5}{12}+\frac{\alpha^2}{2}, 
\\
\widetilde{C}_6&:=\zeta'(-1)-\log G(1+\alpha)+F_Q[\mu_Q]+\frac{1+\alpha}{2}\log(2\pi)+\mathsf{e}_{\ell_{\alpha}}
    +\frac{\alpha^2}{2}\log\bigl(r_1^2\Delta Q(0)\bigr)+C_3(u,a)
\end{align*}
where $C_1(u,a),C_2(u,a),C_3(u,a)$ are given by \eqref{def of mathcal En constant C1}, \eqref{def of mathcal En constant C2}, \eqref{def of mathcal En constant C3}, respectively, and $\zeta(z)$ is the Riemann zeta function. 
\end{theorem}

\begin{remark}
The coefficient $\widetilde{C}_5$ can be rewritten as $\frac{6-\chi}{12}$,  where $\chi$ is the Euler characteristic of the droplet. In our case, since the droplet is a disk, $\chi=1$. The fact that $\widetilde{C}_5$ can be rewritten in this way is consistent with the general conjecture \cite{ZW2006}. 
Similar results have been established in \cite{BKS2023} for centered disks and annuli, in \cite{ACC2023c} for multiple disjoint connected annuli possibly including a central disk, and in \cite{BKSY2025} for the (generalized) spherical ensemble.
\end{remark}


\subsubsection*{Related works}
We conclude by briefly mentioning recent developments in the literature related to Theorem~\ref{theorem:asymptotic expansion of counting statistics}, \ref{theorem:calEn expansion}, and Corollary~\ref{corollary:partition function}.

The study of structured determinants with singularities in one-dimensional point processes has a long history. In the seminal work \cite{FisherHartwig}, Fisher and Hartwig conjectured the asymptotic behavior of Toeplitz determinants in the large-size limit when the weight function is supported on the unit circle and possesses root- and jump-type singularities, which is now commonly referred to as Fisher--Hartwig singularities. 
For early works and more historical background, we refer to \cite{DIKreview} and references therein. 
Recent results on structured determinants with singularities in connection with random matrix theory include \cite{DIK, DIK14, Fahs,CGMY2020,CK2015,BCL2024} for Toeplitz determinants, \cite{DXZ2022,DXZ2023,LYZ24,Xetal,XZZ2024,YZ24,CM2023} for Fredholm determinants, \cite{Krasovsky, ItsKrasovsky, BerWebbWong, Charlier, CFWW2021} for Hankel determinants, \cite{BasorEhrhardt1, DIK} for Toeplitz+Hankel determinants. The above list is not exhaustive, and we therefore refer the reader to the references cited therein.


There has been significant progress on the precise large-$n$ asymptotics of partition functions, moment generating functions, and hole probabilities for the random normal matrices mentioned earlier. 
A seminal development is the works of Charlier \cite{C2021 FH,C2021} on the moment generating function and hole probabilities for random normal matrices with Mittag–Leffler potentials $Q(z)=|z|^{2b}$ for $b>0$.
Further progress was then made by Byun, Kang, and Seo \cite{BKS2023}, who obtained precise large-$n$ asymptotics for the partition functions of two-dimensional Coulomb gases with rotation-invariant potentials. 
Building on these results, subsequent advances include the study of multi-component droplets \cite{ACC2023c}; hole probabilities \cite{C2023,BP24}; partition functions with a fixed hard wall \cite{AFLS25}; partition functions for the case $\Delta Q=0$ along some circle inside the droplet \cite{AL25}; partition functions of spherical Coulomb gases \cite{BKSY2025}; Coulomb gases with Lemniscate-type potentials having a conical singularity at the origin \cite{Byun25}; and precise large-$n$ asymptotics of the Ginibre ensemble with a large point charge insertion breaking rotational symmetry \cite{BSY2025}. The latter is motivated by the strong asymptotics of planar orthogonal polynomials associated with the Ginibre ensemble with a point insertion \cite{BBLM2015}, and is connected to further investigations \cite{BEG18,BKP2023,Byun,LeeYang,LeeYang2,LY2023,BLY2022,BY2023} as well as their applications to Gaussian multiplicative chaos \cite{WebbWong}.
For recent developments on non-Gaussian potentials, such as the truncated unitary ensemble and the spherical ensemble with a point insertion, see \cite{DKMS25,BFL25,BFKL25}. These models correspond to point singularities and require the Riemann–Hilbert problem approach. In contrast, the case \eqref{def of mathcal En} involves circular and jump-type singularities, but our analysis does not require the Riemann–Hilbert problem as in \cite{BC2022}.


\subsubsection*{Outline}
It suffices to prove Theorem~\ref{theorem:calEn expansion}, as Theorem~\ref{theorem:asymptotic expansion of counting statistics} and Corollary~\ref{corollary:partition function} follow as special cases.
The strategy for proving Theorem~\ref{theorem:calEn expansion} combines the Laplace method, the precise Riemann sum approximation of \cite{C2021}, and the decomposition of a large sum into global and local analysis parts as in \cite{BC2022}.

We first split the logarithmic sum in \eqref{def of mathcal En} into a global part and a local part. 
In Section~\ref{section:global analysis}, we establish the large $n$ asymptotics of the global part (Lemmas~\ref{lemma:S0}, \ref{lemma:S1}, and \ref{lemma:S3}), 
where the proofs of Lemmas~\ref{lemma:S1} and~\ref{lemma:S3} require careful analysis of the error. 
In Section~\ref{section:local analysis part}, we establish the large $n$ asymptotics of the summand corresponding to the local part. 
This is first stated in Lemma~\ref{lemma:S2in part1}, and subsequently, each sum is approximated using Lemma~\ref{lemma:Euler-Maclaurin} with more refined error estimates than in the global part.  

%We emphasize that both the global and local analysis parts (Sections~\ref{section:global analysis} and~\ref{section:local analysis part}) are carried out using the Euler--Maclaurin formula, Lemma~\ref{lemma:Euler-Maclaurin}. 
%Alternatively, they could also be derived using the more precise Riemann approximation techniques of \cite[Lemmas~2.9 and~4.10]{BC2022}.



\section{Global analysis part for the proof of Theorem~\ref{theorem:calEn expansion}}
\label{section:global analysis}
Since $|z|^{2\alpha}e^{-nq(|z|)}\omega(z)$ is rotation invariant, $\mathcal{E}_{n}\equiv\mathcal{E}_{n,u,a}$ can be identically expressed in terms of one-fold integrals.
This is well-known fact as {\it Andr\'{e}ief identity} and has already been used in different contexts, see e.g., \cite{BC2022,BKS2023,ACC2023c,ACCL1,ACCL2,ACM2024,BP24,C2021 FH}.
For fixed $u\in\R,a\in(-1,+\infty)$, and $\rho\in(0,r_1)$, we have  
\begin{equation}
\label{def of calEn}
  \mathcal{E}_{n}\equiv\mathbb{E}\Bigl[e^{\frac{u}{\pi}\im\log p_n(\rho)}e^{a\re \log p_n(\rho)}\Bigr]
  =\frac{D_n}{Z_{n}},
\end{equation}
    where 
\begin{align}
\label{def of Dn}
    D_n&:=\prod_{j=0}^{n-1} \int_0^{+\infty} 2v^{2j+2\alpha+1}e^{-nq(v)}\omega(v)\,dv,
    \\
    \nonumber
    Z_{n}&:=\prod_{j=0}^{n-1} \int_0^{+\infty} 2v^{2j+2\alpha+1}e^{-nq(v)}\,dv.
\end{align}
Here, $\omega(|z|)$ is defined by \eqref{def of omega}.
Note that \eqref{def of Dn} is simply written as 
\begin{align*}
    D_n&=
    \prod_{j=0}^{n-1}
    \bigg( 
e^{u}\int_0^{\rho} 2v^{2j+2\alpha+1}e^{-nq(v)}(\rho-v)^{a}\,dv
+
\int_{\rho}^{+\infty} 2v^{2j+2\alpha+1}e^{-nq(v)}(v-\rho)^{a}\,dv
    \bigg). 
\end{align*}
In order to analyze the precise large $n$-asymptotics of \eqref{def of calEn}, we follow the robust strategy done in \cite{ACC2023c,BC2022,BKS2023}. 
We define 
\begin{equation}
\label{def of V tau}
    V_{\tau}(r):=q(r)-2\tau\log r, \qquad \tau\equiv \tau(j):=\frac{j}{n}. 
\end{equation}
By differentiating \eqref{def of V tau} with respect to $r$, from \cite[Eq. (2.4)]{BKS2023}, we have 
\begin{align}
\begin{split}
\label{def of V tau relationship 1}
V_{\tau}'(r)&=q'(r)-\frac{2\tau}{r},\qquad V_{\tau}''(r)=4\Delta Q(r)-\frac{1}{r}V_{\tau}'(r), \\ 
V_{\tau}^{(3)}(r)&=4\partial_{r}\Delta Q(r)-\frac{4}{r}\Delta Q(r)+\frac{2}{r^2}V_{\tau}'(r), \\
V_{\tau}^{(4)}(r)&=4\partial_r^2\Delta Q(r)+\frac{12}{r^2}\Delta Q(r)-\frac{4}{r}\partial_{r}\Delta Q(r)-\frac{6}{r^3}V_{\tau}'(r). 
\end{split}
\end{align}
We denote $r_{\tau}$ for $0\leq \tau\leq 1$ by
\begin{equation}
\label{def of rtau equation}
    r_{\tau}\,q'(r_{\tau})=2\tau,
\end{equation}
which gives rise to $V_{\tau}'(r_{\tau})=0$. 
Note that $r_{\tau}$ satisfies the following differential equation \cite[Eq. (2.6)]{BKS2023}
\begin{equation}
\label{def of rtau ODE}
\frac{dr_{\tau}}{d\tau}=\frac{1}{2r_{\tau}\Delta Q(r_{\tau})}>0,
\end{equation}
where we have used part (3) of Assumptions~\ref{Assumption_Q}.
By \cite[Eq. (3.7)]{BKS2023}, $r_{\tau}$ satisfies the following asymptotics
\begin{equation}
\label{def of rtau to 0}
    r_{\tau}=\Bigl(\frac{\tau}{\Delta Q(0)}\Bigr)^{\frac{1}{2}}+\mathcal{O}(\tau),\qquad \tau\to0. 
\end{equation}
Let $\tau_{\rho}$ be a solution so that 
\begin{equation}
\label{def of tau rho}
    \rho \,q'(\rho)=2\tau_{\rho}. 
\end{equation}
To split the logarithmic sum of \eqref{def of mathcal En}, we define critical indices 
\begin{equation}
\label{def of critical indices}
    g_{1,-}:=\lceil n(\tau_{\rho}-\delta_n) \rceil,\qquad
    g_{1,+}:=\lfloor n(\tau_{\rho}+\delta_n) \rfloor, \qquad 
    \delta_n':=\frac{M}{\sqrt{n}},\qquad M:=n^{\frac{1}{8}}(\log n)^{-\frac{1}{8}},
\end{equation}
where $\lceil x\rceil$ denotes the smallest integer $\geq x$, and $\lfloor x\rfloor$ denotes the largest integer $\leq x$.
%Our choice \eqref{def of critical indices} for the parameter $M$ plays an essential role to estimate error terms, which are essentially same as \cite[Section 4]{BC2022}.
We also define 
\begin{equation*}
    D_n:=\lfloor n^{\frac{1}{6}}\rfloor,\qquad\delta_n:=\frac{\log n}{\sqrt{n}}.
\end{equation*}
%whose choices are taken from \cite[Section 3]{BKS2023}. 
We write 
\begin{align}\label{def of hnj in}
h_{n,j}^{(\mathrm{in})}(\rho)&:=
\int_0^{\rho} 2ve^{\mathsf{k}(v)}e^{-nV_{\tau}(v)}|\rho-v|^{a}\,dv, 
\\
\label{def of hnj out}
h_{n,j}^{(\mathrm{out})}(\rho)&:=\int_{\rho}^{+\infty} 2ve^{\mathsf{k}(v)}e^{-nV_{\tau}(v)}|v-\rho|^{a}\,dv, 
\\
h_{n,j}&:=\int_0^{+\infty} 2ve^{\mathsf{k}(v)}e^{-nV_{\tau}(v)}\,dv,
\end{align}
where $\mathsf{k}(v):=2\alpha\log v$ for $\alpha>-1$.
Then, we can rewrite $\mathcal{E}_{n}$ as 
\begin{equation}
\label{def of log calEn splitted}
 \log\mathcal{E}_n=\sum_{j=0}^{n-1}\log\Bigl[
e^{u}\frac{h_{n,j}^{(\mathrm{in})}(\rho)}{h_{n,j}}+\frac{h_{n,j}^{(\mathrm{out})}(\rho)}{h_{n,j}}
 \Bigr]=S_0+S_1+S_2+S_3,
\end{equation}
where 
\begin{align*}
S_0&:=\sum_{j=0}^{D_n-1}\log\Bigl[
e^{u}\frac{h_{n,j}^{(\mathrm{in})}(\rho)}{h_{n,j}}+\frac{h_{n,j}^{(\mathrm{out})}(\rho)}{h_{n,j}}
 \Bigr],
\qquad
S_1:=\sum_{j=D_n}^{g_{1,-}-1}\log\Bigl[
e^{u}\frac{h_{n,j}^{(\mathrm{in})}(\rho)}{h_{n,j}}+\frac{h_{n,j}^{(\mathrm{out})}(\rho)}{h_{n,j}}
 \Bigr],
\\
S_2&:=\sum_{j=g_{1,-}}^{g_{1,+}}\log\Bigl[
e^{u}\frac{h_{n,j}^{(\mathrm{in})}(\rho)}{h_{n,j}}+\frac{h_{n,j}^{(\mathrm{out})}(\rho)}{h_{n,j}}
 \Bigr],
\qquad
S_3:=\sum_{j=g_{1,+}+1}^{n-1}\log\Bigl[
e^{u}\frac{h_{n,j}^{(\mathrm{in})}(\rho)}{h_{n,j}}+\frac{h_{n,j}^{(\mathrm{out})}(\rho)}{h_{n,j}}
 \Bigr].
\end{align*}
For the later purpose, we split $S_2$ into 
\begin{equation*}
    S_2=S_2^{(\mathrm{in})}+S_2^{(\mathrm{out})}, 
\end{equation*}
where 
\[
S_2^{(\mathrm{in})}:=\sum_{j=g_{1,-}}^{\lfloor n\tau_{\rho}\rfloor-1}\log\Bigl[
 e^{u}\frac{h_{n,j}^{(\mathrm{in})}(\rho)}{h_{n,j}}+\frac{h_{n,j}^{(\mathrm{out})}(\rho)}{h_{n,j}}
 \Bigr],
\qquad
S_2^{(\mathrm{out})}:=\sum_{j=\lfloor n\tau_{\rho}\rfloor}^{g_{1,+}}\log\Bigl[
e^{u}\frac{h_{n,j}^{(\mathrm{in})}(\rho)}{h_{n,j}}+\frac{h_{n,j}^{(\mathrm{out})}(\rho)}{h_{n,j}}
 \Bigr]. 
\]

\begin{lemma}
\label{lemma:0 j Dn minus1}
For $0\leq j\leq D_{n}-1$, there exists $c>0$ independent of $n$ such that as $n\to+\infty$, we have 
\begin{align}
\label{def of 0 j Dn in}
h_{n,j}^{(\mathrm{in})}(\rho)&=
e^{-nq(0)}\Bigl(\frac{2}{nq''(0)}\Bigr)^{j+\alpha+1}
\rho^{a}\Gamma(j+\alpha+1)
\bigg[1+\mathcal{O}\Bigl(\frac{(j+1)^{3/2}(\log n)^3}{\sqrt{n}}\Bigr)\bigg], 
\\
\label{def of 0 j Dn out}
h_{n,j}^{(\mathrm{out})}(\rho)&=e^{-nq(0)}\cdot \mathcal{O}(e^{-cn}),     
\end{align}
uniformly for $a$ in compact subsets of $(-1,+\infty)$. 
\end{lemma}

\begin{proof}
We recall \cite[Lemma 4.1]{ACC2023c}; 
\begin{equation}
h_{n,j}
=e^{-nq(0)}\Bigl(\frac{2}{nq''(0)}\Bigr)^{j+\alpha+1}
\Gamma(j+\alpha+1)
\bigg[1+\mathcal{O}\Bigl(\frac{(j+1)^{3/2}(\log n)^3}{\sqrt{n}}\Bigr)\bigg]. 
\end{equation}
By a similar manner to the proof of \cite[Lemma 4.1]{ACC2023c}, as $n\to+\infty$, we obtain \eqref{def of 0 j Dn in} and \eqref{def of 0 j Dn out}. 
In particular, the error terms do not depend on $a$. 
This completes the proof. 
\end{proof}

\begin{lemma}\label{lemma:S0}
There exists $\delta>0$ such that as $n\to+\infty$, we have
\begin{equation}
\label{def of S0 asymptotics}
   S_0=D_n\log(\rho^a e^{u})+\mathcal{O}((\log n)^3n^{-\frac{1}{12}}),  
\end{equation}
uniformly for $u\in\{z\in\C:|z-x|\leq \delta\}$ and $a$ in compact subsets of $(-1,+\infty)$. 
\end{lemma}
\begin{proof}
By Lemma~\ref{lemma:0 j Dn minus1} and \cite[Lemma 4.1]{ACC2023c}, 
there exists $\delta>0$ such that we obtain \eqref{def of S0 asymptotics}, uniformly for $u\in\{z\in\C:|z-x|\leq \delta\}$ and $a$ in compact subsets of $(-1,+\infty)$. 
\end{proof}

We next turn to $S_1$ and $S_3$. 
We use the following lemma from \cite[Lemma 4.3]{ACC2023c}.
\begin{lemma}\label{lemma:hnj asymptotics}
For $D_n\leq j\leq n-1$ and $\mathsf{k}(r)=2\alpha\log r$, as $n\to+\infty$, we have 
\begin{equation}
\label{def of Dn leq j leq n-1}
h_{n,j}
=\sqrt{\frac{2\pi}{n}} \frac{r_{\tau}}{\sqrt{\Delta Q(r_{\tau})}}e^{\mathsf{k}(r_{\tau})}e^{-nV_{\tau}(r_{\tau})}\cdot\Bigl(1+\frac{\mathcal{A}(r_{\tau})}{n}+\mathcal{O}(\frac{(\log n)^{\nu}}{j^{3/2}})\Bigr)    
\end{equation}
for some $\nu>0$, where
\begin{align*}
\mathcal{A}(r_{\tau})&:=\mathcal{B}(r_{\tau})+\frac{\mathsf{k}'(r_{\tau})^2}{2}\frac{1}{d_2}+\frac{\mathsf{k}''(r_{\tau})}{2}\frac{1}{d_2}+\frac{\mathsf{k}'(r_{\tau})}{r_{\tau}}\frac{1}{d_2}-\frac{\mathsf{k}'(r_{\tau})}{2}\frac{d_3}{d_2^2},   
\\
\mathcal{B}(r)&:=-\frac{1}{32}\frac{\partial_r^2\Delta Q(r)}{(\Delta Q(r))^2}-\frac{19}{96}\frac{\partial_r\Delta Q(r)}{(\Delta Q(r))^2}+\frac{5}{96}\frac{(\partial_r \Delta Q(r))^2}{(\Delta Q(r))^3}+\frac{1}{12}\frac{1}{r^2\Delta Q(r)},      \\
d_m&:=V_{\tau}^{(m)}(r_{\tau}).
\end{align*}
Here, $\mathcal{O}(j^{-3/2}(\log n)^{\nu})$ can be replaced with $\mathcal{O}(n^{-2})$ for large $j$, i.e., for $j\geq c_0 n$ with some $c_0>0$. 
\end{lemma}
We first establish the large $n$-asymptotics of $S_1$.
To this end, for a sufficiently small $\epsilon>0$, let 
\[
j_{1,-}:=\lceil n(\tau_{\rho}-\epsilon)\rceil.    
\]
We write 
\begin{align*}
\theta_{1,-}^{(\epsilon)}&:=j_{1,-}-n(\tau_{\rho}-\epsilon),
\quad
\theta_{D_n}:=n^{1/6}-D_n, 
\\
\theta_{1,-}&:=\lceil n(\tau_{\rho}-\delta_n)\rceil-n(\tau_{\rho}-\delta_n),\quad
\theta_{1,+}:=n(\tau_{\rho}+\delta_n)-\lfloor n(\tau_{\rho}+\delta_n)\rfloor. 
\end{align*}
Let us denote
\begin{align}
\label{def of Delta nplus}
    \Delta_{n,+}&:=\sqrt{\frac{n}{\Delta Q(\rho)}}\frac{\delta_n'-\tfrac{\theta_{1,+}}{n}}{\rho}=\frac{1}{\rho\sqrt{\Delta Q(\rho)}}\Bigl(M-\frac{\theta_{1,+}}{\sqrt{n}}\Bigr),
    \\
    \label{def of Delta nminus}
    \Delta_{n,-}&:=\sqrt{\frac{n}{\Delta Q(\rho)}}\frac{\delta_n'-\tfrac{\theta_{1,-}}{n}}{\rho}=\frac{1}{\rho\sqrt{\Delta Q(\rho)}}\Bigl(M-\frac{\theta_{1,-}}{\sqrt{n}}\Bigr).
\end{align}
By the implicit function theorem, Assumptions~\ref{Assumption_Q}, and \eqref{def of rtau ODE}, as $n\to+\infty$, $r_{\tau(g_{\pm})}$ and $r_{\tau(j_{1,-})}$ are expanded as  
\begin{align}
\label{def of rtau g1plus}
    r_{\tau(g_{1,+})}
&=
\rho+\frac{1}{2\sqrt{\Delta Q(\rho)}}\frac{\Delta_{n,+}}{\sqrt{n}}
-\frac{\Delta Q(\rho)+\rho\partial_r\Delta Q(\rho)}{8\rho \Delta Q(\rho)^2}\frac{\Delta_{n,+}^2}{n}+\mathcal{O}\Bigl(\frac{\Delta_{n,+}^3}{n^{3/2}}\Bigr), 
\\
\label{def of rtau g1minus}
r_{\tau(g_{1,-})}
&=\rho-\frac{1}{2\sqrt{\Delta Q(\rho)}}\frac{\Delta_{n,-}}{\sqrt{n}}
-\frac{\Delta Q(\rho)+\rho\partial_r\Delta Q(\rho)}{8\rho \Delta Q(\rho)^2}\frac{\Delta_{n,-}^2}{n}+\mathcal{O}\Bigl(\frac{\Delta_{n,-}^3}{n^{3/2}}\Bigr).
%\\
%\label{def of rtau j1minus}
%r_{\tau(j_{1,-})}
%&=
%\rho-\frac{1}{2\sqrt{\Delta Q(\rho)}}\frac{\Delta_{n,-}^{(\epsilon)}}{\sqrt{n}}
%-\frac{\Delta Q(\rho)+\rho\partial_r\Delta Q(\rho)}{8\rho \Delta Q(\rho)^2}\frac{{\Delta_{n,-}^{(\epsilon)}}^2}{n}+\mathcal{O}\Bigl(\frac{{\Delta_{n,-}^{(\epsilon)}}^3}{n^{3/2}}\Bigr),
\end{align}
\begin{lemma}\label{lemma:S1}
There exists $\delta>0$ such that as $n\to+\infty$, we have 
\begin{equation}
\label{def of S1 asymptotics}
S_1= C_1^{(1)}n + C_2^{(1)}\sqrt{n}+C_3^{(1)}+C_n^{(1)}+
\mathcal{O}\Bigl(\frac{(\log n)^{3}}{n^{\frac{1}{12}}}\Bigr),    
\end{equation}
uniformly for $u\in\{z\in\C:|z-x|\leq \delta\}$ and $a$ in compact subsets of $(-1,+\infty)$, where 
\begin{align*}
C_1^{(1)}&:=
\tau_{\rho}u
+
\int_{0}^{\rho}2a u\Delta Q(u)\log(\rho-u)\,du
\qquad
C_2^{(1)}:=0,
\\
C_3^{(1)}&:=
\frac{a}{2}\log 2
+\frac{a}{4}\log\Delta Q(\rho)
+\frac{a}{2}\log\rho
-
\frac{a(a-1)}{8}\Bigl(3+\frac{\rho\,\partial_r\Delta Q(\rho)}{\Delta Q(\rho)}\Bigr)
\\
&\quad
+
\frac{a}{4}\int_0^{\rho}\Bigl(\frac{1}{\rho-u} 
\frac{u\partial_r\Delta Q(u)}{\Delta Q(u)}
-
\frac{1}{\rho-u} 
\frac{\rho\partial_r\Delta Q(\rho)}{\Delta Q(\rho)}
\Bigr)\,du
\\
&\quad
-\frac{a}{4}
\log(2\rho\sqrt{\Delta Q(\rho)})
\Bigl(4\alpha+a+2-\frac{\rho\,\partial_r\Delta Q(\rho)}{\Delta Q(\rho)}\Bigr),
\\
C_n^{(1)}&:=
-\rho\sqrt{\Delta Q(\rho)}\sqrt{n}\Delta_{n,-}u
-D_nu
-aD_{n}\log\rho
\\
&\quad
-\frac{a}{2}\rho\sqrt{\Delta Q(\rho)}
\Bigl(2\log\Delta_{n,-}-2\log 2 -\log\Delta Q(\rho)-\log n-2\Bigr)\Delta_{n,-}\sqrt{n}
\\
&\quad
-\frac{a}{8}\Bigl(1+\frac{\rho\partial_r\Delta Q(\rho)}{\Delta Q(\rho)}\Bigr)\Delta_{n,-}^2
-\frac{a}{2}\log\Delta_{n,-}
+\frac{a}{4}\log n
+\frac{a(a-1)}{2}
\frac{\sqrt{n}}{\Delta_{n,-}}\rho\sqrt{\Delta Q(\rho)}
\\
&\quad
-\frac{a(a-1)(2a-3)}{12}\rho\sqrt{\Delta Q(\rho)}\frac{\sqrt{n}}{\Delta_{n,-}^3}
+\frac{a}{4}\Bigl(\log\Delta_{n,-}-\frac{1}{2}\log n\Bigr)
\Bigl( 4\alpha+a+2 -\frac{\rho\partial_r\Delta Q(\rho)}{\Delta Q(\rho)} \Bigr).
\end{align*}
\end{lemma}
\begin{proof}
By Assumptions~\ref{Assumption_Q} on the potential $Q$ and for $j\in\{D_n,\dots,g_{1,-}-1\}$, and a sufficiently large $n\in\N$, there exists a unique critical point $r_{\tau}\in(0,\rho)$, which satisfies $r_{\tau}q'(r_{\tau})=2\tau$. 
This implies that there is no critical point inside $(\rho,+\infty)$. 
By the integrability of the exponential term $e^{-nq(r)}$, 
we find that there exists $c>0$ such that 
\begin{equation}
    h_{n,j}^{(\mathrm{out})}(\rho)=e^{-nV_{\tau}(r_{\tau})}\cdot \mathcal{O}(e^{-cn}). 
\end{equation}
Therefore, we focus on \eqref{def of hnj in}. We split the integral
\begin{align*}
h_{n,j}^{(\mathrm{in})}(\rho)
&=
\int_{[0,\rho]\cap \{r:|r-r_{\tau}|<\delta_n\}}2r^{2\alpha+1}e^{-nV_{\tau}(r)}(\rho-r)^{a}\,dr
+
\int_{[0,\rho]\cap \{r:|r-r_{\tau}|\geq \delta_n\}}2r^{2\alpha+1}e^{-nV_{\tau}(r)}(\rho-r)^{a}\,dr
\\
&=
\int_{[0,\rho]\cap \{r:|r-r_{\tau}|<\delta_n\}}2r e^{\mathsf{k}(r)}e^{-nV_{\tau}(r)}(\rho-r)^{a}\,dr
+
e^{-nV_{\tau}(r_{\tau})}\cdot \mathcal{O}(e^{-c(\log n)^2}), 
\end{align*}
where $\mathsf{k}(r):=2\alpha\log r$.
Here, the exponential small error follows from a similar manner to \cite[Proof of Lemma 2.1]{BKS2023}.
For the first term in the last line, we apply the Laplace method, and then we have 
\begin{align*}
&\quad
\int_{[0,\rho]\cap \{r:|r-r_{\tau}|<\delta_n\}}2r e^{\mathsf{k}(r)}e^{-nV_{\tau}(r)}|r-\rho|^{a}\,dr
\\
&=
\frac{2r_{\tau}e^{\mathsf{k}(r_{\tau})}e^{-nV_{\tau}(r_{\tau})}}{\sqrt{nd_2}}\int_{-\sqrt{d_2n}\delta_n}^{\sqrt{d_2n}\delta_n}
e^{-\frac{1}{2}u^2}
\Bigl|r_{\tau}-\rho+\frac{u}{\sqrt{nd_2}}\Bigr|^{a}
\Bigl( 
1+\frac{c_1(u)}{\sqrt{n}}+\frac{c_2(u)}{n}+\frac{c_3(u)}{n^{3/2}}+\mathcal{O}(\frac{c_4(u)}{n^2})
\Bigr)\,du
\\
&=
\frac{\sqrt{2\pi}r_{\tau}e^{\mathsf{k}(r_{\tau})}e^{-nV_{\tau}(r_{\tau})}}{\sqrt{n\Delta Q(r_{\tau})}}
(\rho-r_{\tau})^a
\Bigl[
1+\frac{1}{n}\Bigl(\mathcal{M}_1^{(\rm in)}(r_{\tau})+\mathcal{A}_1(r_{\tau})\Bigr)
\\
&\quad
+\frac{1}{n^2}
\Bigl\{
\frac{a(a-1)(a-2)(a-3)u^4}{24d_2^2(\rho-r_{\tau})^4}
+\mathcal{A}_2(r_{\tau})
+\mathcal{O}\Bigl(\frac{1}{(\rho-r_{\tau})^3}\Bigr)
\Bigr\}
\Bigr]
+
\mathcal{O}(e^{-c(\log n)^2}), 
\end{align*}
where the error term depends on $a$, but it does not affect the order of the error. Here, 
\begin{align*}
\mathcal{M}_1^{(\rm in)}(r_{\tau})
&=
\frac{a(a-1)(\rho-r_{\tau})^{-2}}{8\Delta Q(r_{\tau})}
-\frac{a(\rho-r_{\tau})^{-1}}{4\Delta Q(r_{\tau})}
\Bigl( 
-\frac{\partial_r\Delta Q(r_{\tau})}{2\Delta Q(r_{\tau})}
+\frac{4\alpha+3}{2r_{\tau}}
\Bigr).
\end{align*}
and, we safely extended the integral region to $(-\infty,\infty)$ with the exponential error $\mathcal{O}(e^{-c(\log n)^2})$ for some $c>0$, and for $d_{j}:=V_{\tau}^{(j)}(r_{\tau})$ given by \eqref{def of V tau relationship 1}, $c_k(u)$ for $k=1,2,3$ are given by 
\begin{align}
\label{def of c_1}
c_1(u)&:=-\frac{d_3}{6d_2^{3/2}}u^3+\Bigl(\mathsf{k}'(r_{\tau})+\frac{1}{r_{\tau}}\Bigr)\frac{1}{d_2^{1/2}}u, 
\\
\label{def of c_2}
c_2(u)&:=\frac{d_3^2}{72d_2^3}u^6-\Bigl(d_4+4d_3\mathsf{k}'(r_{\tau})+\frac{4d_3}{r_{\tau}}\Bigr)\frac{u^4}{24d_2^2}+\Bigl(\frac{\mathsf{k}''(r_{\tau})}{2}+\frac{\mathsf{k}'(r_{\tau})^2}{2}+\frac{\mathsf{k}'(r_{\tau})}{r_{\tau}}\Bigr)\frac{u^2}{d_2}
\\
\begin{split}
\label{def of c_3}
c_3(u)&:=-\frac{d_3^3}{1296d_2^{9/2}}u^9+\Bigl(d_3d_4+2d_3^2\mathsf{k}'(r_{\tau})+\frac{2d_3^2}{r_{\tau}} \Bigr)\frac{u^7}{144d_2^{7/2}}
\\
&\quad
-\Bigl(
d_5+\frac{5}{r_{\tau}}d_4+5d_4\mathsf{k}'(r_{\tau})+\frac{20d_3}{r_{\tau}}\mathsf{k}'(r_{\tau})+10d_3\mathsf{k}'(r_{\tau})^2+10d_3\mathsf{k}''(r_{\tau})
\Bigr)
\frac{u^5}{120d_2^{5/2}}
\\
&\quad
+
\Bigl(\frac{3\mathsf{k}'(r_{\tau})^2}{r_{\tau}}
+\frac{3\mathsf{k}''(r_{\tau})}{r_{\tau}}+\mathsf{k}'(r_{\tau})^3+3\mathsf{k}'(r_{\tau})\mathsf{k}''(r_{\tau})+\mathsf{k}^{(3)}(r_{\tau})\Bigr)\frac{u^3}{6d_2^{3/2}},
\end{split}
\end{align}
and $c_4(u)$ is a polynomial consisting of $u^{12},u^{10},u^{8},u^6,u^4$.
Since $\rho-r_{\tau}>0$ and $\rho-r_{\tau}\searrow c n^{-1/2}M$ with some $c>0$ for $D_{n}\leq j\leq g_{1,-}-1$, by Lemma~\ref{lemma:hnj asymptotics}, we have 
\begin{align*}
S_1
&=
\sum_{j=D_n}^{g_{1,-}-1}
\log\bigl(e^{u}(\rho-r_{\tau})^{a}\bigr)
+\frac{1}{n}\sum_{j=j_{1,-}}^{g_{1,-}-1}
\Bigl[
\frac{a(a-1)(\rho-r_{\tau})^{-2}}{8\Delta Q(r_{\tau})}
-\frac{a(\rho-r_{\tau})^{-1}}{4\Delta Q(r_{\tau})}
\Bigl( 
-\frac{\partial_r\Delta Q(r_{\tau})}{2\Delta Q(r_{\tau})}+\frac{4\alpha+3}{2r_{\tau}}
\Bigr)
\Bigr]
\\
&\quad
-\frac{1}{n^2}
\sum_{j=D_n}^{g_{1,-}-1}
\frac{a(a-1)(2a-3)}{64\Delta Q(r_{\tau})^2(\rho-r_{\tau})^4}
+\sum_{j=D_n}^{g_{1,-}-1}
\mathcal{O}\Bigl(\frac{1}{n^{2}}\frac{1}{(\rho-r_{\tau})^3}\Bigr)
+
\mathcal{O}\Bigl(\frac{(\log n)^{3}}{n^{\frac{1}{12}}}\Bigr),
\end{align*}
where the second error term is independent of $a,u$.
By Lemma~\ref{lemma:Euler-Maclaurin} and by change of variable $r_{\tau(j)}=t$, we have 
\[
\sum_{j=D_n}^{g_{1,-}-1}
\mathcal{O}\Bigl(\frac{1}{n^{2}}\frac{1}{(\rho-r_{\tau})^3}\Bigr)
=
\mathcal{O}\Bigl(\frac{(\log n)^{\frac{3}{8}}}{n^{\frac{7}{8}}}\Bigr). 
\]
Consequently, the expansion of $S_1$ takes the following form:
\begin{align*}
S_1&=(g_{1,-}-D_n)u+a\int_{D_n}^{g_{1,-}}\log(\rho-r_{\tau(t)})\,dt
-\frac{a}{2}\log(\rho-r_{\tau(g_{1,-})})
+\frac{a}{2}\log(\rho-r_{\tau(D_n)})
\\
&\quad
+\frac{1}{n}\sum_{j=D_n}^{g_{1,-}-1}
\Bigl[
\frac{a(a-1)(\rho-r_{\tau})^{-2}}{8\Delta Q(r_{\tau})}
-\frac{a(\rho-r_{\tau})^{-1}}{4\Delta Q(r_{\tau})}
\Bigl( 
-\frac{\partial_r\Delta Q(r_{\tau})}{2\Delta Q(r_{\tau})}+\frac{4\alpha+3}{2r_{\tau}}
\Bigr)
\Bigr]
\\
&\quad
-\frac{1}{n^2}
\sum_{j=D_n}^{g_{1,-}-1}
\frac{a(a-1)(2a-3)}{64\Delta Q(r_{\tau})^2(\rho-r_{\tau})^4}
+\mathcal{O}\Bigl(\frac{(\log n)^{3}}{n^{\frac{1}{12}}}\Bigr), 
\end{align*}
where we have used Lemma~\ref{lemma:Euler-Maclaurin}, and the error term is independent of $a,u$.
By \eqref{def of rtau to 0} and \eqref{def of rtau g1minus}, we have 
\begin{align*}
a\int_{D_n}^{g_{1,-}}\log(\rho-r_{\tau(t)})\,dt
&=
n\int_{0}^{\rho}2a u\Delta Q(u)\log(\rho-u)\,du
-aD_{n}\log\rho
\\
&\quad
-\frac{a}{2}\rho\sqrt{\Delta Q(\rho)}
\Bigl(2\log\Delta_{n,-}-2\log 2 -\log\Delta Q(\rho)-\log n-2\Bigr)\Delta_{n,-}\sqrt{n}
\\
&\quad
-\frac{a}{8}\Bigl(1+\frac{\rho\partial_r\Delta Q(\rho)}{\Delta Q(\rho)}\Bigr)\Delta_{n,-}^2
+\mathcal{O}\Bigl(\frac{(\log n)^{1/4}}{n^{1/4}}\Bigr),
\\
-\frac{a}{2}\log(\rho-r_{\tau(g_{1,-})})+\frac{a}{2}\log(\rho-r_{\tau(D_n)})
&=
-\frac{a}{2}\log\Delta_{n,-}
+\frac{a}{2}\log 2
+\frac{a}{4}\log\Delta Q(\rho)
+\frac{a}{4}\log n
+\frac{a}{2}\log\rho
+\mathcal{O}\Bigl(\frac{(\log n)^{1/4}}{n^{1/4}}\Bigr). 
\end{align*}
The Taylor theorem gives rise to
\begin{align*}
&\quad
a\int_{D_n}^{g_{1,-}}\log(\rho-r_{\tau(t)})\,dt
-\frac{a}{2}\log(\rho-r_{\tau(g_{1,-})})
+\frac{a}{2}\log(\rho-r_{\tau(D_n)})
\\
&=
n\int_{0}^{\rho}2a u\Delta Q(u)\log(\rho-u)\,du
+\frac{a}{2}\log 2
+\frac{a}{4}\log\Delta Q(\rho)
+\frac{a}{2}\log\rho
\\
&\quad
-\frac{a}{2}\rho\sqrt{\Delta Q(\rho)}
\Bigl(2\log\Delta_{n,-}-2\log 2 -\log\Delta Q(\rho)-\log n-2\Bigr)\Delta_{n,-}\sqrt{n}
\\
&\quad
-aD_{n}\log\rho
-\frac{a}{8}\Bigl(1+\frac{\rho\partial_r\Delta Q(\rho)}{\Delta Q(\rho)}\Bigr)\Delta_{n,-}^2
-\frac{a}{2}\log\Delta_{n,-}
+\frac{a}{4}\log n
+\mathcal{O}\Bigl(\frac{(\log n)^{1/4}}{n^{1/4}}\Bigr). 
\end{align*}
By Lemma~\ref{lemma:Euler-Maclaurin}, we have 
\begin{align*}
-\frac{1}{n^2}
\sum_{j=D_n}^{g_{1,-}-1}
\frac{a(a-1)(2a-3)}{64\Delta Q(r_{\tau})^2(\rho-r_{\tau})^4}
&=
-\frac{1}{n}\int_{r_{\tau(D_n)}}^{r_{\tau(g_{1,-})}}
\frac{a(a-1)(2a-3)}{32\Delta Q(u)(\rho-u)^4}u\,du+\mathcal{O}(\Delta_{n,-}^{-4}).
\end{align*}
Integration by parts leads to
\[
-\frac{1}{n}\int_{r_{\tau(D_n)}}^{r_{\tau(g_{1,-})}}
\frac{a(a-1)(2a-3)}{32\Delta Q(u)(\rho-u)^4}u\,du
=
-
\frac{1}{n}\frac{a(a-1)(2a-3)r_{\tau(g_{1,-})}}{96\Delta Q(r_{\tau(g_{1,-})})(\rho-r_{\tau(g_{1,-})})^3}
+
\widetilde{\epsilon}_{n,-}
+
\mathcal{O}(n^{-1}), 
\]
where 
\[
\widetilde{\epsilon}_{n,-}:=\frac{1}{n}\int_{r_{\tau(D_n)}}^{r_{\tau(g_{1,-})}}
\frac{a(a-1)(2a-3)}{96(\rho-u)^3}\frac{\Delta Q(u)-u\partial_u\Delta Q(u)}{\Delta Q(u)^2}\,du. 
\]
Note that 
\begin{align*}
-
\frac{1}{n}\frac{a(a-1)(2a-3)r_{\tau(g_{1,-})}}{96\Delta Q(r_{\tau(g_{1,-})})(\rho-r_{\tau(g_{1,-})})^3}
&=-\frac{a(a-1)(2a-3)}{12}\rho\sqrt{\Delta Q(\rho)}\frac{\sqrt{n}}{\Delta_{n,-}^3}+\mathcal{O}(\Delta_{n,-}^{-2}), 
\\
|\widetilde{\epsilon}_{n,-}|&\lesssim \frac{1}{n}\bigg|\int_{r_{\tau(D_n)}}^{r_{\tau(g_{1,-})}}\frac{du}{(\rho-u)^3}\bigg|=\mathcal{O}(\Delta_{n,-}^{-2}).
\end{align*}
By Lemma~\ref{lemma:Euler-Maclaurin}, we have 
\begin{align*}
&\quad\frac{1}{n}\sum_{j=D_n}^{g_{1,-}-1}
\Bigl[
\frac{a(a-1)(\rho-r_{\tau})^{-2}}{8\Delta Q(r_{\tau})}
-\frac{a(\rho-r_{\tau})^{-1}}{4\Delta Q(r_{\tau})}
\Bigl( 
-\frac{\partial_r\Delta Q(r_{\tau})}{2\Delta Q(r_{\tau})}+\frac{4\alpha+3}{2r_{\tau}}
\Bigr)
\Bigr]
\\
&=
\frac{a(a-1)}{4}\bigg[\frac{\rho}{\rho-r_{\tau(g_{1,-})}}+\log(\rho-r_{\tau(g_{1,-})})-\frac{\rho}{\rho-r_{\tau(D_n)}}-\log(\rho-r_{\tau(D_n)})\bigg]
\\
&\quad
+\frac{a}{4}\Bigl(\log(\rho-r_{\tau(g_{1,-})})-\log(\rho-r_{\tau(D_n)})\Bigr)
\Bigl(4\alpha+3-\frac{\rho\,\partial_r\Delta Q(\rho)}{\Delta Q(\rho)}\Bigr)
\\
&\quad
+\frac{a}{4}\int_{r_{\tau(D_n)}}^{r_{\tau(g_{1,-})}}\Bigl(\frac{1}{\rho-u}\frac{u\,\partial_r\Delta Q(u)}{\Delta Q(u)}-\frac{1}{\rho-u}\frac{\rho\,\partial_r\Delta Q(\rho)}{\Delta Q(\rho)}\Bigr)\,du
+\mathcal{O}(\Delta_{n,-}^{-2}).
\end{align*}
Using the asymptotics by the Taylor theorem, 
\begin{align*}
\frac{a}{4}\int_{r_{\tau(D_n)}}^{r_{\tau(g_{1,-})}}\Bigl(\frac{1}{\rho-u}\frac{u\,\partial_r\Delta Q(u)}{\Delta Q(u)}-\frac{1}{\rho-u}\frac{\rho\,\partial_r\Delta Q(\rho)}{\Delta Q(\rho)}\Bigr)\,du
&=
\frac{a}{4}\int_{0}^{\rho}\Bigl(\frac{1}{\rho-u}\frac{u\,\partial_r\Delta Q(u)}{\Delta Q(u)}-\frac{1}{\rho-u}\frac{\rho\,\partial_r\Delta Q(\rho)}{\Delta Q(\rho)}\Bigr)\,du
\\
&\quad
+\mathcal{O}\Bigl(\frac{\Delta_{n,-}}{\sqrt{n}}+\frac{D_n}{n}\Bigr), 
\end{align*}
\begin{align*}
&\quad\frac{\rho}{\rho-r_{\tau(g_{1,-})}}+\log(\rho-r_{\tau(g_{1,-})})-\frac{\rho}{\rho-r_{\tau(D_n)}}-\log(\rho-r_{\tau(D_n)})
\\
&=
2\rho\sqrt{\Delta Q(\rho)}\frac{\sqrt{n}}{\Delta_{n,-}}-\frac{1}{2}\Bigl(3+\frac{\rho\,\partial_r\Delta Q(\rho)}{\Delta Q(\rho)}\Bigr)-\log(2\rho\sqrt{\Delta Q(\rho)})-\frac{1}{2}\log n+\log\Delta_{n,-}+\mathcal{O}\Bigl(\sqrt{\frac{D_n}{n}}\Bigr), 
\end{align*}
and 
\[
\log(\rho-r_{\tau(g_{1,-})})-\log(\rho-r_{\tau(D_n)})
=-\frac{1}{2}\log n+\log\Delta_{n,-}-\log(2\rho\sqrt{\Delta Q(\rho)})
+\mathcal{O}\Bigl(\sqrt{\frac{D_n}{n}}\Bigr),
\]
we have 
\begin{align*}
&\quad\frac{1}{n}\sum_{j=D_n}^{g_{1,-}-1}
\Bigl[
\frac{a(a-1)(\rho-r_{\tau})^{-2}}{8\Delta Q(r_{\tau})}
-\frac{a(\rho-r_{\tau})^{-1}}{4\Delta Q(r_{\tau})}
\Bigl( 
-\frac{\partial_r\Delta Q(r_{\tau})}{2\Delta Q(r_{\tau})}+\frac{4\alpha+3}{2r_{\tau}}
\Bigr)
\Bigr]
\\
&=
-
\frac{a(a-1)}{4}\bigg[
\frac{1}{2}\Bigl(3+\frac{\rho\,\partial_r\Delta Q(\rho)}{\Delta Q(\rho)}\Bigr)
+\log(2\rho\sqrt{\Delta Q(\rho)})
\bigg]
-\frac{a}{4}
\log(2\rho\sqrt{\Delta Q(\rho)})
\Bigl(4\alpha+3-\frac{\rho\,\partial_r\Delta Q(\rho)}{\Delta Q(\rho)}\Bigr)
\\
&\quad
+\frac{a}{4}\int_{0}^{\rho}\Bigl(\frac{1}{\rho-u}\frac{u\,\partial_r\Delta Q(u)}{\Delta Q(u)}-\frac{1}{\rho-u}\frac{\rho\,\partial_r\Delta Q(\rho)}{\Delta Q(\rho)}\Bigr)\,du
\\
&\quad
+\frac{a}{4}\Bigl(
-\frac{1}{2}\log n+\log\Delta_{n,-}
\Bigr)
\Bigl(4\alpha+a+2-\frac{\rho\,\partial_r\Delta Q(\rho)}{\Delta Q(\rho)}\Bigr)
+\frac{a(a-1)}{2}
\rho\sqrt{\Delta Q(\rho)}\frac{\sqrt{n}}{\Delta_{n,-}}
+\mathcal{O}(\Delta_{n,-}^{-2}). 
\end{align*}
Combining all the above, we obtain \eqref{def of S1 asymptotics}.
\end{proof}
Next, we compute the large $n$-asymptotics of $S_3$. 
\begin{lemma}\label{lemma:S3}
There exists $\delta>0$ such that as $n\to+\infty$, we have 
\begin{equation}
\label{def of S3 asymptotics}
 S_3=
 C_1^{(3)}n+C_2^{(3)}\sqrt{n}+C_3^{(3)}+C_n^{(3)}+\mathcal{O}\Bigl(\frac{(\log n)^{\frac{1}{4}}}{n^{\frac{1}{4}}}\Bigr),
\end{equation}
uniformly for $u\in\{z\in\C:|z-x|\leq \delta\}$ and $a$ in compact subsets of $(-1,+\infty)$, where 
\begin{align*}
C_1^{(3)}&:=
a\int_{\rho}^{r_1}\log (t-\rho)\,2t\Delta Q(t)\,dt,
\\
C_2^{(3)}&:=0, 
\\
C_3^{(3)}&:=
-\frac{a}{2}\log(r_1-\rho)
+\frac{a}{2}\log2
+\frac{a}{4}\log(\Delta Q(\rho))
+\frac{a(a-1)}{4}\bigg[\frac{1}{2}\Bigl(1+\frac{\rho\,\partial_r\Delta Q(\rho)}{\Delta Q(\rho)}\Bigr)-\frac{\rho}{r_1-\rho}\bigg]
\\
&\quad
+\frac{a}{4}\Bigl(4\alpha+a+2-\frac{\rho\,\partial_r\Delta Q(\rho)}{\Delta Q(\rho)}\Bigr)\log(2(r_1-\rho)\sqrt{\Delta Q(\rho)})
\\
&\quad
-\frac{a}{4}\int_{\rho}^{r_1}\Bigl(\frac{1}{t-\rho}\frac{t\,\partial_t\Delta Q(t)}{\Delta Q(t)}-\frac{1}{t-\rho}\frac{\rho\,\partial_r\Delta Q(\rho)}{\Delta Q(\rho)}\Bigr)\,dt, 
\\
\begin{split}
C_n^{(3)}&:=
-\frac{a}{2}\rho\sqrt{\Delta Q(\rho)}
\Bigl( 
2\log\Delta_{n,+}-2\log 2-\log\Delta Q(\rho)-\log n-2
\Bigr)\sqrt{n}\Delta_{n,+}
\\
&\quad
+\frac{a}{8}\Bigl(1+\frac{\rho\partial_r\Delta Q(\rho)}{\Delta Q(\rho)}\Bigr)\Delta_{n,+}^2
-\frac{a}{2}\log(\Delta_{n,+})
+\frac{a}{4}\log n
+\frac{a(a-1)}{2}\rho\sqrt{\Delta Q(\rho)}\frac{\sqrt{n}}{\Delta_{n,+}}
\\
&\quad
-\frac{a}{4}
\Bigl( 
\log\Delta_{n,+}-\frac{1}{2}\log n
\Bigr)
\Bigl(
4\alpha+a+2-\frac{\rho\partial_r\Delta Q(\rho)}{\Delta Q(\rho)}
\Bigr)
-\frac{a(a-1)(2a-3)}{12}\rho\sqrt{\Delta Q(\rho)}\frac{\sqrt{n}}{\Delta_{n,+}^3}. 
\end{split}
\end{align*}


\end{lemma}

\begin{proof}
Similar manner to Lemma~\ref{lemma:S1}, there exists $c>0$ such that as $n\to+\infty$, we have  
\begin{align}
\begin{split}
\label{def of hnj out in S3}
h_{n,j}^{(\mathrm{out})}(\rho)
&=  
\frac{\sqrt{2\pi}r_{\tau}e^{\mathsf{k}(r_{\tau})}e^{-nV_{\tau}(r_{\tau})}}{\sqrt{n\Delta Q(r_{\tau})}}
(r_{\tau}-\rho)^a
\bigg[
1+\frac{1}{n}
\Bigl(
\mathcal{A}_1(r_{\tau})
+\mathcal{M}_1^{(\mathrm{out})}(r_{\tau})
\Bigr)
\\
&\quad
+\frac{1}{n^2}
\Bigl\{
\frac{a(a-1)(a-2)(a-3)u^4}{24d_2^2(r_{\tau}-\rho)^4}
+\mathcal{A}_2(r_{\tau})
+\mathcal{O}\Bigl(\frac{1}{(r_{\tau}-\rho)^3}\Bigr)
\Bigr\}
\bigg]
+
\mathcal{O}(e^{-c(\log n)^2}),
\end{split}
\end{align}
where the constant in the error term might depend on $a$, but the order of the error term is independent of $a$. Here, 
\[
\mathcal{M}_1^{(\mathrm{out})}(r_{\tau})
=
\frac{a(a-1)(r_{\tau}-\rho)^{-2}}{8\Delta Q(r_{\tau})}
+\frac{a(r_{\tau}-\rho)^{-1}}{8\Delta Q(r_{\tau})}
\Bigl(-\frac{\partial_r\Delta Q(r_{\tau})}{\Delta Q(r_{\tau})}
+\frac{4\alpha+3}{r_{\tau}}
\Bigr). 
\]
Also, there exists $c>0$ independent of $n$ such that we have 
\begin{equation}
\label{def of hnj in in S3}    
e^{u}\frac{h_{n,j}^{(\mathrm{in})}(\rho)}{h_{n,j}}=e^{-nV_{\tau}(r_{\tau})}\cdot\mathcal{O}(e^{-cn}), 
\end{equation}
where we can take the error term to be independent of $u$.
Thus, by Lemma~\ref{lemma:hnj asymptotics}, \eqref{def of hnj out in S3}, and \eqref{def of hnj in in S3}, let us choose $\delta>0$ sufficiently small so that 
\[
e^{u}\frac{h_{n,j}^{(\mathrm{in})}(\rho)}{h_{n,j}}+\frac{h_{n,j}^{(\mathrm{out})}(\rho)}{h_{n,j}}
\]
remains bounded away from the interval $(-\infty,0]$ as $n\to+\infty$ uniformly for $u\in\{z\in\C:|z-x|\leq \delta\}$ and $a$ in compact subsets of $(-1,+\infty)$. 
Therefore, as $n\to+\infty$, we have 
\begin{align*}
&\quad\sum_{j=g_{1,+}+1}^{n-1}\log\Bigl[
e^{u}\frac{h_{n,j}^{(\mathrm{in})}(\rho)}{h_{n,j}}+\frac{h_{n,j}^{(\mathrm{out})}(\rho)}{h_{n,j}}
 \Bigr]   
 \\
&= 
a\sum_{j=g_{1,+}+1}^{n-1}\log(r_{\tau}-\rho)
+
\frac{1}{n}
\sum_{j=g_{1,+}+1}^{n-1}
\Bigl\{ 
\frac{a(a-1)(r_{\tau}-\rho)^{-2}}{8\Delta Q(r_{\tau})}
+\frac{a(r_{\tau}-\rho)^{-1}}{8\Delta Q(r_{\tau})}
\Bigl(-\frac{\partial_r\Delta Q(r_{\tau})}{\Delta Q(r_{\tau})}
+\frac{4\alpha+3}{r_{\tau}}
\Bigr)
\Bigr\}
\\
&\quad
-\frac{1}{n^2}
\sum_{j=g_{1,+}+1}^{n-1}
\frac{a(a-1)(2a-3)}{64\Delta Q(r_{\tau})^2(r_{\tau}-\rho)^4}
+
\sum_{j=g_{1,+}+1}^{n-1}
\mathcal{O}\Bigl(\frac{(r_{\tau}-\rho)^{-3}}{n^2}\Bigr), 
\end{align*} 
uniformly for $u\in\{z\in\C:|z-x|\leq \delta\}$ and $a$ in compact subsets of $(-1,+\infty)$. 
By change of variable $r_{\tau(j)}=t$ and \eqref{def of rtau ODE}, we have 
\begin{equation}
\label{def of summation in 4 S3}
\sum_{j=g_{1,+}+1}^{n-1}
\mathcal{O}\Bigl(\frac{(r_{\tau}-\rho)^{-3}}{n^2}\Bigr)
=
\mathcal{O}\Bigl(\frac{(\log n)^{\frac{3}{8}}}{n^{\frac{7}{8}}}\Bigr). 
\end{equation}
By Lemma~\ref{lemma:Euler-Maclaurin}, we have 
\begin{align*}
\sum_{j=g_{1,+}+1}^{n-1}a\log(r_{\tau}-\rho)
&=  
\int_{g_{1,+}}^{n}a\log(r_{\tau(t)}-\rho)\,dt
-\frac{a}{2}\log(r_1-\rho)-\frac{a}{2}\log(r_{\tau(g_{1,+})}-\rho)
+\mathcal{O}(n^{-\frac{3}{8}}). 
\end{align*} 
Note that by \eqref{def of rtau g1plus}, as $n\to+\infty$, we have 
\begin{align*}
\int_{g_{1,+}}^{n}a\log(r_{\tau(t)}-\rho)\,dt
&=n\,a\int_{\rho}^{r_1}\log (u-\rho)\,2u\Delta Q(u)\,du
\\
&\quad
-a\rho\sqrt{\Delta Q(\rho)}
\Bigl( 
\log\Delta_{n,+}-\log 2-\frac{1}{2}\log\Delta Q(\rho)-\frac{1}{2}\log n-1
\Bigr)\sqrt{n}\Delta_{n,+}
\\
&\quad
+\frac{a}{8}\Bigl(1+\frac{\rho\partial_r\Delta Q(\rho)}{\Delta Q(\rho)}\Bigr)\Delta_{n,+}^2
+
\mathcal{O}\Bigl( 
\frac{\Delta_{n,+}^3}{\sqrt{n}}
\Bigr).
\end{align*}
Thus, the Taylor theorem together with \eqref{def of rtau g1plus} gives rise to
\begin{align}
\begin{split}
\label{def of summation in 1 S3}
\sum_{j=g_{1,+}+1}^{n-1}a\log(r_{\tau}-\rho)
&=  
n\,a\int_{r_{\tau}(g_{1,+})}^{r_1}\log (u-\rho)\,2u\Delta Q(u)\,du
\\
&\quad
+\frac{a}{2}\rho\sqrt{\Delta Q(\rho)}
\Bigl( 
2\log 2+2+\log\Delta Q(\rho)+\log n
-2\log\Delta_{n,+}
\Bigr)\sqrt{n}\Delta_{n,+}
\\
&\quad
+\frac{a}{8}\Bigl(1+\frac{\rho\partial_r\Delta Q(\rho)}{\Delta Q(\rho)}\Bigr)\Delta_{n,+}^2
-\frac{a}{2}\log(r_1-\rho)
-\frac{a}{2}\log(\Delta_{n,+})
+\frac{a}{2}\log2
\\
&\quad
+\frac{a}{4}\log(\Delta Q(\rho))
+\frac{a}{4}\log n
+
\mathcal{O}\Bigl( 
\frac{1}{\Delta_{n,+}}
\Bigr). 
\end{split}
\end{align}
The term of order $n^{-1}$ can be calculated as
\begin{align}
\begin{split}
\label{def of summation in 2 S3}
&\quad
\frac{1}{n}
\sum_{j=g_{1,+}+1}^{n-1}
\Bigl\{ 
\frac{a(a-1)(r_{\tau}-\rho)^{-2}}{8\Delta Q(r_{\tau})}
+\frac{a(r_{\tau}-\rho)^{-1}}{8\Delta Q(r_{\tau})}
\Bigl(-\frac{\partial_r\Delta Q(r_{\tau})}{\Delta Q(r_{\tau})}
+\frac{4\alpha+3}{r_{\tau}}
\Bigr)
\Bigr\}
\\
&=
\frac{a(a-1)}{4}\bigg[\frac{1}{2}\Bigl(1+\frac{\rho\,\partial_r\Delta Q(\rho)}{\Delta Q(\rho)}\Bigr)+\log(2(r_1-\rho)\sqrt{\Delta Q(\rho)})-\frac{\rho}{r_1-\rho}\bigg]
\\
&\quad
+\frac{a}{4}\Bigl(4\alpha+3-\frac{\rho\,\partial_r\Delta Q(\rho)}{\Delta Q(\rho)}\Bigr)\log(2(r_1-\rho)\sqrt{\Delta Q(\rho)})
\\
&\quad
-\frac{a}{4}\int_{\rho}^{r_1}\Bigl(\frac{1}{u-\rho}\frac{u\,\partial_u\Delta Q(u)}{\Delta Q(u)}-\frac{1}{u-\rho}\frac{\rho\,\partial_r\Delta Q(\rho)}{\Delta Q(\rho)}\Bigr)\,du
\\
&\quad
+\frac{a(a-1)}{2}\rho\sqrt{\Delta Q(\rho)}\frac{\sqrt{n}}{\Delta_{n,+}}
+\frac{a}{4}\Bigl(\frac{1}{2}\log n-\log\Delta_{n,+}\Bigr)\Bigl(4\alpha+a+2-\frac{\rho\,\partial_r\Delta Q(\rho)}{\Delta Q(\rho)}\Bigr)+\mathcal{O}(\Delta_{n,+}^{-2}). 
\end{split}
\end{align}
Finally, it is straightforward to see that by \eqref{def of rtau g1plus} and Lemma~\ref{lemma:Euler-Maclaurin}, 
\begin{align}
\label{def of summation in 3 S3}
-\frac{1}{n^2}
\sum_{j=g_{1,+}+1}^{n-1}
\frac{a(a-1)(2a-3)}{64\Delta Q(r_{\tau})^2(r_{\tau}-\rho)^4}
&=
-\frac{a(a-1)(2a-3)}{12}\rho\sqrt{\Delta Q(\rho)}\frac{\sqrt{n}}{\Delta_{n,+}^3}
+\mathcal{O}\Bigl(\frac{(\log n)^{1/4}}{n^{1/4}}\Bigr). 
\end{align}
Combining \eqref{def of summation in 4 S3} with \eqref{def of summation in 1 S3}, \eqref{def of summation in 2 S3}, and \eqref{def of summation in 3 S3}, we obtain \eqref{def of S3 asymptotics}.
\end{proof}









   



\section{Local asymptotic analysis for the proof of Theorem~\ref{theorem:calEn expansion}}
\label{section:local analysis part}
\subsection{Asymptotic expansion of \texorpdfstring{$S_2$}{S_2}}
\label{subsection:local analysis S2}
We begin with the asymptotic expansion of the summand \eqref{def of log calEn splitted}. 
We recall \eqref{def of parabolic cylinder function} here. 
\begin{lemma}\label{lemma:S2in part1}
Let
\begin{equation}
\label{def of xi}
     \xi\equiv\xi_{j}:=\frac{\sqrt{n}}{\sqrt{\Delta Q(\rho)}}\frac{\tau-\tau_{\rho}}{\rho},
\end{equation}
where $\tau_{\rho}$ is given by \eqref{def of tau rho}. 
There exists $\delta>0$ such that as $n\to+\infty$, we have 
\begin{align}
\begin{split}
\label{def of S2 asymptotics part 1}
S_2
&=
\sum_{j=g_{1,-}}^{g_{1,+}}\log\Bigl(\frac{1}{(4n)^{\frac{a}{2}}}\Bigr)
+
\frac{a}{2}\sum_{j=g_{1,-}}^{g_{1,+}}\log\Bigl(\frac{1}{\Delta Q(\rho)}\Bigr)
+
\sum_{j=g_{1,-}}^{g_{1,+}}\log \bigl[\mathcal{H}_{a,u}(\xi)\bigr]
\\
&\quad
+\frac{1}{\sqrt{n}} \sum_{j=g_{1,-}}^{g_{1,+}}\frac{\widetilde{g}_{a,u}(\xi)}{12\rho\sqrt{\Delta Q(\rho)}g_{a,u}(\xi)}
+\mathcal{O}\Bigl(\frac{(\log n)^{\frac{7}{8}}}{n^{\frac{1}{8}}}\Bigr),
\end{split}
\end{align}    
uniformly for $u\in\{z\in\C:|z-x|\leq \delta\}$ and $a$ in compact subsets of $(-1,+\infty)$.
Here, $\mathcal{H}_{a,u}(\xi),g_{a,u}(\xi)$ are given by \eqref{def of cal H au}, and 
\begin{align}
\begin{split}
\label{def of tilde gau}
\widetilde{g}_{a,u}(\xi)&:=
-a\xi\Bigl(1-\frac{\rho\,\partial_r\Delta Q(\rho)}{\Delta Q(\rho)}\Bigr)
(e^{u}D_{-a-1}(\xi)+D_{-a-1}(-\xi))
\\
&\quad
+(e^u D_{-a}(\xi)-D_{-a}(-\xi))
\xi^2\Bigl(2+\frac{\rho\,\partial_r\Delta Q(\rho)}{\Delta Q(\rho)}\Bigr)
\\
&\quad
-(e^u D_{-a}(\xi)-D_{-a}(-\xi))
\bigg[
6(\rho \mathsf{k}'(\rho)+1)+(2+a)\Bigl(1-\frac{\rho\,\partial_r\Delta Q(\rho)}{\Delta Q(\rho)}\Bigr)
\bigg].
\end{split}
\end{align}
\end{lemma}
\begin{proof}
We begin with analyzing $S_2^{(\mathrm{in})}$. 
There exists $c>0$ such that 
\begin{align}
\begin{split}
\label{def of hin rho 1}
h_{n,j}^{(\mathrm{in})}(\rho)
&=
\int_{r_{\tau}-\delta_n'}^{\rho}
2r e^{\mathsf{k}(r)}e^{-nV_{\tau}(r)}|\rho-r|^{\alpha}\,dr
+\int_{0}^{r_{\tau}-\delta_n'}
2r e^{\mathsf{k}(r)}e^{-nV_{\tau}(r)}|\rho-r|^{\alpha}\,dr
\\
%&=\int_{r_{\tau}-\delta_n'}^{\rho}
%2r e^{\mathsf{k}(r)}e^{-nV_{\tau}(r)}(\rho-r)^{\alpha}\,dr
%+e^{-nV_{\tau}(r_{\tau})}\cdot\mathcal{O}\bigl(e^{-cM^2}\bigr)
%\\
&=
\frac{2r_{\tau}e^{\mathsf{k}(r_{\tau})}e^{-nV_{\tau}(r_{\tau})}}{(nd_2)^{\frac{a+1}{2}}}
\int_{0}^{\sqrt{d_2n}(\rho-r_{\tau}+\delta_n')}
v^{a}
e^{-\frac{1}{2}(v-\sqrt{nd_2}(\rho-r_{\tau}))^2}
\\
&\qquad\times
\Bigl( 
1+\frac{c_1(u_v)}{\sqrt{n}}+\frac{c_2(u_v)}{n}+\frac{c_3(u_v)}{n^{3/2}}+\mathcal{O}(\frac{c_4(u_v)}{n^2})
\Bigr)\,dv
+e^{-nV_{\tau}(r_{\tau})}\cdot\mathcal{O}\bigl(e^{-cM^2}\bigr), 
\end{split}
\end{align}
where $u_v=\sqrt{nd_2}(\rho-r_{\tau})-v=\chi(\tau)-v$ with $\chi(\tau):=\sqrt{nd_2}(\rho-r_{\tau})$.
Since $\sqrt{d_2n}(\rho-r_{\tau}+\delta_n')>\sqrt{nd_2}(\rho-r_{\tau})$ for $g_{1,-}\leq j\leq \lfloor \tau_{\rho_1}\rfloor-1$, one can safely extend the integral region to $[0,+\infty)$ with an error $\mathcal{O}(e^{-cM^2})$ for some $c>0$. 
Note that for $a> -1$ and $k\in\mathbb{Z}_{\geq 0}$, by \eqref{def of parabolic cylinder function}, we have 
\begin{equation}
\int_{0}^{+\infty}
v^{a+k}e^{-\frac{1}{2}(v-\chi(\tau))^2}\,dv
=
\vartheta_{k,a}
e^{-\frac{\chi(\tau)^2}{4}}D_{-a-k-1}(-\chi(\tau)). 
\end{equation}
Here, $\vartheta_{k,a}:=\Gamma(a+k+1)$.
Therefore, we have 
\begin{align}
\begin{split}
\label{def of hin rho 2}
&\quad 
\int_{0}^{\sqrt{d_2n}(\rho-r_{\tau}+\delta_n')}
v^{a}
e^{-\frac{1}{2}(v-\sqrt{nd_2}(\rho-r_{\tau}))^2}
\Bigl( 
1+\frac{c_1(u_v)}{\sqrt{n}}+\frac{c_2(u_v)}{n}+\frac{c_3(u_v)}{n^{3/2}}+\mathcal{O}(\frac{c_4(u_v)}{n^2})
\Bigr)\,dv
\\
&=
\mathsf{c}_0+\frac{\mathsf{c}_1}{\sqrt{n}}+\frac{\mathsf{c}_2}{n}+\frac{\mathsf{c}_3}{n^{\frac{3}{2}}}+\mathcal{O}\Bigl(\frac{\mathsf{c}_4}{n^2}\Bigr),
\end{split}
\end{align}
where $c_1,c_2,c_3$ are given by \eqref{def of c_1}, \eqref{def of c_2}, and \eqref{def of c_3}, respectively, and 
\begin{equation}
\label{def of sfc0}
\mathsf{c}_0
:= 
\vartheta_{0,a}
e^{-\frac{\chi(\tau)^2}{4}}D_{-a-1}(-\chi(\tau))
=
\vartheta_{0,a}
e^{-\frac{\phi(\tau)^2}{4}}D_{-a-1}(\phi(\tau)), 
\end{equation}
\begin{align}
\begin{split}
\label{def of sfc1}
\mathsf{c}_1
&:=
\int_{0}^{+\infty}
v^{a}
e^{-\frac{1}{2}(v-\sqrt{nd_2}(\rho-r_{\tau}))^2}c_1(u_v)\,dv 
\\
&\,\,\,=
\frac{d_3}{6d_2^{3/2}}
\sum_{\ell=0}^{3}\binom{3}{\ell}\phi(\tau)^{3-\ell}\vartheta_{\ell,a}e^{-\frac{\phi(\tau)^2}{4}}D_{-a-\ell-1}(\phi(\tau))
\\
&\quad
-
\frac{1}{d_2^{1/2}}\Bigl(\mathsf{k}'(r_{\tau})+\frac{1}{r_{\tau}}\Bigr)
\sum_{\ell=0}^{1}\binom{1}{\ell}\phi(\tau)^{1-\ell}\vartheta_{\ell,a}e^{-\frac{\phi(\tau)^2}{4}}D_{-a-\ell-1}(\phi(\tau)), 
\end{split}
\end{align}
\begin{align}
\begin{split}
\label{def of sfc2}
\mathsf{c}_2
&:=\int_{0}^{+\infty}
v^{a}
e^{-\frac{1}{2}(v-\sqrt{nd_2}(\rho-r_{\tau}))^2}c_2(u_v)\,dv  
\\
&\,\,\,=
\frac{d_3^2}{72d_2^3}\sum_{\ell=0}^{6}\binom{6}{\ell}\phi(\tau)^{6-\ell}\vartheta_{\ell,a}e^{-\frac{\chi(\tau)^2}{4}}D_{-a-\ell-1}(\phi(\tau))
\\
&\quad
-\frac{1}{24d_2^2}
\Bigl(d_4+4d_3\mathsf{k}'(r_{\tau})+\frac{4d_3}{r_{\tau}}\Bigr)
\sum_{\ell=0}^{4}\binom{4}{\ell}\phi(\tau)^{4-\ell}
\vartheta_{\ell,a}e^{-\frac{\phi(\tau)^2}{4}}D_{-a-\ell-1}(\phi(\tau))
\\
&\quad
+\frac{1}{d_2}\Bigl(\frac{\mathsf{k}''(r_{\tau})}{2}+\frac{\mathsf{k}'(r_{\tau})^2}{2}+\frac{\mathsf{k}'(r_{\tau})}{r_{\tau}}\Bigr)
\sum_{\ell=0}^{2}\binom{2}{\ell}\phi(\tau)^{2-\ell}
\vartheta_{\ell,a}e^{-\frac{\phi(\tau)^2}{4}}D_{-a-\ell-1}(\phi(\tau)), 
\end{split}
\end{align}
\begin{align}
\begin{split}
\label{def of sfc3}
\mathsf{c}_3
&:=\int_{0}^{+\infty}
v^{a}
e^{-\frac{1}{2}(v-\sqrt{nd_2}(\rho-r_{\tau}))^2}c_3(u_v)\,dv  
\\
&\,\,\,=
\frac{d_3^3}{1296d_2^{9/2}}
\sum_{\ell=0}^{9}\binom{9}{\ell}\phi(\tau)^{9-\ell}
\vartheta_{\ell,a}e^{-\frac{\phi(\tau)^2}{4}}D_{-a-\ell-1}(\phi(\tau))
\\
&\quad
-\Bigl(d_3d_4+2d_3^2\mathsf{k}'(r_{\tau})+\frac{2d_3^2}{r_{\tau}} \Bigr)\frac{1}{144d_2^{7/2}}
\sum_{\ell=0}^{7}\binom{7}{\ell}\phi(\tau)^{7-\ell}
\vartheta_{\ell,a}e^{-\frac{\phi(\tau)^2}{4}}D_{-a-\ell-1}(\phi(\tau))
\\
&\quad
+
\Bigl(
d_5+\frac{5}{r_{\tau}}d_4+5d_4\mathsf{k}'(r_{\tau})+\frac{20d_3}{r_{\tau}}\mathsf{k}'(r_{\tau})+10d_3\mathsf{k}'(r_{\tau})^2+10d_3\mathsf{k}''(r_{\tau})
\Bigr)
\sum_{\ell=0}^{5}\binom{5}{\ell}
\frac{\phi(\tau)^{5-\ell}\vartheta_{\ell,a}e^{-\frac{\phi(\tau)^2}{4}}}{120d_2^{5/2}}D_{-a-\ell-1}(\phi(\tau))
\\
&\quad
-
\Bigl(\frac{3\mathsf{k}'(r_{\tau})^2}{r_{\tau}}
+\frac{3\mathsf{k}''(r_{\tau})}{r_{\tau}}+\mathsf{k}'(r_{\tau})^3+3\mathsf{k}'(r_{\tau})\mathsf{k}''(r_{\tau})+\mathsf{k}^{(3)}(r_{\tau})\Bigr)
\sum_{\ell=0}^{3}\binom{3}{\ell}
\frac{\phi(\tau)^{3-\ell}
\vartheta_{\ell,a}e^{-\frac{\phi(\tau)^2}{4}}}{6d_2^{3/2}}D_{-a-\ell-1}(\phi(\tau)).
\end{split}
\end{align}
We next consider $h_{n,j}^{(\mathrm{out})}(\rho)$. 
In this integral region and for $g_{1,-}\leq j\leq \lfloor n\tau_{\rho} \rfloor -1$, there is no critical point, but we still proceed with the expansion in terms of $r_{\tau}$. Indeed, 
\[
h_{n,j}^{(\mathrm{out})}(\rho)=\int_{\rho}^{+\infty} 2ve^{\mathsf{k}(v)}e^{-nV_{\tau}(v)}|v-\rho|^{a}\,dv =
\int_{0}^{+\infty} 2ve^{\mathsf{k}(v)}e^{-nV_{\tau}(v)}|v-\rho|^{a}\,dv
-
\int_0^{\rho}2ve^{\mathsf{k}(v)}e^{-nV_{\tau}(v)}|v-\rho|^{a}\,dv.
\]
We already know the asymptotics of the second terms in the last line. 
For the first term, we split the integral $[0,+\infty)$ into $[0,+\infty)=\{r\in[0,+\infty): |r-r_{\tau}|<\delta_n'\}\cup\{r\in[0,+\infty): |r-r_{\tau}|\geq\delta_n'\}$. 
In the later region, the corresponding integral can be neglected with an error $e^{-nV_{\tau}(r_{\tau})}\cdot \mathcal{O}(e^{-cM^2})$ for some $c>0$. 
The Gaussian integral gives 
\begin{align*}
&\quad 
\frac{2r_{\tau}e^{\mathsf{k}(r_{\tau})}e^{-nV_{\tau}(r_{\tau})}}{(nd_2)^{\frac{1}{2}}}
\int_{-\sqrt{d_2n}\delta_n'}^{\sqrt{d_2n}\delta_n'}
e^{-\frac{1}{2}u^2}
\Bigl|r_{\tau}-\rho+\frac{u}{\sqrt{nd_2}}\Bigr|^{a}
\Bigl( 
1+\frac{c_1(u)}{\sqrt{n}}+\frac{c_2(u)}{n}+\frac{c_3(u)}{n^{3/2}}+\mathcal{O}(\frac{c_4(u)}{n^2})
\Bigr)\,du
\\
&=
\frac{2r_{\tau}e^{\mathsf{k}(r_{\tau})}e^{-nV_{\tau}(r_{\tau})}}{(nd_2)^{\frac{a+1}{2}}}
\int_{-\sqrt{d_2n}\delta_n'+\sqrt{d_2n}(r_{\tau}-\rho)}^{\sqrt{nd_2}\delta_n'+\sqrt{nd_2}(r_{\tau}-\rho)}
e^{-\frac{1}{2}(v-\sqrt{nd_2}(r_{\tau}-\rho))^2}
|v|^{a}
\Bigl( 
1+\frac{c_1(u_v)}{\sqrt{n}}+\frac{c_2(u_v)}{n}+\frac{c_3(u_v)}{n^{3/2}}+\mathcal{O}(\frac{c_4(u_v)}{n^2})
\Bigr)\,dv,
\end{align*}
where $u_v=v-\sqrt{nd_2}(r_{\tau}-\rho)$. 
One can safely extend the integral region to $(-\infty,+\infty)$ with an exponential error $\mathcal{O}(e^{-cM^2})$ for some $c>0$. 
Each integral is given by the following: for $\phi(\tau):=\sqrt{nd_2}(r_{\tau}-\rho)$, 
\begin{align*}
&\quad\int_{-\infty}^{+\infty}
e^{-\frac{1}{2}(v-\sqrt{nd_2}(r_{\tau}-\rho))^2}
|v|^{a}\,dv
\\
&=
\int_{0}^{+\infty}
e^{-\frac{1}{2}(v-\phi(\tau))^2}
v^{a}\,du
+\int_{0}^{+\infty}
e^{-\frac{1}{2}(v+\phi(\tau))^2}
v^{a}\,du
=\vartheta_{0,a}e^{-\frac{1}{4}\phi(\tau)^2}\bigl( 
D_{-a-1}(-\phi(\tau))+D_{-a-1}(\phi(\tau))
\bigr). 
\end{align*}
Since by \eqref{def of parabolic cylinder function}, 
\begin{align*}
\int_{0}^{+\infty}t^{a}e^{-\frac{1}{2}(t+z)^2}\,dt
&=\Gamma(a+1)e^{-\frac{1}{4}z^2}D_{-a-1}(z), 
\\
\int_{-\infty}^{+\infty}e^{-\frac{1}{2}(u-\phi(\tau))^2}|v|^{a}v^{\ell}\,dv
&=
\Gamma(a+\ell+1)e^{-\frac{1}{4}\phi(\tau)^2}
\Bigl( 
(-1)^{\ell}D_{-a-\ell-1}(\phi(\tau))
+
D_{-a-\ell-1}(-\phi(\tau))
\Bigr), 
\end{align*}
we have 
\begin{equation}
\label{def of sfcd0}
\mathsf{c}_0'=   
\int_{-\infty}^{+\infty}e^{-\frac{1}{2}(u-\phi(\tau))^2}|v|^{a}\,dv
=
\vartheta_{0,a}e^{-\frac{1}{4}\phi(\tau)^2}
\Bigl( 
D_{-a-1}(\phi(\tau))
+
D_{-a-1}(-\phi(\tau))
\Bigr), 
\end{equation}
\begin{align}
\begin{split}
\label{def of sfcd1}
\mathsf{c}_1'&=\int_{-\infty}^{+\infty}
e^{-\frac{1}{2}(v-\sqrt{nd_2}(r_{\tau}-\rho))^2}
|v|^{a}
c_1(u_v)\,dv
\\
&=
-\frac{d_3}{6d_2^{3/2}}\sum_{\ell=0}^{3}\binom{3}{\ell}(-\phi(\tau))^{3-\ell}
\vartheta_{\ell,a}e^{-\frac{1}{4}\phi(\tau)^2}
\Bigl( 
(-1)^{\ell}D_{-a-\ell-1}(\phi(\tau))
+
D_{-a-\ell-1}(-\phi(\tau))
\Bigr)
\\ &\quad
+\frac{1}{d_2^{1/2}}\Bigl(\mathsf{k}'(r_{\tau})+\frac{1}{r_{\tau}}\Bigr)
\sum_{\ell=0}^{1}\binom{1}{\ell}(-\phi(\tau))^{1-\ell}
\vartheta_{\ell,a}e^{-\frac{1}{4}\phi(\tau)^2}
\Bigl( 
(-1)^{\ell}D_{-a-\ell-1}(\phi(\tau))
+
D_{-a-\ell-1}(-\phi(\tau))
\Bigr),
\end{split}
\end{align}
\begin{align}
\begin{split}
\label{def of sfcd2}
\mathsf{c}_2'&= \int_{-\infty}^{+\infty}
e^{-\frac{1}{2}(v-\sqrt{nd_2}(r_{\tau}-\rho))^2}
|v|^{a}
c_2(u_v)\,du
\\
&=
\frac{d_3^2}{72d_2^3}
\sum_{\ell=0}^{6}\binom{6}{\ell}(-\phi(\tau))^{6-\ell}
\vartheta_{\ell,a}e^{-\frac{1}{4}\phi(\tau)^2}
\Bigl( 
(-1)^{\ell}D_{-a-\ell-1}(\phi(\tau))
+
D_{-a-\ell-1}(-\phi(\tau))
\Bigr)
\\
&\quad
-
\frac{1}{24d_2^2}\Bigl(d_4+4d_3\mathsf{k}'(r_{\tau})+\frac{4d_3}{r_{\tau}}\Bigr)
\sum_{\ell=0}^{4}\binom{4}{\ell}(-\phi(\tau))^{4-\ell}
\vartheta_{\ell,a}e^{-\frac{1}{4}\phi(\tau)^2}
\Bigl( 
(-1)^{\ell}D_{-a-\ell-1}(\phi(\tau))
+
D_{-a-\ell-1}(-\phi(\tau))
\Bigr)
\\
&\quad
+
\frac{1}{d_2}\Bigl(\frac{\mathsf{k}''(r_{\tau})}{2}+\frac{\mathsf{k}'(r_{\tau})^2}{2}+\frac{\mathsf{k}'(r_{\tau})}{r_{\tau}}\Bigr)
\sum_{\ell=0}^{2}\binom{2}{\ell}(-\phi(\tau))^{2-\ell}
\vartheta_{\ell,a}e^{-\frac{1}{4}\phi(\tau)^2}
\Bigl( 
(-1)^{\ell}D_{-a-\ell-1}(\phi(\tau))
+
D_{-a-\ell-1}(-\phi(\tau))
\Bigr). 
\end{split}
\end{align}
Note that $\mathsf{c}'_3$ is given as well as \eqref{def of sfc3}, but we omit the explicit expression. 
Therefore, we obtain 
\begin{align}
\begin{split}
\label{def of hrho out 2}
h_{n,j}^{(\mathrm{out})}(\rho)
&=
\int_{0}^{+\infty} 2ve^{\mathsf{k}(v)}e^{-nV_{\tau}(v)}|v-\rho|^{a}\,dv
-
\int_0^{\rho}2ve^{\mathsf{k}(v)}e^{-nV_{\tau}(v)}|v-\rho|^{a}\,dv
\\
&=
\frac{2r_{\tau}e^{\mathsf{k}(r_{\tau})}e^{-nV_{\tau}(r_{\tau})}}{(nd_2)^{\frac{a+1}{2}}}
\Bigl( 
\mathsf{c}_0'-\mathsf{c}_0+\frac{\mathsf{c}_1'-\mathsf{c}_1}{\sqrt{n}}+\frac{\mathsf{c}_2'-\mathsf{c}_2}{n}+\mathcal{O}\Bigl(\frac{\mathsf{c}_3'-\mathsf{c}_3}{n^{3/2}}\Bigr)
\Bigr),
\end{split}
\end{align}
where by \eqref{def of sfc0}, \eqref{def of sfc1}, \eqref{def of sfc2}, \eqref{def of sfc3}, \eqref{def of sfcd0}, \eqref{def of sfcd1}, and \eqref{def of sfcd2}, 
\begin{align*}
\mathsf{c}_0'-\mathsf{c}_0
%&=
%\vartheta_{0,a}e^{-\frac{1}{4}\phi(\tau)^2}
%\Bigl( 
%D_{-a-1}(\phi(\tau))
%+
%D_{-a-1}(-\phi(\tau))
%\Bigr)-\vartheta_{0,a}
%e^{-\frac{\phi(\tau)^2}{4}}D_{-a-1}(\phi(\tau))
&=
\vartheta_{0,a}e^{-\frac{1}{4}\phi(\tau)^2}
D_{-a-1}(-\phi(\tau)), 
\\
\mathsf{c}_1'-\mathsf{c}_1
&=
-\frac{d_3}{6d_2^{3/2}}\sum_{\ell=0}^{3}\binom{3}{\ell}(-\phi(\tau))^{3-\ell}
\vartheta_{\ell,a}e^{-\frac{1}{4}\phi(\tau)^2}
D_{-a-\ell-1}(-\phi(\tau))
\\ &\quad
+\frac{1}{d_2^{1/2}}\Bigl(\mathsf{k}'(r_{\tau})+\frac{1}{r_{\tau}}\Bigr)
\sum_{\ell=0}^{1}\binom{1}{\ell}(-\phi(\tau))^{1-\ell}
\vartheta_{\ell,a}e^{-\frac{1}{4}\phi(\tau)^2}
D_{-a-\ell-1}(-\phi(\tau)), 
\\
\mathsf{c}_2'-\mathsf{c}_2
&=    
\frac{d_3^2}{72d_2^3}
\sum_{\ell=0}^{6}\binom{6}{\ell}(-\phi(\tau))^{6-\ell}
\vartheta_{\ell,a}e^{-\frac{1}{4}\phi(\tau)^2}
D_{-a-\ell-1}(-\phi(\tau))
\\
&\quad
-
\frac{1}{24d_2^2}\Bigl(d_4+4d_3\mathsf{k}'(r_{\tau})+\frac{4d_3}{r_{\tau}}\Bigr)
\sum_{\ell=0}^{4}\binom{4}{\ell}(-\phi(\tau))^{4-\ell}
\vartheta_{\ell,a}e^{-\frac{1}{4}\phi(\tau)^2}
D_{-a-\ell-1}(-\phi(\tau))
\\
&\quad
+
\frac{1}{d_2}\Bigl(\frac{\mathsf{k}''(r_{\tau})}{2}+\frac{\mathsf{k}'(r_{\tau})^2}{2}+\frac{\mathsf{k}'(r_{\tau})}{r_{\tau}}\Bigr)
\sum_{\ell=0}^{2}\binom{2}{\ell}(-\phi(\tau))^{2-\ell}
\vartheta_{\ell,a}e^{-\frac{1}{4}\phi(\tau)^2}
D_{-a-\ell-1}(-\phi(\tau)). 
\end{align*}
By \eqref{def of hin rho 1}, \eqref{def of hin rho 2}, and \eqref{def of hrho out 2}, let us choose $\delta>0$ sufficiently small so that $e^{u}\frac{h_{n,j}^{(\mathrm{in})}(\rho)}{h_{n,j}}+\frac{h_{n,j}^{(\mathrm{out})}(\rho)}{h_{n,j}}$ remains bounded away from the interval $(-\infty,0]$ as $n\to+\infty$ uniformly for $u\in\{z\in\C:|z-x|\leq \delta\}$ and $a$ in compact subsets of $(-1,+\infty)$. 
Then, we have 
\begin{align}
\begin{split}
\label{def of S2 first part}    
\sum_{j=g_{1,-}}^{\lfloor n\tau_{\rho}\rfloor -1}
\log\Bigl(e^{u}\frac{h_{n,j}^{(\mathrm{in})}(\rho)}{h_{n,j}}+\frac{h_{n,j}^{(\mathrm{out})}(\rho)}{h_{n,j}}\Bigr)
&=
\sum_{j=g_{1,-}}^{\lfloor n\tau_{\rho}\rfloor -1}\log\Bigl(\frac{1}{(4n)^{\frac{a}{2}}}\Bigr)
+
\frac{a}{2}\sum_{j=g_{1,-}}^{\lfloor n\tau_{\rho}\rfloor -1}\log\Bigl(\frac{1}{\Delta Q(\rho)}\Bigr)
+
\sum_{j=g_{1,-}}^{\lfloor n\tau_{\rho}\rfloor -1}\log \bigl[\mathcal{H}_{a,u}(\xi)\bigr]
\\
&\quad
+\frac{1}{\sqrt{n}} \sum_{j=g_{1,-}}^{\lfloor n\tau_{\rho}\rfloor -1}\frac{\widetilde{g}_{a,u}(\xi)}{12\rho\sqrt{\Delta Q(\rho)}g_{a,u}(\xi)}
+\mathcal{O}\Bigl(\frac{(\log n)^{\frac{7}{8}}}{n^{\frac{1}{8}}}\Bigr),
\end{split}
\end{align}
uniformly for $u\in\{z\in\C:|z-x|\leq \delta\}$ and $a$ in compact subsets of $(-1,+\infty)$.
Here, $g_{a,u}(x),\widetilde{g}_{a,u}(x)$ are given by \eqref{def of cal H au} and \eqref{def of tilde gau}, respectively. 
\begin{comment}
\begin{align}
g_0(\xi)&:=
e^{u}D_{-a-1}(\xi)+D_{-a-1}(-\xi),
\\
\label{def of tilde gau}
\widetilde{g}_{a,u}(\xi)&:=
-a\xi\Bigl(1-\frac{\rho\,\partial_r\Delta Q(\rho)}{\Delta Q(\rho)}\Bigr)g_0(\xi)
+(e^u D_{-a}(\xi)-D_{-a}(-\xi))
\xi^2\Bigl(2+\frac{\rho\,\partial_r\Delta Q(\rho)}{\Delta Q(\rho)}\Bigr)
\\
&\quad
-(e^u D_{-a}(\xi)-D_{-a}(-\xi))
\bigg[
6(\rho \mathsf{k}'(\rho)+1)+(2+a)\Bigl(1-\frac{\rho\,\partial_r\Delta Q(\rho)}{\Delta Q(\rho)}\Bigr)
\bigg].
\end{align}
\end{comment}

We consider the case $\lfloor n\tau_{\rho}\rfloor\leq j\leq g_{1,+}$. 
Similar manner to \eqref{def of hrho out 2}, there exists $c>0$ such that we have 
\begin{align*}
h_{n,j}^{(\mathrm{out})}(\rho)
&=
\frac{2r_{\tau}e^{\mathsf{k}(r_{\tau})}e^{-nV_{\tau}(r_{\tau})}}{(nd_2)^{\frac{a+1}{2}}}
\int_{0}^{\sqrt{nd_2}\delta_n'+\sqrt{nd_2}(r_{\tau}-\rho)}
y^a
e^{-\frac{1}{2}(y-\sqrt{nd_2}(r_{\tau}-\rho))^2}
\\
&\quad\times
\Bigl( 
1+\frac{c_1(u_y)}{\sqrt{n}}+\frac{c_2(u_y)}{n}+\frac{c_3(u)}{n^{3/2}}+\mathcal{O}(\frac{c_4(u_y)}{n^2})
\Bigr)\,dy
+\mathcal{O}(e^{-cM^2}),
\end{align*}
where $u_y:=y-\sqrt{nd_2}(r_{\tau}-\rho)=y-\phi(\tau)$ with $\phi(\tau):=\sqrt{nd_2}(r_{\tau}-\rho)$. 
Since $\sqrt{nd_2}\delta_n'+\sqrt{nd_2}(r_{\tau}-\rho)>\sqrt{nd_2}(r_{\tau}-\rho)$ for $\lfloor n\tau_{\rho_1}\rfloor\leq j \leq g_{1,+}$, we can extend the integral region to $[0+\infty)$ with an exponential error $\mathcal{O}(e^{-cM^2})$. 
The Gaussian integral gives 
\begin{equation}
\label{def of hrho out second part}
h_{n,j}^{(\mathrm{out})}(\rho)
%&=
%\frac{2r_{\tau}e^{\mathsf{k}(r_{\tau})}e^{-nV_{\tau}(r_{\tau})}}{\sqrt{nd_2}}
%\int_{\sqrt{d_2n}(\rho-r_{\tau})}^{\sqrt{nd_2}\delta_n'}
%\Bigl(r_{\tau}-\rho+\frac{u}{\sqrt{nd_2}}\Bigr)^a
%e^{-\frac{1}{2}u^2}
%\\
%&\quad\times
%\Bigl( 
%1+\frac{c_1(u)}{\sqrt{n}}+\frac{c_2(u)}{n}+\frac{c_3(u)}{n^{3/2}}+\mathcal{O}(\frac{c_4(u)}{n^2})
%\Bigr)\,du
%+\mathcal{O}(e^{-cM^2})
%\\
=
\frac{2r_{\tau}e^{\mathsf{k}(r_{\tau})}e^{-nV_{\tau}(r_{\tau})}}{(nd_2)^{\frac{a+1}{2}}}
\Bigl( 
\widetilde{\mathsf{c}}_0
+\frac{\widetilde{\mathsf{c}}_1}{\sqrt{n}}
+\frac{\widetilde{\mathsf{c}}_2}{n}
+\frac{\widetilde{\mathsf{c}}_3}{n^{\frac{3}{2}}}
+\mathcal{O}\Bigl(\frac{\widetilde{\mathsf{c}}_4}{n^2}\Bigr)
\Bigr)
+\mathcal{O}(e^{-cM^2}),
\end{equation}
where 
%and since $(y-\phi(\tau))^m=\sum_{\ell=0}^{m}\binom{m}{\ell}(-\phi(\tau))^{m-\ell}y^{\ell}$, we have 
\begin{align*}
\widetilde{\mathsf{c}}_0
&=
\int_{0}^{+\infty}
y^ae^{-\frac{1}{2}(y-\phi(\tau))^2}\,dy
=
\vartheta_{0,a}
e^{-\frac{\phi(\tau)^2}{4}}D_{-a-1}(-\phi(\tau)), 
\\
\widetilde{\mathsf{c}}_1
&=\int_{0}^{+\infty}
y^ae^{-\frac{1}{2}(y-\phi(\tau))^2}c_1(u_y)\,dy
\\
&=
-\frac{d_3}{6d_2^{3/2}}
\sum_{\ell=0}^{3}\binom{3}{\ell}(-\phi(\tau))^{3-\ell}\vartheta_{\ell,a}e^{-\frac{\phi(\tau)^2}{4}}D_{-a-\ell-1}(-\phi(\tau))
\\
&\quad
+
\frac{1}{d_2^{1/2}}\Bigl(\mathsf{k}'(r_{\tau})+\frac{1}{r_{\tau}}\Bigr)
\sum_{\ell=0}^{1}\binom{1}{\ell}(-\phi(\tau))^{1-\ell}\vartheta_{\ell,a}e^{-\frac{\phi(\tau)^2}{4}}D_{-a-\ell-1}(-\phi(\tau)),
\\
\widetilde{\mathsf{c}}_2
&=
\int_{0}^{+\infty}
y^ae^{-\frac{1}{2}(y-\phi(\tau))^2}c_2(u_y)\,dy 
\\
&=
\frac{d_3^2}{72d_2^3}
\sum_{\ell=0}^{6}\binom{6}{\ell}(-\phi(\tau))^{6-\ell}\vartheta_{\ell,a}e^{-\frac{\phi(\tau)^2}{4}}D_{-a-\ell-1}(-\phi(\tau))
\\
&\quad
-\frac{1}{24d_2^2}\Bigl(d_4+4d_3\mathsf{k}'(r_{\tau})+\frac{4d_3}{r_{\tau}}\Bigr)
\sum_{\ell=0}^{4}\binom{4}{\ell}(-\phi(\tau))^{4-\ell}\vartheta_{\ell,a}e^{-\frac{\phi(\tau)^2}{4}}D_{-a-\ell-1}(-\phi(\tau))
\\
&\quad 
+\frac{1}{d_2}\Bigl(\frac{\mathsf{k}''(r_{\tau})}{2}+\frac{\mathsf{k}'(r_{\tau})^2}{2}+\frac{\mathsf{k}'(r_{\tau})}{r_{\tau}}\Bigr)
\sum_{\ell=0}^{2}\binom{2}{\ell}(-\phi(\tau))^{2-\ell}\vartheta_{\ell,a}e^{-\frac{\phi(\tau)^2}{4}}D_{-a-\ell-1}(-\phi(\tau)). 
\end{align*}
On the other hand, 
\begin{align}
\begin{split}
\label{def of hrho in second part}    
h_{n,j}^{(\mathrm{in})}(\rho)
&:=
\int_{0}^{+\infty} 2ve^{\mathsf{k}(v)}e^{-nV_{\tau}(v)}|v-\rho|^{a}\,dv
-
\int_{\rho}^{+\infty} 2ve^{\mathsf{k}(v)}e^{-nV_{\tau}(v)}(v-\rho)
^{a}\,dv
\\
&\,=
\frac{2r_{\tau}e^{\mathsf{k}(r_{\tau})}e^{-nV_{\tau}(r_{\tau})}}{(nd_2)^{\frac{a+1}{2}}}
\Bigl( 
\mathsf{c}_0'-\widetilde{\mathsf{c}}_0+\frac{\mathsf{c}_1'-\widetilde{\mathsf{c}}_1}{\sqrt{n}}+\frac{\mathsf{c}_2'-\widetilde{\mathsf{c}}_2}{n}+\mathcal{O}\Bigl(\frac{\mathsf{c}_3'-\widetilde{\mathsf{c}}_3}{n^{3/2}}\Bigr)
\Bigr),
\end{split}
\end{align}
where 
\begin{align*}
\mathsf{c}_0'-\widetilde{\mathsf{c}}_0
&=
\vartheta_{0,a}e^{-\frac{1}{4}\phi(\tau)^2}
D_{-a-1}(\phi(\tau)), 
\\
\mathsf{c}_1'-\widetilde{\mathsf{c}}_1
&=
-\frac{d_3}{6d_2^{3/2}}\sum_{\ell=0}^{3}\binom{3}{\ell}(-\phi(\tau))^{3-\ell}
\vartheta_{\ell,a}e^{-\frac{1}{4}\phi(\tau)^2}
(-1)^{\ell}D_{-a-\ell-1}(\phi(\tau))
\\ &\quad
+\frac{1}{d_2^{1/2}}\Bigl(\mathsf{k}'(r_{\tau})+\frac{1}{r_{\tau}}\Bigr)
\sum_{\ell=0}^{1}\binom{1}{\ell}(-\phi(\tau))^{1-\ell}
\vartheta_{\ell,a}e^{-\frac{1}{4}\phi(\tau)^2}
(-1)^{\ell}D_{-a-\ell-1}(\phi(\tau)), 
\\
\mathsf{c}_2'-\widetilde{\mathsf{c}}_2
&=
\frac{d_3^2}{72d_2^3}
\sum_{\ell=0}^{6}\binom{6}{\ell}(-\phi(\tau))^{6-\ell}
\vartheta_{\ell,a}e^{-\frac{1}{4}\phi(\tau)^2}
(-1)^{\ell}D_{-a-\ell-1}(\phi(\tau))
\\
&\quad
-
\frac{1}{24d_2^2}\Bigl(d_4+4d_3\mathsf{k}'(r_{\tau})+\frac{4d_3}{r_{\tau}}\Bigr)
\sum_{\ell=0}^{4}\binom{4}{\ell}(-\phi(\tau))^{4-\ell}
\vartheta_{\ell,a}e^{-\frac{1}{4}\phi(\tau)^2} 
(-1)^{\ell}D_{-a-\ell-1}(\phi(\tau))
\\
&\quad
+
\frac{1}{d_2}\Bigl(\frac{\mathsf{k}''(r_{\tau})}{2}+\frac{\mathsf{k}'(r_{\tau})^2}{2}+\frac{\mathsf{k}'(r_{\tau})}{r_{\tau}}\Bigr)
\sum_{\ell=0}^{2}\binom{2}{\ell}(-\phi(\tau))^{2-\ell}
\vartheta_{\ell,a}e^{-\frac{1}{4}\phi(\tau)^2}
(-1)^{\ell}D_{-a-\ell-1}(\phi(\tau)). 
\end{align*}
Rewriting the above expansion in terms of \eqref{def of xi}, by \eqref{def of hrho out second part} and \eqref{def of hrho in second part} we obtain the same expansion with \eqref{def of S2 first part}.  
\begin{comment}

First note that 
\begin{align*}
\phi(\tau)
=\sqrt{nd_2}(r_{\tau}-\rho)    
=\xi_{\rho}-\mathsf{q}_1(\rho)\frac{\xi_{\rho}^2}{\sqrt{n}}
-
\mathsf{q}_2(\rho)
\frac{\xi_{\rho}^3}{n}
+\mathcal{O}\Bigl(\frac{\xi_{\rho}^4}{n^{\frac{3}{2}}}\Bigr).
\end{align*}

\end{comment}

%Since 
%\[
%D_{1-a}(\xi)=\xi D_{-a}(\xi)+aD_{-a-1}(\xi),
%\]
\end{proof}
To apply Lemma~\ref{lemma:Euler-Maclaurin} to each sum in Lemma~\ref{lemma:S2in part1}, we will use
\begin{equation}
\label{def of d relationship between g0 g1}
\partial_x\log\bigl[\mathcal{H}_{a,u}(x)\bigr]
=
-\frac{e^{u}D_{-a}(x)-D_{-a}(-x)}{e^uD_{-a-1}(x)+D_{-a-1}(-x)},
\end{equation}
where we have used \cite[Subsections~12.8 and~12.9]{NIST}.
Now, we establish the asymptotic expansion of $S_2$. 
\begin{lemma}\label{lemma:Analysis of S2}
There exists $\delta>0$ such that as $n\to+\infty$, we have  
\begin{align*}
S_2&=
C_2^{(2)}\sqrt{n}
+C_3^{(2)}
+C_n^{(2)}
+\mathcal{O}\Bigl(\frac{(\log n)^{\frac{1}{4}}}{n^{\frac{1}{4}}}\Bigr), 
\end{align*}
uniformly for $u\in\{z\in\C:|z-x|\leq \delta\}$ and $a$ in compact subsets of $(-1,+\infty)$, where 
\begin{align*}
C_2^{(2)}
&:=
\rho\sqrt{\Delta Q(\rho)}
\int_{-\infty}^{\infty}
\bigg(\log\bigl[\mathcal{H}_{a,u}(x)\bigr]
-a\log|x|-u\mathbf{1}_{(-\infty,0)}(x)
\bigg)
\,dx,
\\
C_3^{(2)}
&:=
-
\Bigl(\alpha+\frac{1}{2}\Bigr)u
+
\frac{1}{6}\Bigl(2+\frac{\rho\partial_r\Delta Q(\rho)}{\Delta Q(\rho)}\Bigr)u
-
\frac{a}{12}\Bigl(1-\frac{\rho\partial_r\Delta Q(\rho)}{\Delta Q(\rho)}\Bigr)u
\\
&\quad 
+\frac{1}{6}\Bigl(2+\frac{\rho\partial_r\Delta Q(\rho)}{\Delta Q(\rho)}\Bigr)
\int_{-\infty}^{+\infty}
\bigg[
x\Bigl(\log\bigl[\mathcal{H}_{a,u}(x)\bigr] 
- u\mathbf{1}_{(-\infty,0)}
\Big)
-ax\log|x| 
-\frac{a(a-1)x}{2(x^2+1)}
\bigg]\,dx,
\nonumber
\\
C_n^{(2)}
&:=
-\frac{a}{2}\Bigl(\rho\sqrt{\Delta Q(\rho)}\sqrt{n}\Delta_{n,+}
+\rho\sqrt{\Delta Q(\rho)}\sqrt{n}\Delta_{n,-}
+1\Bigr)\log(4n)
\\
&\quad
+
\sqrt{n}\rho\sqrt{\Delta Q(\rho)}
\bigg( 
a\,\Delta_{n,+}(\log \Delta_{n,+}-1)
+a\, \Delta_{n,-}(\log \Delta_{n,-}-1)
+u\Delta_{n,-}
\bigg)
+\frac{a}{2}\log\Delta_{n,+}
+\frac{a}{2}\log\Delta_{n,-}. 
\end{align*}
\end{lemma}

\begin{proof}
\begin{comment}
Note that
\begin{align*}
r_{\tau(g_{1,-})}
&=
\rho
-\frac{1}{2\rho\Delta Q(\rho)}m_{n,-}
-\frac{\Delta Q(\rho)+\rho\partial_r\Delta Q(\rho)}{8(\rho\Delta Q(\rho))^3}
m_{n,-}^2
+\mathcal{O}\bigl(m_{n,-}^3\bigr)
\\
&=\rho-\frac{1}{2\sqrt{\Delta Q(\rho)}}\frac{\Delta_{n,-}}{\sqrt{n}}
-\frac{\Delta Q(\rho)+\rho\partial_r\Delta Q(\rho)}{8\rho \Delta Q(\rho)^2}\frac{\Delta_{n,-}^2}{n}+\mathcal{O}\Bigl(\frac{\Delta_{n,-}^3}{n^{3/2}}\Bigr), 
\\
r_{\tau(g_{1,+})}
&=
\rho+\frac{1}{2\rho\Delta Q(\rho)}m_{n,+}
-\frac{\Delta Q(\rho)+\rho\partial_r\Delta Q(\rho)}{8(\rho\Delta Q(\rho))^3}m_{n,+}^2
+\mathcal{O}\bigl(m_{n,+}^3\bigr)
\\
&=
\rho+\frac{1}{2\sqrt{\Delta Q(\rho)}}\frac{\Delta_{n,+}}{\sqrt{n}}
-\frac{\Delta Q(\rho)+\rho\partial_r\Delta Q(\rho)}{8\rho \Delta Q(\rho)^2}\frac{\Delta_{n,+}^2}{n}+\mathcal{O}\Bigl(\frac{\Delta_{n,+}^3}{n^{3/2}}\Bigr). 
\end{align*}
\end{comment}
%Recall that $g_{1,+}=n(\tau_{\rho}+\delta_n)-\theta_{1,+}$ and $g_{1,-}=n(\tau_{\rho}-\delta_n)+\theta_{1,-}$, and then we have 
By \eqref{def of critical indices}, as $n\to+\infty$, we have 
\[
g_{1,+}-g_{1,-}+1
=
\rho\sqrt{\Delta Q(\rho)}\sqrt{n}\Delta_{n,+}
+\rho\sqrt{\Delta Q(\rho)}\sqrt{n}\Delta_{n,-}
+1.
\]
Therefore, the first and second terms of \eqref{def of S2 asymptotics part 1} satisfy 
\begin{align*}
&\quad
\sum_{j=g_{1,-}}^{g_{1,+}}\log\Bigl(\frac{1}{(4n)^{\frac{a}{2}}}\Bigr)
+
\frac{a}{2}
\sum_{j=g_{1,-}}^{g_{1,+}}\log\Bigl(\frac{1}{\Delta Q(\rho)}\Bigr)
\\
&=
-\Bigl( \rho\sqrt{\Delta Q(\rho)}\sqrt{n}\Delta_{n,+}
+\rho\sqrt{\Delta Q(\rho)}\sqrt{n}\Delta_{n,-}
+1\Bigr)
a\log2
\\
&\quad
-\Bigl( \rho\sqrt{\Delta Q(\rho)}\sqrt{n}\Delta_{n,+}
+\rho\sqrt{\Delta Q(\rho)}\sqrt{n}\Delta_{n,-}
+1\Bigr)\frac{a}{2}\log n
\\
&\quad
-\Bigl( \rho\sqrt{\Delta Q(\rho)}\sqrt{n}\Delta_{n,+}
+\rho\sqrt{\Delta Q(\rho)}\sqrt{n}\Delta_{n,-}
+1\Bigr)\frac{a}{2}\log \Delta Q(\rho).
\end{align*}
By Lemma~\ref{lemma:Euler-Maclaurin}, as $n\to+\infty$, we have 
\begin{align*}
\sum_{j=g_{1,-}}^{g_{1,+}}
\log \bigl[\mathcal{H}_{a,u}(\xi)\bigr] 
&=
\int_{g_{1,-}}^{g_{1,+}}\log\bigl[\mathcal{H}_{a,u}(\xi_t)\bigr]\,dt
+\frac{1}{2}\log\bigl[\mathcal{H}_{a,u}(-\Delta_{-,n})\bigr]
+\frac{1}{2}\log\bigl[\mathcal{H}_{a,u}(\Delta_{+,n})\bigr]
+\mathcal{O}\Bigl(\frac{(\log n)^{\frac{1}{4}}}{n^{\frac{1}{4}}}\Bigr). 
\end{align*}
By change of variables \eqref{def of rtau ODE}, we get 
\[
\int_{g_{1,-}}^{g_{1,+}}\log\bigl[\mathcal{H}_{a,u}(\xi_t)\bigr]\,dt
=
\sqrt{n}\,\rho\sqrt{\Delta Q(\rho)}
\int_{-\Delta_{n,-}}^{\Delta_{n,+}}
\log\bigl[\mathcal{H}_{a,u}(x)\bigr]\,dx. 
\] 
From \cite[Section 12.9]{NIST}, we have 
\begin{comment}
\begin{align}
\mathcal{H}_{a,u}(x)
&=
x^a\Bigl(1+\frac{a(a-1)}{2x^2}+\mathcal{O}(x^{-3})\Bigr)+\mathcal{O}(e^{-\frac{x^2}{2}}), \qquad x\to+\infty, 
\\
\mathcal{H}_{a,u}(x)
&=e^u(-x)^a\Bigl(1+\frac{a(a-1)}{2x^2}+\mathcal{O}((-x)^{-3})\Bigr)+\mathcal{O}(e^{-\frac{x^2}{2}}),\qquad x\to-\infty. 
\end{align}
\end{comment}
\begin{equation}
\label{def of asymptotics of mathcal H au x}
\log\bigl[\mathcal{H}_{a,u}(x)\bigr]
=
a\log |x|+u\mathbf{1}_{(-\infty,0)}(x)+\frac{a(a-1)}{2x^2}-\frac{a(a-1)(2a-3)}{4x^4}+\mathcal{O}(x^{-6}),\quad x\to\pm\infty.
\end{equation}
By the above and regularizing the integral, we have 
\begin{align*}
&\quad\sqrt{n}\,\rho\sqrt{\Delta Q(\rho)}
\int_{-\Delta_{n,-}}^{\Delta_{n,+}}
\log\bigl[\mathcal{H}_{a,u}(x)\bigr]\,dx
\\
&=
\sqrt{n}\,\rho\sqrt{\Delta Q(\rho)}
\int_{-\infty}^{+\infty}
\bigg(\log\bigl[\mathcal{H}_{a,u}(x)\bigr]
-a\log|x|-u\mathbf{1}_{(-\infty,0)}(x)
\bigg)
\,dx
\\
&\quad
+\sqrt{n}\rho\sqrt{\Delta Q(\rho)}
\bigg( 
a\,\Delta_{n,+}(\log \Delta_{n,+}-1)
+a\, \Delta_{n,-}(\log \Delta_{n,-}-1)
+u\Delta_{n,-}
\bigg)
\\
&\quad
-\sqrt{n}\rho\sqrt{\Delta Q(\rho)}\Bigl( 
\frac{a(a-1)}{2\Delta_{n,+}}
+\frac{a(a-1)}{2\Delta_{n,-}}
-\frac{a(a-1)(2a-3)}{12\Delta_{n,+}^3}
-\frac{a(a-1)(2a-3)}{12\Delta_{n,-}^3}
+\mathcal{O}(\Delta_{n,+}^{-5}+\Delta_{n,-}^{-5})
\Bigr).
\end{align*}
\begin{comment}
where we have used 
\begin{align*}
&\quad\sqrt{n}\,\rho\sqrt{\Delta Q(\rho)}
\int_{-\Delta_{n,-}}^{\Delta_{n,+}}
\Bigl(a\log|x|+u\mathbf{1}_{(-\infty,0)}(x)\Bigr)\,dx
\\
&=\sqrt{n}\rho\sqrt{\Delta Q(\rho)}
\bigg( 
a\,\Delta_{n,+}(\log \Delta_{n,+}-1)
+
a\,\Delta_{n,-}(\log \Delta_{n,-}-1)
+u\Delta_{n,-}
\bigg). 
\end{align*}
Note that 
\begin{align*}
&\quad
-\sqrt{n}\,\rho\sqrt{\Delta Q(\rho)}
\int_{\Delta_{n,+}}^{+\infty}\Bigl(\frac{a(a-1)}{2x^2}-\frac{a(a-1)(2a-3)}{4x^4}+\mathcal{O}(x^{-6})\Bigr)\,dx
\\
&\quad
-\sqrt{n}\,\rho\sqrt{\Delta Q(\rho)}
\int_{-\infty}^{-\Delta_{n,-}}\Bigl(\frac{a(a-1)}{2x^2}-\frac{a(a-1)(2a-3)}{4x^4}+\mathcal{O}(x^{-6})\Bigr)\,dx
\\
&=
-\sqrt{n}\rho\sqrt{\Delta Q(\rho)}\Bigl( 
\frac{a(a-1)}{2\Delta_{n,+}}
+\frac{a(a-1)}{2\Delta_{n,-}}
-\frac{a(a-1)(2a-3)}{12\Delta_{n,+}^3}
-\frac{a(a-1)(2a-3)}{12\Delta_{n,-}^3}
+\mathcal{O}(\Delta_{n,+}^{-5}+\Delta_{n,-}^{-5})
\Bigr). 
\end{align*}
\end{comment}
\begin{comment}
Furthermore, observe that as $n\to+\infty$, 
\begin{align*}
\frac{1}{2}\log\Bigl[\frac{\Gamma(a+1)}{\sqrt{2\pi}}e^{-\frac{\xi(g_{1,+})^2}{4}}g_{0}(\xi(g_{1,+}))\Bigr]    
&=
\frac{a}{2}\log\Delta_{n,+}+\mathcal{O}\Bigl(\frac{1}{\Delta_{n,+}^2}\Bigr),
\\
\frac{1}{2}\log\Bigl[\frac{\Gamma(a+1)}{\sqrt{2\pi}}e^{-\frac{\xi(g_{1,-})^2}{4}}g_{0}(\xi(g_{1,-}))\Bigr]
&=\frac{u}{2}+\frac{a}{2}\log\Delta_{n,-}+\mathcal{O}\Bigl(\frac{1}{\Delta_{n,-}^2}\Bigr).
\end{align*}
\end{comment}
Next we observe that 
\begin{align*}
\frac{1}{\sqrt{n}} \sum_{j=g_{1,-}}^{g_{1,+}}
\frac{\widetilde{g}_{a,u}(\xi)}{12\rho\sqrt{\Delta Q(\rho)}g_{a,u}(\xi)}
&=
\frac{1}{12}\int_{-\Delta_{n,-}}^{\Delta_{n,+}}
\frac{\widetilde{g}_{a,u}(x)}{g_{a,u}(x)}\,dx
+\mathcal{O}\Bigl(\frac{(\log n)^{\frac{1}{4}}}{n^{\frac{1}{4}}}\Bigr),
%+
%\frac{1}{\sqrt{n}} 
%\frac{g_1(\xi(g_{1,+}))}{12\rho\sqrt{\Delta Q(\rho)}g_0(\xi(g_{1,+}))}
%\\
%&\quad
%+
%\frac{1}{\sqrt{n}} 
%\frac{g_1(\xi(g_{1,-}))}{12\rho\sqrt{\Delta Q(\rho)}g_0(\xi(g_{1,-}))},
\end{align*}
Note that by \eqref{def of tilde gau}, \eqref{def of cal H au}, and \eqref{def of d relationship between g0 g1}, we have 
\begin{align*}
\frac{1}{12}\int_{-\Delta_{n,-}}^{\Delta_{n,+}}
\frac{\widetilde{g}_{a,u}(x)}{g_{a,u}(x)}\,dx
&=
\frac{1}{6}\Bigl(2+\frac{\rho\,\partial_r\Delta Q(\rho)}{\Delta Q(\rho)}\Bigr)
\int_{-\Delta_{n,-}}^{\Delta_{n,+}}
x\log\bigl[\mathcal{H}_{a,u}(x)\bigr]\,dx
-\frac{a}{24}\Bigl(1-\frac{\rho\,\partial_r\Delta Q(\rho)}{\Delta Q(\rho)}\Bigr)(\Delta_{n,+}^2-\Delta_{n,-}^2)
\\
&
\quad
-\frac{1}{12}\Bigl(2+\frac{\rho\,\partial_r\Delta Q(\rho)}{\Delta Q(\rho)}\Bigr)\Bigl[\Delta_{n,+}^2
\log\bigl[\mathcal{H}_{a,u}(\Delta_{n,+})\bigr]
-\Delta_{n,-}^2
\log\bigl[\mathcal{H}_{a,u}(-\Delta_{n,-})\bigr]\Bigr]
\\
&
\quad
+\Bigl[
\alpha+\frac{1}{2}+\frac{2+a}{12}\Bigl(1-\frac{\rho\partial_r\Delta Q(\rho)}{\Delta Q(\rho)}\Bigr)\Bigr]
\Bigl[\log\bigl[\mathcal{H}_{a,u}(\Delta_{n,+})\bigr]-\log\bigl[\mathcal{H}_{a,u}(-\Delta_{n,-})\bigr]\Bigr]. 
\end{align*}
\begin{comment}
Note that 
\begin{align*}
\frac{1}{12}\frac{g_1(x)}{g_0(x)}
&=
-\frac{a}{12}\Bigl(1+\frac{2\rho\partial_r\Delta Q(\rho)}{\Delta Q(\rho)}\Bigr)
x
-
\frac{1}{12}\Bigl(2+\frac{\rho\partial_r\Delta Q(\rho)}{\Delta Q(\rho)}\Bigr)
x^2
\partial_x\log\Bigl[\frac{\Gamma(a+1)}{\sqrt{2\pi}}e^{-\frac{x^2}{4}}g_0(x)\Bigr]
\\
&\quad 
+\Bigl(\alpha+\frac{1}{2}\Bigr)
\partial_x\log\Bigl[\frac{\Gamma(a+1)}{\sqrt{2\pi}}e^{-\frac{x^2}{4}}g_0(x)\Bigr]
+
\frac{2+a}{12}\Bigl(1-\frac{\rho\partial_r\Delta Q(\rho)}{\Delta Q(\rho)}\Bigr)
\partial_x\log\Bigl[\frac{\Gamma(a+1)}{\sqrt{2\pi}}e^{-\frac{x^2}{4}}g_0(x)\Bigr].   
\end{align*}
\end{comment}
By regularizing the integral, we obtain 
\begin{align*}
&\quad \frac{1}{12}\int_{-\Delta_{n,-}}^{\Delta_{n,+}}
\frac{\widetilde{g}_{a,u}(x)}{g_{a,u}(x)}\,dx
\\
&=
-\frac{a}{24}\Bigl(1+\frac{2\rho\partial_r\Delta Q(\rho)}{\Delta Q(\rho)}\Bigr)(\Delta_{n,+}^2-\Delta_{n,-}^2)
+\Bigl(\alpha+\frac{1}{2}\Bigr)
\log\bigl[\mathcal{H}_{a,u}(\Delta_{n,+})\bigr]
-
\Bigl(\alpha+\frac{1}{2}\Bigr)
\log\bigl[\mathcal{H}_{a,u}(-\Delta_{n,+})\bigr]
\\
&\quad
+
\frac{2+a}{12}\Bigl(1-\frac{\rho\partial_r\Delta Q(\rho)}{\Delta Q(\rho)}\Bigr)
\log\bigl[\mathcal{H}_{a,u}(\Delta_{n,+})\bigr]
-
\frac{2+a}{12}\Bigl(1-\frac{\rho\partial_r\Delta Q(\rho)}{\Delta Q(\rho)}\Bigr)
\log\bigl[\mathcal{H}_{a,u}(-\Delta_{n,-})\bigr]
\\
&\quad
-
\frac{1}{12}\Bigl(2+\frac{\rho\partial_r\Delta Q(\rho)}{\Delta Q(\rho)}\Bigr)
\Delta_{n,+}^2\log\bigl[\mathcal{H}_{a,u}(\Delta_{n,+})\bigr]
+
\frac{1}{12}\Bigl(2+\frac{\rho\partial_r\Delta Q(\rho)}{\Delta Q(\rho)}\Bigr)
\Delta_{n,-}^2\log\bigl[\mathcal{H}_{a,u}(-\Delta_{n,-})\bigr]
\\
&\quad 
+\frac{1}{6}\Bigl(2+\frac{\rho\partial_r\Delta Q(\rho)}{\Delta Q(\rho)}\Bigr)
\int_{-\infty}^{+\infty}
\bigg[
x\Bigl(\log\bigl[\mathcal{H}_{a,u}(x)\bigr] 
- u\mathbf{1}_{(-\infty,0)}
\Big)
-ax\log|x| 
-\frac{a(a-1)x}{2(x^2+1)}
\bigg]\,dx
\\
&\quad 
-\frac{u}{12}\Bigl(2+\frac{\rho\partial_r\Delta Q(\rho)}{\Delta Q(\rho)}\Bigr)
\Delta_{n,-}^2
+\frac{1}{6}\Bigl(2+\frac{\rho\partial_r\Delta Q(\rho)}{\Delta Q(\rho)}\Bigr)
\int_{-\Delta_{n,-}}^{\Delta_{n,+}}
\Bigl(ax\log|x| + \frac{a(a-1)x}{2(x^2+1)}\Bigr)\,dx
+\mathcal{O}\Bigl(\frac{(\log n)^{\frac{1}{4}}}{n^{\frac{1}{4}}}\Bigr).
\end{align*}
\begin{comment}
Since 
\begin{align*}
\log\Bigl[\frac{\Gamma(a+1)}{\sqrt{2\pi}}e^{-\frac{\Delta_{n,+}^2}{4}}g_0(\Delta_{n,+})\Bigr]
&=a\log \Delta_{n,+}+\frac{a(a-1)}{2\Delta_{n,+}^2}+\mathcal{O}(\Delta_{n,+}^{-4}),
\\
\log\Bigl[\frac{\Gamma(a+1)}{\sqrt{2\pi}}e^{-\frac{\Delta_{n,-}^2}{4}}g_0(\Delta_{n,-})\Bigr]
&=u+a\log \Delta_{n,-}+\frac{a(a-1)}{2\Delta_{n,-}^2}+\mathcal{O}(\Delta_{n,-}^{-4}), 
\end{align*}
By \eqref{def of asymptotics of mathcal H au x}, we have 
\begin{align*}
&\quad \frac{1}{12}\int_{-\Delta_{n,-}}^{\Delta_{n,+}}
\frac{\widetilde{g}_{a,u}(x)}{g_{a,u}(x)}\,dx
\\
&=
\frac{1}{6}\Bigl(2+\frac{\rho\partial_r\Delta Q(\rho)}{\Delta Q(\rho)}\Bigr)
\int_{-\infty}^{+\infty}
\bigg[
x\Bigl(\log\bigl[\mathcal{H}_{a,u}(x)\bigr] 
- u\mathbf{1}_{(-\infty,0)}
\Big)
-ax\log|x| 
-\frac{a(a-1)x}{2(x^2+1)}
\bigg]\,dx
\\
&\quad
-
\Bigl(\alpha+\frac{1}{2}\Bigr)u
-
\frac{2+a}{12}\Bigl(1-\frac{\rho\partial_r\Delta Q(\rho)}{\Delta Q(\rho)}\Bigr)u 
+\mathcal{O}\Bigl(\frac{(\log n)^{\frac{1}{4}}}{n^{\frac{1}{4}}}\Bigr). 
\end{align*}
\end{comment}
Combining all of the above together with \eqref{def of asymptotics of mathcal H au x}, we obtain the result. 
\end{proof}

\subsection{Proof of Theorem~\ref{theorem:calEn expansion}}
We now complete the proof of Theorem~\ref{theorem:calEn expansion}. 
Combining the error terms from Lemma~\ref{lemma:S1},~\ref{lemma:Analysis of S2}, and~\ref{lemma:S3}, we have $C_n^{(1)}+C_{n}^{(2)}+C_{n}^{(3)}=
-a\log2-\frac{a}{2}\log\Delta Q(\rho)$, which is added to the term of order $\mathcal{O}(1)$.
Thus, by Lemma~\ref{lemma:S0},~\ref{lemma:S1},~\ref{lemma:Analysis of S2}, and~\ref{lemma:S3}, we get $C_3^{(1)}+C_3^{(2)}+C_3^{(3)}-a\log2-\frac{a}{2}\log\Delta Q(\rho)=C_3(u,a)$.
By Lemma~\ref{lemma:S0},~\ref{lemma:S1},~\ref{lemma:Analysis of S2}, and~\ref{lemma:S3}, $C_1(u,a),C_2(u,a)$ are similarly computed.
This completes the proof of Theorem~\ref{theorem:calEn expansion}. 
\begin{comment}
    
\begin{align*}
C_3^{(1)}&:=
\end{align*}

\begin{align*}
C_3^{(3)}&:=
-\frac{a}{2}\log\Bigl(\frac{r_1}{\rho}-1\Bigr)
-\frac{a(a-1)}{4}
\frac{r_1}{r_1-\rho}
-
\frac{a}{4}
\int_{0}^{r_1}
\Bigl(
\frac{1}{u-\rho}
\frac{u\partial_u\Delta Q(u)}{\Delta Q(u)}
-\frac{1}{u-\rho}
\frac{\rho\partial_r\Delta Q(\rho)}{\Delta Q(\rho)}
\Bigr)
\,du
\\
&\quad
+
\frac{a}{4}\Bigl(
4\alpha+a+2-\frac{\rho\partial_r\Delta Q(\rho)}{\Delta Q(\rho)}
\Bigr)\log\Bigl(\frac{r_1}{\rho}-1\Bigr)
\end{align*}


\begin{align*}
    C_3^{(2)}
&:=
-
\Bigl(\alpha+\frac{1}{2}\Bigr)u
+
\frac{1}{6}\Bigl(2+\frac{\rho\partial_r\Delta Q(\rho)}{\Delta Q(\rho)}\Bigr)u
-
\frac{a}{12}\Bigl(1-\frac{\rho\partial_r\Delta Q(\rho)}{\Delta Q(\rho)}\Bigr)u
\\
&\quad 
+\frac{1}{6}\Bigl(2+\frac{\rho\partial_r\Delta Q(\rho)}{\Delta Q(\rho)}\Bigr)
\int_{-\Delta_{n,-}}^{\Delta_{n,+}}
\bigg[
x\Bigl(\log\Bigl[\frac{\Gamma(a+1)}{\sqrt{2\pi}}e^{-\frac{x^2}{4}}g_0(x)\Bigr] 
- u\mathbf{1}_{(-\infty,0)}
\Big)
-ax\log|x| 
-\frac{a(a-1)x}{2(x^2+1)}
\bigg]\,dx
\end{align*}

\begin{align*}
&\quad C_3^{(1)}+C_3^{(2)}+C_3^{(3)}+C_n^{(3)}
\\
&=-\frac{a}{2}\log\Bigl(\frac{r_1}{\rho}-1\Bigr)
-\frac{a(a-1)}{4}
\frac{r_1}{r_1-\rho}
-
\frac{a}{4}
\int_{0}^{r_1}
\Bigl(
\frac{1}{u-\rho}
\frac{u\partial_u\Delta Q(u)}{\Delta Q(u)}
-\frac{1}{u-\rho}
\frac{\rho\partial_r\Delta Q(\rho)}{\Delta Q(\rho)}
\Bigr)
\,du
\\
&\quad
+
\frac{a}{4}\Bigl(
4\alpha+a+2-\frac{\rho\partial_r\Delta Q(\rho)}{\Delta Q(\rho)}
\Bigr)\log\Bigl(\frac{r_1}{\rho}-1\Bigr)
\\
&\quad
-
\Bigl(\alpha+\frac{1}{2}\Bigr)u
+
\frac{1}{6}\Bigl(2+\frac{\rho\partial_r\Delta Q(\rho)}{\Delta Q(\rho)}\Bigr)u
-
\frac{a}{12}\Bigl(1-\frac{\rho\partial_r\Delta Q(\rho)}{\Delta Q(\rho)}\Bigr)u
\\
&\quad 
+\frac{1}{6}\Bigl(2+\frac{\rho\partial_r\Delta Q(\rho)}{\Delta Q(\rho)}\Bigr)
\int_{-\Delta_{n,-}}^{\Delta_{n,+}}
\bigg[
x\Bigl(\log\Bigl[\frac{\Gamma(a+1)}{\sqrt{2\pi}}e^{-\frac{x^2}{4}}g_0(x)\Bigr] 
- u\mathbf{1}_{(-\infty,0)}
\Big)
-ax\log|x| 
-\frac{a(a-1)x}{2(x^2+1)}
\bigg]\,dx. 
\end{align*}
\end{comment}

\subsubsection*{Acknowledgements}
%The author is grateful to Christophe Charlier for their valuable feedback to improve the paper and for engaging in numerous stimulating discussions.
The author is grateful to Sung-Soo Byun and Seong-Mi Seo for valuable feedbacks.
The author acknowledges support from the European Research Council (ERC), Grant Agreement No. 101115687.% by his promoter Christophe Charlier who is a Research Associate of the Fonds de la Recherche Scientifique FNRS.



\appendix
\section{Consistency between Theorem~\ref{theorem:asymptotic expansion of counting statistics} and \cite[Theorem~1.1]{C2021 FH}}
\label{section: appendix counting}
In this appendix, we confirm the consistency between Theorem~\ref{theorem:asymptotic expansion of counting statistics} and \cite[Theorem~1.1]{C2021 FH} mentioned in Remark~\ref{remark:counting statistics}. 
We recall that $C_1,C_2,C_3$ for $Q(z)=|z|^{2b}$ ($b>0$) $r_1=\rho$, and $r_{k}\equiv0$ for $k=2,3,\dots,m+1$ in \cite[Theorem 1.1]{C2021 FH} are given by 
\begin{align}
\label{def of C1 count MF}
C_1&=b\rho^{2b}u,
\\
\label{def of C2 count MF}
C_2&=\sqrt{2}b\rho^b\int_0^{+\infty}\Bigl(
\mathcal{F}(t,e^u)+\mathcal{F}(t,e^{-t})
\Bigr)\,dt, 
\\
\label{def of C3 count MF}
C_3&=-\Bigl(\frac{1}{2}+\alpha\Bigr)u+4b\int_0^{+\infty}t\bigl(\mathcal{F}(t,e^u)-\mathcal{F}(t,e^{-u})\bigr)\,dt
+b\int_{-\infty}^{+\infty}\mathcal{G}(t,e^u)\frac{5t^2-1}{3}\,dt, 
\end{align}
where $\mathcal{F}(t,s)$ for $t\in\R$ and $s\in\C\backslash(-\infty,0]$ is given by \eqref{def of mathcalFts} and 
\begin{equation*}
\mathcal{G}(t,s):= \frac{1-s}{1+\frac{s-1}{2}\mathrm{erfc}(t)}\frac{e^{-t^2}}{\sqrt{\pi}} =\frac{d}{dt}\mathcal{F}(t,s). 
\end{equation*}
It is straightforward to see that $C_{1}(u)$ and $C_2(u)$ in Theorem~\ref{theorem:asymptotic expansion of counting statistics} are consistent with \eqref{def of C1 count MF} and \eqref{def of C2 count MF}, respectively.
To see that $C_3$ in \cite[Theorem 1.1]{C2021 FH} is consistent with \eqref{def of C3 count MF}, by integration by parts, the last term in \eqref{def of C3 count MF} can be rewritten as
\begin{align*}
\int_{-\infty}^{+\infty}\mathcal{G}(t,e^u)\frac{5t^2-1}{3}\,dt
&=
b\lim_{M\to+\infty}\int_0^{M}\mathcal{G}(t,e^u)\frac{5t^2-1}{3}\,dt
-
b\lim_{M\to+\infty}\int_0^{M}\mathcal{G}(t,e^{-u})\frac{5t^2-1}{3}\,dt
\\
&=\frac{b}{3}u-\frac{10b}{3}\int_0^{+\infty}t\bigl( \mathcal{F}(t,e^{u})-\mathcal{F}(t,e^{-u})\bigr)\,dt. 
\end{align*}
Therefore, \eqref{def of C3 count MF} can be deduced to 
\[
C_3=-\Bigl(\frac{1}{2}+\alpha\Bigr)u+\frac{b}{3}u+\frac{2b}{3}\int_0^{+\infty}t\Bigl(\mathcal{F}(t,e^{u})-\mathcal{F}(t,e^{-u})\Bigr)\,dt. 
\]
If we informally Theorem~\ref{theorem:asymptotic expansion of counting statistics} to the case $Q(z)=|z|^{2b}$ ($b>0$), by $2+\frac{\rho\partial_r\Delta Q(\rho)}{\Delta Q(\rho)}=2b$, we obtain $C_3(u)=C_3$.
This shows the consistency between Theorem~\ref{theorem:asymptotic expansion of counting statistics} and \cite[Theorem~1.1]{C2021 FH}.

\section{Consistency between Theorem~\ref{theorem:calEn expansion} and 
\cite[Theorem~1.1]{BC2022}}
\label{section:appendix parabolic cylinder function}
We confirm the consistency between Theorem~\ref{theorem:calEn expansion} and 
\cite[Theorem~1.1]{BC2022}. 
We begin with collecting some facts on the parabolic cylinder function and the associated Hermite polynomials from \cite{NIST,Temme,Wunsche,BC2022}.  
The $\nu$-th associated Hermite polynomials $\{\mathrm{He}_k^{(\nu)};k=0,1,\dots\}$ are defined recursively by 
\begin{equation}
\begin{cases}
    \mathrm{He}_{k+1}^{(\nu)}(x)=x\mathrm{He}_{k}^{(\nu)}(x)-(k+\nu)\mathrm{He}_{k-1}^{(\nu)}(x), & k\geq 1, \\
    \mathrm{He}_{0}^{(\nu)}(x)=1,\qquad \mathrm{He}_{1}^{(\nu)}(x)=x,
\end{cases}    
\end{equation}
and satisfy the orthogonality relations 
\begin{equation}
\int_{-\infty}^{+\infty}\mathrm{He}_{k}^{(\nu)}(x)\mathrm{He}_{\ell}^{(\nu)}(x)\,\frac{dx}{|D_{-\nu}(ix)|^2}=\sqrt{2\pi}(k+\nu)!\delta_{k,\ell}, 
\end{equation}
see \cite[Eq. (12.7.2)]{NIST}, where the parabolic cylinder function $D_{-\nu}(x)$ is given by \eqref{def of parabolic cylinder function}.  
\begin{comment}
For instance, for $\nu\in\R$, 
\begin{align}
\begin{split}
D_{\nu}(z)&:=2^{\frac{\nu}{2}}e^{-\frac{z^2}{4}}\Bigl[
\frac{(-\tfrac{1}{2})!}{(-\tfrac{1+\nu}{2})!} {}_1F_1\Bigl(-\frac{\nu}{2};\frac{1}{2};\frac{z^2}{2}\Bigr)
+\frac{z}{\sqrt{2}}\frac{(-\tfrac{3}{2})!}{(-1-\tfrac{\nu}{2})!}{}_1F_1\Bigl(\frac{1-\nu}{2};\frac{3}{2};\frac{z^2}{2}\Bigr)
\Bigr]    
\\
&\,=
2^{\frac{\nu}{2}}e^{\frac{z^2}{4}}\Bigl[
\frac{(-\tfrac{1}{2})!}{(-\tfrac{1+\nu}{2})!} {}_1F_1\Bigl(\frac{1+\nu}{2};\frac{1}{2};-\frac{z^2}{2}\Bigr)
+\frac{z}{\sqrt{2}}\frac{(-\tfrac{3}{2})!}{(-1-\tfrac{\nu}{2})!}{}_1F_1\Bigl(1+\frac{\nu}{2};\frac{3}{2};-\frac{z^2}{2}\Bigr)
\Bigr].     
\end{split}
\end{align}
\end{comment}
It is known that the parabolic cylinder function is related to a family of Hermite polynomials $\{H_n\}_{n\in\N}$, i.e., 
\begin{equation}
    D_{n}(z)=e^{-\frac{1}{4}z^2}2^{-\frac{n}{2}}H_n\Bigl(\frac{z}{\sqrt{2}}\Bigr)
    =e^{-\frac{1}{4}z^2}\mathrm{He}_n(z), \qquad n=0,1,2,\dots,
\end{equation}
where $H_n(x)=2^{\frac{n}{2}}\mathrm{He}_n(\sqrt{2}x)$ is defined by 
\[
H_n(x)=(-1)^ne^{x^2}\frac{d^n}{dx^n}e^{-x^2},
\]
see \cite[Eq. (18.5.5)]{NIST}.
For $\nu\in\R$ with $\nu\notin\mathbb{Z}_{<0}$ and $n\in\mathbb{Z}_{\geq 0}$, 
\begin{equation}
D_{-\nu-n}(z):=\frac{\Gamma(-\nu-n+1)}{\Gamma(-\nu+1)}
\Bigl[
(-i)^n\mathrm{He}_{n}^{(\nu-1)}(iz)D_{-\nu}(z)-(-i)^{n-1}\mathrm{He}_{n-1}^{(\nu)}(iz)D_{-\nu+1}(z)
\Bigr].
\end{equation}
Particularly, 
\begin{equation}
\label{def of D0 Dminus1}
D_0(z):=U(-\tfrac{1}{2},z)=e^{-\frac{z^2}{4}},\qquad
D_{-1}(z):=U(\tfrac{1}{2},z)=e^{\frac{z^2}{4}}\sqrt{\frac{\pi}{2}}\mathrm{erfc}\Bigl(\frac{z}{\sqrt{2}}\Bigr). 
\end{equation}
If $a\in\mathbb{Z}_{>0}$ and $\nu=1$, then by $\lim_{\nu\to1}\frac{\Gamma(-\nu-a+1)}{\Gamma(-\nu+1)}=\frac{(-1)^a}{a!}$, one can write 
\[
e^{-\frac{z^2}{4}}D_{-1-a}(z)
=\frac{(-1)^a}{a!}
\Bigl[
\frac{1}{i^{a}}\mathrm{He}_{a}^{(0)}(iz)\sqrt{\frac{\pi}{2}}\mathrm{erfc}\Bigl(\frac{z}{\sqrt{2}}\Bigr)
-\frac{1}{i^{a-1}}\mathrm{He}_{a-1}^{(1)}(iz)e^{-\frac{z^2}{2}}
\Bigr].
\]
Next, we recall some functionals from \cite[Eqs. (1.8), (1.9), (1.10), (1.11), (1.13), and (1.14)]{BC2022} (here we mainly focus on $a\geq 3$ for the simplicity), 
\begin{align*}
    p_{0,a}(x)&:=\frac{1}{i^a} \mathrm{He}_{a}(ix), \qquad  q_{0,a}(x):=\frac{1}{i^{a-1}}\mathrm{He}_{a-1}^{(1)}(ix), \\
    p_{1,a}(x)&:=-\frac{a}{2}p_{0,a+1}(x)-ab \Bigl(p_{0,a+1}(x)-(3a-1)p_{0,a-1}(x)+\frac{5}{3}(a-1)(a-2)p_{0,a-3}(x)\Bigr),
    \\
    q_{1,a}(x)&:=-\frac{a}{2}q_{0,a+1}(x)-ab \Bigl(q_{0,a+1}(x)-(3a-1)q_{0,a-1}(x)+\frac{5}{3}(a-1)(a-2)q_{0,a-3}(x)\Bigr), 
    \\
    \mathcal{G}_{0}(y;u,a)&:=p_{0,a}(-\sqrt{2}y)\Bigl((-1)^a+\frac{e^u-(-1)^a}{2}\mathrm{erfc}(y)\Bigr)+q_{0,a}(-\sqrt{2}y)(e^u-(-1)^a)\frac{e^{-y^2}}{\sqrt{2\pi}}, 
    \\
    \mathcal{G}_{1}(y;u,a)&:=p_{1,a}(-\sqrt{2}y)\Bigl((-1)^{a}+\frac{e^u-(-1)^a}{2}\mathrm{erfc}(y)\Bigr)
    +q_{1,a}(-\sqrt{2}y)(e^{u}-(-1)^{a})\frac{e^{-y^2}}{\sqrt{2\pi}}. 
\end{align*}
The fact that $\mathrm{He}_{a}^{(\nu)}(-iy)=(-1)^a\mathrm{He}_{a}^{(\nu)}(iy)$ for $y\in\R$ and $a\in\N$ gives rise to 
\begin{align*}
&\quad
e^{u}\frac{\Gamma(a+1)}{\sqrt{2\pi}}e^{-\frac{y^2}{4}}D_{-1-a}(y)
+\frac{\Gamma(a+1)}{\sqrt{2\pi}}e^{-\frac{y^2}{4}}D_{-1-a}(-y)
\\
&=
\frac{1}{i^{a}}\mathrm{He}_{a}^{(0)}(-iy)
\Bigl[
(-1)^a+(e^{u}-(-1)^a)\frac{1}{2}\mathrm{erfc}\Bigl(\frac{y}{\sqrt{2}}\Bigr)
\Bigr]
+
\bigl(e^{u}-(-1)^{a}\bigr)\frac{1}{i^{a-1}}\mathrm{He}_{a-1}^{(1)}(-iy)\frac{e^{-\frac{y^2}{2}}}{\sqrt{2\pi}}.
\end{align*}
This shows the consistency between $\mathcal{H}_{a,u}(y)$ given by \eqref{def of cal H au} and \cite[Eq. (1.13)]{BC2022} , i.e., 
\[
\mathcal{G}_{0}\Bigl(\frac{y}{\sqrt{2}};u,a\Bigr)
=
 e^{u}\frac{\Gamma(a+1)}{\sqrt{2\pi}}e^{-\frac{y^2}{4}}D_{-1-a}(y)
+\frac{\Gamma(a+1)}{\sqrt{2\pi}}e^{-\frac{y^2}{4}}D_{-1-a}(-y)
=
\mathcal{H}_{a,u}(y).
\]
By a similar manner to the above, since by \cite[Eq. (12.8.1)]{NIST}, 
\[
D_{2-a}(y)=(y^2+a-1)D_{-a}(y)+ayD_{-a-1}(y), \qquad
(a+1)D_{-a-2}(y)=-yD_{-a-1}(y)+D_{-a}(y), 
\]
we have 
\begin{align*}
\mathcal{G}_{1}\Bigl(\frac{y}{\sqrt{2}};u,a\Bigr)
&=
-\frac{\Gamma(a+1)}{\sqrt{2\pi}}e^{-\frac{y^2}{4}}
\bigg[
\Bigl(\frac{a}{2}+ab \Bigr)(a+1)\bigl(e^{u}D_{-2-a}(y)-D_{-2-a}(-y)\bigr)
\\
&\quad
-b(3a-1)\bigl(e^uD_{-a}(y)-D_{-a}(-y)\bigr)
+\frac{5}{3}b\bigl(e^u D_{2-a}(y)-D_{2-a}(-y)\bigr)
\bigg]
\\
&=
-\frac{\Gamma(a+1)}{\sqrt{2\pi}}e^{-\frac{y^2}{4}}
\bigg[
\Bigl(\frac{a}{2}+ab-b(3a-1)+\frac{5}{3}b(y^2+a-1)\Bigr)
(e^{u}D_{-a}(y)-D_{-a}(-y))
\\
&\qquad
+\Bigl(\frac{2}{3}ab-\frac{a}{2}\Bigr)(e^{u}D_{-a-1}(y)+D_{-a-1}(-y))
\bigg]. 
\end{align*}
This gives rise to 
\begin{equation}
\label{def of wG1 over wG1}
\frac{\mathcal{G}_{1}(\tfrac{y}{\sqrt{2}};u,a)}{\mathcal{G}_{0}(\tfrac{y}{\sqrt{2}};u,a)}
=
-\Bigl[
\frac{a}{2}+ab-b(3a-1)+\frac{5}{3}b(y^2+a-1)
\Bigr]\partial_y\bigl(\log\bigl[\mathcal{H}_{a,u}(y)\bigr]-u\mathbf{1}_{(-\infty,0)}(y)\bigr)
-\Bigl(-\frac{a}{2}+\frac{2}{3}ab\Bigr)y,
\end{equation}
where we have used \eqref{def of d relationship between g0 g1}. 
By substituting \eqref{def of wG1 over wG1} into \cite[$\tfrac{\mathcal{G}_1(y;u,a)}{\mathcal{G}_0(y;u,a)}$ of $C_3$ in Theorem 1.1]{BC2022} and straightforward computations, we find that \cite[$C_3$ in Theorem 1.1]{BC2022} recovers $C_3$ in Theorem~\ref{theorem:calEn expansion}. 
\begin{comment}

We observe that the constant $C_3$ in \cite[Theorem 1.1]{BC2022} can be furthermore simplified. Namely, we will simply 
\[
\int_{-\infty}^{+\infty}\frac{1}{\sqrt{2}}\frac{\mathcal{G}_{1}(y:u,a)}{\mathcal{G}_{0}(y:u,a)}\,dy
=
\int_{-\infty}^{+\infty}\frac{1}{2}\frac{\mathcal{G}_{1}(-\tfrac{y}{\sqrt{2}}:u,a)}{\mathcal{G}_{0}(-\tfrac{y}{\sqrt{2}}:u,a)}\,dy
=
\frac{1}{2}
\lim_{M\to+\infty}
\int_{-M}^{M}\frac{\mathcal{G}_{1}(-\tfrac{y}{\sqrt{2}}:u,a)}{\mathcal{G}_{0}(-\tfrac{y}{\sqrt{2}}:u,a)}\,dy. 
\]


We remark that
\[
 \widetilde{\mathcal{G}}_0(-y)
 =
 e^{u}\frac{\Gamma(a+1)}{\sqrt{2\pi}}e^{-\frac{y^2}{4}}D_{-1-a}(y)
+\frac{\Gamma(a+1)}{\sqrt{2\pi}}e^{-\frac{y^2}{4}}D_{-1-a}(-y)
\]
Note that
\begin{align*}
&\quad
p_{1,a}(y)\Bigl((-1)^{a}-\frac{(-1)^a}{2}\mathrm{erfc}(-\tfrac{y}{\sqrt{2}})\Bigr)
-q_{1,a}(y)(-1)^{a}\frac{e^{-\frac{y^2}{2}}}{\sqrt{2\pi}}
\\
&=
\Bigl(\frac{a}{2}+ab \Bigr)(-1)^{a+1}
\Bigl[
p_{0,a+1}(y)
\frac{1}{2}\mathrm{erfc}(\tfrac{y}{\sqrt{2}})
-
q_{0,a+1}(y)\frac{e^{-\frac{y^2}{2}}}{\sqrt{2\pi}}
\Bigr]
\\
&\quad
-ab(3a-1)
(-1)^{a-1}
\Bigl[
p_{0,a-1}(y)
\frac{1}{2}\mathrm{erfc}(\tfrac{y}{\sqrt{2}})
-q_{0,a-1}(y)\frac{e^{-\frac{y^2}{2}}}{\sqrt{2\pi}}
\Bigr]
\\
&\quad
+\frac{5}{3}ba(a-1)(a-2)(-1)^{a-3}
\Bigl[p_{0,a-3}(y)
\frac{1}{2}\mathrm{erfc}(\tfrac{y}{\sqrt{2}})
-q_{0,a-3}(y)\frac{e^{-\frac{y^2}{2}}}{\sqrt{2\pi}}
\Bigr]
\\
&=
\Bigl(\frac{a}{2}+ab \Bigr)(a+1)
\frac{\Gamma(a+1)}{\sqrt{2\pi}}e^{-\frac{y^2}{4}}D_{-2-a}(y)
-b(3a-1)\frac{\Gamma(a+1)}{\sqrt{2\pi}}e^{-\frac{y^2}{4}}D_{-a}(y)
+\frac{5}{3}b
\frac{\Gamma(a+1)}{\sqrt{2\pi}}e^{-\frac{y^2}{4}}D_{2-a}(y)
\end{align*}
where we have used that $1-\frac{1}{2}\mathrm{erfc}(-\tfrac{y}{\sqrt{2}})=\frac{1}{2}\mathrm{erfc}(\tfrac{y}{\sqrt{2}})$. 
Note also that 
\begin{align*}
&\quad
e^{u}\Bigl(p_{1,a}(y)\frac{1}{2}\mathrm{erfc}(-\tfrac{y}{\sqrt{2}})
+q_{1,a}(y)\frac{e^{-\frac{y^2}{2}}}{\sqrt{2\pi}}
\Bigr)
\\
&=
-e^{u}
(-1)^{a+1}
\Bigl(\frac{a}{2}+ab \Bigr)
\Bigl[p_{0,a+1}(-y)\frac{1}{2}\mathrm{erfc}(-\tfrac{y}{\sqrt{2}})
-q_{0,a+1}(-y)\frac{e^{-\frac{y^2}{2}}}{\sqrt{2\pi}}
\Bigr]
\\
&\quad
+e^{u}ab (3a-1)(-1)^{a-1}
\Bigl[p_{0,a-1}(-y)
\frac{1}{2}\mathrm{erfc}(-\tfrac{y}{\sqrt{2}})
-q_{0,a-1}(-y)\frac{e^{-\frac{y^2}{2}}}{\sqrt{2\pi}}
\Bigr]
\\
&\quad
-e^{u}\frac{5}{3}ba(a-1)(a-2)(-1)^{a-3}
\Bigl[p_{0,a-3}(-y)
\frac{1}{2}\mathrm{erfc}(-\tfrac{y}{\sqrt{2}})
-q_{0,a-3}(-y)
\frac{e^{-\frac{y^2}{2}}}{\sqrt{2\pi}}
\Bigr]
\\
&\quad
\\
&=
-e^{u}\Bigl(\frac{a}{2}+ab \Bigr)
(a+1)\frac{\Gamma(a+1)}{\sqrt{2\pi}}e^{-\frac{y^2}{4}}D_{-2-a}(-y)
\\
&\quad
+e^{u}b (3a-1)
\frac{\Gamma(a+1)}{\sqrt{2\pi}}e^{-\frac{y^2}{4}}D_{-a}(-y)
-e^{u}\frac{5}{3}b
\frac{\Gamma(a+1)}{\sqrt{2\pi}}e^{-\frac{y^2}{4}}D_{2-a}(-y).
\end{align*}
Since by \cite[Eq. (12.8.1)]{NIST}, 
\begin{align*}
 (a+1)D_{-a-2}(y)&=D_{-a}(y)-yD_{-a-1}(y),\quad
 D_{2-a}(y)=(y^2+a-1)D_{-a}(y)+ayD_{-a-1}(y), \\
 (a+1)D_{-a-2}(-y)&=D_{-a}(-y)+yD_{-a-1}(-y),\quad
 D_{2-a}(-y)=(y^2+a-1)D_{-a}(-y)-ayD_{-a-1}(-y),
 \\
D_{1-a}(\xi)&=\xi D_{-a}(\xi)+aD_{-a-1}(\xi),\quad
D_{1-a}(-\xi)=-\xi D_{-a}(-\xi)+aD_{-a-1}(-\xi). 
\end{align*}
we have 
\begin{align*}
\frac{\widetilde{\mathcal{G}}_1(y)}{\widetilde{\mathcal{G}}_0(y)}
&=
\Bigl[
\frac{a}{2}+ab-b(3a-1)+\frac{5}{3}b(y^2+a-1)
\Bigr] \frac{D_{-a}(y)-e^{u}D_{-a}(-y)}{D_{-1-a}(y)+e^{u}D_{-a-1}(-y)}
+\Bigl(-\frac{a}{2}+\frac{2}{3}ab\Bigr)y,
\end{align*}
which gives 
\begin{align*}
&\quad\frac{1}{2}
\lim_{M\to+\infty}
\int_{-M}^{M}\frac{\mathcal{G}_{1}(-\tfrac{y}{\sqrt{2}}:u,a)}{\mathcal{G}_{0}(-\tfrac{y}{\sqrt{2}}:u,a)}\,dy
\\
&=
\frac{1}{2}
\lim_{M\to+\infty}
\int_{-M}^{M}
\Bigl[
\frac{a}{2}+ab-b(3a-1)+\frac{5}{3}b(y^2+a-1)
\Bigr] \frac{D_{-a}(-y)-e^{u}D_{-a}(y)}{D_{-1-a}(-y)+e^{u}D_{-a-1}(y)}
\,dy
\\
&\quad
+\frac{1}{2}
\lim_{M\to+\infty}
\Bigl(-\frac{a}{2}+\frac{2}{3}ab\Bigr)
\int_{-M}^{M}y\,dy
\\
&=
\lim_{M\to+\infty}
\int_{-M}^{M}
\Bigl[
\frac{a}{4}+\frac{1}{2}ab-\frac{1}{2}b(3a-1)+\frac{5}{6}b(y^2+a-1)
\Bigr]\partial_y\log\bigl[\mathcal{G}_0(\tfrac{y}{\sqrt{2}};u,a)\bigr]
\,dy
\\
&=
\lim_{M\to+\infty}
\Bigl[
\Bigl(
\frac{a}{4}+\frac{1}{2}ab-\frac{1}{2}b(3a-1)+\frac{5}{6}b(y^2+a-1)
\Bigr)
\log\bigl[\mathcal{G}_0(\tfrac{y}{\sqrt{2}};u,a)\bigr]\Bigr]\Bigr|_{y=-M}^{y=M}
\\
&\quad
-
\lim_{M\to+\infty}
\int_{-M}^{M}
\frac{10}{3}by
\log\bigl[\mathcal{G}_0(y;u,a)\bigr]
\,dy
\end{align*}
Recall 
\begin{align*}
&\int_{-\infty}^{+\infty}
\bigg[
\frac{1}{\sqrt{2}}\frac{\mathcal{G}_1(y:u,a)}{\mathcal{G}_0(y:u,a)}
+4by\Bigl(\log\bigl(\mathcal{G}_0(y:u,a)\bigr)-u\mathbf{1}_{(-\infty,0)}(y)\Bigr)
\\
&
\qquad
-\frac{a}{2}y\Bigl(1+2b+8b+8b\log(\sqrt{2}|y|)\Bigr)
+
\frac{2ab-a^2}{4}\frac{y}{1+y^2}
\bigg]],dy. 
\end{align*}
\end{comment}


\section{Euler-Maclaurin formula}
In this work, we have used Euler-Maclaurin formula \cite{BKS2023,AFLS25,ACC2023c} to establish the precise large $n$-asymptotics of \eqref{def of mathcal En}, which we state below for reference.
\begin{lemma}[Euler-Maclaurin formula]\label{lemma:Euler-Maclaurin}
Let $f(x)$ be $2m$ times differentiable function on the interval $[p,q]$. Then we have 
\begin{align}
\begin{split}
\sum_{j=p+1}^{q-1}f(j)&=\int_{p}^{q}f(x)\,dx-\frac{f(p)+f(q)}{2}+\sum_{j=1}^{m-1}\frac{B_{2j}}{(2j)!}\Bigl( 
f^{(2j-1)}(q)-f^{(2j-1)}(p)
\Bigr)
+R_{2m}, 
\\
\sum_{j=p}^{q}f(j)&=\int_{p}^{q}f(x)\,dx+\frac{f(p)+f(q)}{2}+\sum_{j=1}^{m-1}\frac{B_{2j}}{(2j)!}\Bigl( 
f^{(2j-1)}(q)-f^{(2j-1)}(p)
\Bigr)
+R_{2m},
\\
\sum_{j=p}^{q-1}f(j)
&=\int_{p}^{q}f(x)\,dx-\frac{f(q)-f(p)}{2}+\sum_{j=1}^{m-1}\frac{B_{2j}}{(2j)!}\Bigl(f^{(2j-1)}(q)-f^{(2j-1)}(p)\Bigr)+R_{2m}, 
\end{split}    
\end{align}
where the sequence $\{B_{2k}\}_k$ are even indexed Bernoulli numbers $B_{2}=\frac{1}{6}$, $B_4=-\frac{1}{30},\dots$, and the remainder term $R_{2m}$ satisfies the bound $|R_{2m}|\leq c_{2m}\int_{p}^{q}|f^{(2m)}(x)|\,dx$ with $c_{2m}=\frac{2\zeta(2m)}{(2\pi)^{2m}}$. 
Here, $\zeta(s)$ is the Riemann zeta function. 
\end{lemma}









































\begin{thebibliography}{99}

\bibitem{ABE2023} G. Akemann, S.-S. Byun and M. Ebke, \emph{Universality of the number variance in rotational invariant two-dimensional Coulomb gases}, \textit{J. Stat. Phys.} \textbf{190} (2023), no.1, Paper No. 9, 34 pp.

\bibitem{ABES2023} G. Akemann, S.-S. Byun, M. Ebke and G. Schehr, \emph{Universality in the number variance and counting statistics of the real and symplectic Ginibre ensemble}, \textit{J. Phys. A} \textbf{56}, doi:10.1088/1751-8121/ad0885.

\bibitem{ADM24}
G. Akemann, M. Duits, and L. D. Molag, \emph{Fluctuations in various regimes of non-Hermiticity and a holographic principle}, arXiv preprint. arXiv:2412.15854

%%\bibitem{AIS2014} G. Akemann, J.R. Ipsen and E. Strahov, Permanental processes from products of complex and quaternionic induced Ginibre ensembles,  \textit{Random Matrices Theory Appl.} \textbf{3} (2014), no. 4, 1450014, 54 pp.

%\bibitem{APS2009} G. Akemann, M.J. Phillips and L. Shifrin, Gap probabilities in non-Hermitian random matrix theory, \textit{J. Math. Phys.} \textbf{50} (2009), no. 6, 063504, 32 pp. 

%\bibitem{AS2013} G. Akemann and E. Strahov, Hole probabilities and overcrowding estimates for products of complex Gaussian matrices, \textit{J. Stat. Phys.} \textbf{151} (2013), no. 6, 987--1003.

\bibitem{AFLS25}
M. Allard, P. J. Forrester, S. Lahiry and B.-J. Shen, \emph{Partition function of 2D Coulomb gases with radially symmetric potentials and a hard wall}, arXiv:2506.14738.

\bibitem{AL25}
M. Allard and S. Lahiry,
\emph{Birth of a gap: Critical phenomena in 2D Coulomb gas},  arXiv:2509.24529.



\bibitem{ACC2023a} Y. Ameur, C. Charlier, and J. Cronvall, \emph{The two-dimensional Coulomb gas: fluctuations through a spectral gap}, arXiv:2210.13959.

\bibitem{ACC2023b} Y. Ameur, C. Charlier and J. Cronvall, \emph{Random normal matrices: eigenvalue correlations near a hard wall}, arXiv:2306.14166.

\bibitem{ACC2023c}
Y. Ameur, C. Charlier and J. Cronvall, \emph{Free Energy and Fluctuations in the Random Normal Matrix Model with Spectral Gaps}, Constr Approx (2025). https://doi.org/10.1007/s00365-025-09720-9


\bibitem{ACCL1} Y. Ameur, C. Charlier, J. Cronvall and J. Lenells, \emph{Exponential moments for disk counting statistics at the hard edge of random normal matrices}, \textit{J. Spectr. Theory} \textbf{13} (2023), no. 3, pp. 841--902.



\bibitem{ACCL2} Y. Ameur, C. Charlier, J. Cronvall and J. Lenells, \emph{Disk counting statistics near hard edges of random normal matrices: the multi-component regime}, Adv. Math. \textbf{441} (2024), 109549.

\bibitem{ACM2024}%NEW ADD
Y. Ameur, C. Charlier and P. Moreillon, \emph{Eigenvalues of truncated unitary matrices: disk counting statistics}, Monatsh. Math. {\bf 204} (2024), no.~2, 197--216; MR4745781

\bibitem{AKS2023} Y. Ameur, N.-G. Kang and S.-M. Seo, \emph{The random normal matrix model: insertion of a point charge}, \textit{Potential Anal.} \textbf{58} (2023), no.2, 331--372.

\bibitem{AW1984} R. Askey and J. Wimp, \emph{Associated Laguerre and Hermite polynomials}, \textit{Proc. Roy. Soc. Edinburgh Sect. A} \textbf{96} (1-2) (1984), 15--37.


\bibitem{BBLM2015} F. Balogh, M. Bertola, S.-Y. Lee and D.T.-R. McLaughlin, \emph{Strong asymptotics of the orthogonal polynomials with respect to a measure supported on the plane}, \textit{Comm. Pure Appl. Math.} \textbf{68} (2015), no.1, 112--172.

\bibitem{BGM2017} F. Balogh, T. Grava and D. Merzi, \emph{Orthogonal polynomials for a class of measures with discrete rotational symmetries in the complex plane}, \textit{Constr. Approx.} \textbf{46} (2017), no.1, 109--169.


\bibitem{Basor} E. Basor, \emph{Asymptotic formulas for Toeplitz determinants}, {\em Trans.
Amer. Math. Soc.} {\bf 239}  (1978), 33--65.

\bibitem{BasorEhrhardt1} E. Basor and T. Ehrhardt, \emph{Asymptotic formulas for determinants of a special class of Toeplitz + Hankel matrices. In \textit{Large truncated Toeplitz matrices, Toeplitz operators, and related topics}}, \textit{Oper. Theory Adv. Appl.}, \textbf{259}, Birkhäuser/Springer, Cham, 2017. 


%\bibitem{BasMor} E. Basor and K.E. Morrison, The Fisher-Hartwig conjecture and Toeplitz eigenvalues, \textit{Linear Algebra Appl.} \textbf{202} (1994), 129--142.

%\bibitem{BasTra} E. L. Basor and C. A. Tracy. The Fisher-Hartwig conjecture and generalizations. {\em Phys. A} {\bf 177} (1991), 167–173.

\bibitem{BerWebbWong} N. Berestycki, C. Webb and M.D. Wong, \emph{Random Hermitian Matrices and Gaussian Multiplicative Chaos}, \textit{Probab. Theory Related Fields}, \textbf{172} (2018), 103--189.



\bibitem{BEG18} M. Bertola, J.G. Elias Rebelo and T. Grava, \emph{Painlev\'{e} IV critical asymptotics for orthogonal polynomials in the complex plane}, \textit{SIGMA Symmetry Integrability Geom. Methods Appl.} \textbf{14} (2018), Paper No. 091, 34 pp.

%\bibitem{BS1985} A. B\"ottcher and B. Silbermann, Toeplitz matrices and determinants with Fisher-Hartwig symbols, \textit{J. Funct. Anal.} \textbf{63} (1985), no. 2, 178--214. 


%\bibitem{BGL2022} A.I. Bufetov, D. Garc\'{i}a-Zelada, and Z. Lin, Fluctuations of the process of moduli for the Ginibre and hyperbolic ensembles, arXiv:2202.11687.

%\bibitem{Berezin2023} S. Berezin, Functional Limit Theorems for Constrained Mittag-Leffler Ensemble in Hard Edge Scaling, arXiv:2308.12658.

\bibitem{BKP2023} S. Berezin, A.B.J. Kuijlaars and I. Parra, \emph{Planar orthogonal polynomials as type I multiple orthogonal polynomials}, \textit{SIGMA Symmetry Integrability Geom. Methods Appl.} \textbf{19} (2023), Paper No. 020, 18 pp.

\bibitem{BCL2024}
E.~J. Blackstone, C. Charlier and J. Lenells, \emph{Toeplitz determinants with a one-cut regular potential and Fisher-Hartwig singularities I. Equilibrium measure supported on the unit circle}, Proc. Roy. Soc. Edinburgh Sect. A {\bf 154} (2024), no.~5, 1431--1472; MR4806283

\bibitem{Byun} S.-S. Byun, \emph{Planar equilibrium measure problem in the quadratic fields with a point charge}, \textit{Comput. Methods Funct. Theory} (2023), \texttt{https://doi.org/10.1007/s40315-023-00494-4}.


\bibitem{Byun25} S.-S. Byun, \emph{Anomalous free energy expansions of planar Coulomb gases: multi-component and conformal singularity}, arXiv:2508.00316



\bibitem{BC2022} S.-S. Byun and C. Charlier, \emph{On the characteristic polynomial of the eigenvalue moduli of random normal matrices}, Constr Approx (2024). https://doi.org/10.1007/s00365-024-09689-x

\bibitem{BFreview} S.-S. Byun and P. J. Forrester, \emph{Progress on the study of the Ginibre ensembles}, KIAS Springer Ser. Math. \textbf{3} Springer, 2025, 221pp.


\bibitem{BFL25}
S.-S. Byun, P. J. Forrester and S. Lahiry, \emph{Properties of the one-component Coulomb gas on a sphere with two macroscopic external charges}, arXiv:2501.05061.



\bibitem{BFKL25}
S.-S. Byun, P. J. Forrester A. B. J. Kuijlaars and S. Lahiry, \emph{Orthogonal polynomials in the spherical ensemble with two insertions}, arXiv:2503.15732.

\bibitem{BKS2023} S.-S. Byun, N.-G. Kang and S.-M. Seo, \emph{Partition functions of determinantal and Pfaffian Coulomb gases with radially symmetric potentials}, \textit{Comm. Math. Phys.} \textbf{401} (2023), no.2, 1627--1663.



\bibitem{BKSY2025}
S.-S. Byun, N.-G. Kang, S.-M. Seo, M. Yang, \emph{Free energy of spherical Coulomb gases with point charges}, arXiv:2501.07284.

\bibitem{BLY2022}
S.-S. Byun, S.-Y. Lee and M. Yang, \emph{Lemniscate ensembles with spectral singularity}, arXiv:2107.07221v2.

\bibitem{BP24} 
 S.-S. Byun and S. Park, \emph{Large gap probabilities of complex and symplectic spherical ensembles with point charges}, arXiv:2405.00386.

\bibitem{BSY2025}
S.-S. Byun, S.-M. Seo and M. Yang, \emph{Free energy expansions of a conditional GinUE and large deviations of the smallest eigenvalue of the LUE}, Comm. Pure Appl. Math. (Online), https://doi.org/10.1002/cpa.70005.


\bibitem{BY2023}
S.-S. Byun and M. Yang, \emph{Determinantal Coulomb gas ensembles with a class of discrete rotational symmetric potentials}, SIAM J. Math. Anal. \textbf{55} (2023), 6867–-6897.

%\bibitem{CE2020} L. Charles and B. Estienne, Entanglement entropy and Berezin-Toeplitz operators, \textit{Comm. Math. Phys.} \textbf{376} (2020), no. 1, 521--554.

\bibitem{Charlier} C. Charlier, \emph{Asymptotics of Hankel determinants with a one-cut regular potential and Fisher-Hartwig singularities}, \textit{Int. Math. Res. Not. IMRN} \textbf{2019} (2019), 7515--7576.

%\bibitem{Charlier 2021a}%new add
%C. Charlier, \emph{Large gap asymptotics for the generating function of the sine point process}, Proc. Lond. Math. Soc. (3) {\bf 123} (2021), no.~2, 103--152; MR4307032



%\bibitem{Charlier 2021b}%new add
%C. Charlier, \emph{Exponential moments and piecewise thinning for the Bessel point process}, Int. Math. Res. Not. IMRN {\bf 2021}, no.~21, 16009--16073; MR4338214

%\bibitem{CharlierMB} C. Charlier, \emph{Asymptotics of Muttalib–Borodin determinants with Fisher–Hartwig singularities}, Selecta Mathematica, \textbf{28} (2022).

\bibitem{C2021 FH} C. Charlier, \emph{Asymptotics of determinants with a rotation-invariant weight and discontinuities along circles}, \textit{Adv. Math.} \textbf{408} (2022), Paper No. 108600, 36 pp.

\bibitem{C2021} C. Charlier, \emph{Large gap asymptotics on annuli in the random normal matrix model}, \textit{Math. Ann.} \textbf{388}(4), 3529--3587 (2024). \texttt{https://doi.org/10.1007/s00208-023-02603-z}.

\bibitem{C2023}
C. Charlier, Hole probabilities and balayage of measures for planar Coulomb gases, arXiv:2311.15285.

%\bibitem{CC2020}%new add
%C. Charlier and T. Claeys, Large gap asymptotics for Airy kernel determinants with discontinuities, Comm. Math. Phys. {\bf 375} (2020), no.~2, 1299--1339; MR4083883

\bibitem{CFWW2021} C. Charlier, B. Fahs, C. Webb and M.D. Wong, \emph{Asymptotics of Hankel determinants with a multi-cut regular potential and Fisher-Hartwig singularities}, Mem. Amer. Math. Soc. \textbf{310} (2025), 1567.


%\bibitem{CharlierGharakhloo} C. Charlier and R. Gharakhloo, \emph{Asymptotics of Hankel determinants with a Laguerre-type or Jacobi-type potential and Fisher-Hartwig singularities}, \textit{Adv. Math.} \textbf{383} (2021), 107672, 69 pp.


\bibitem{CM2023}%new add
C. Charlier and P. Moreillon, \emph{On the generating function of the Pearcey process}, Ann. Appl. Probab. {\bf 33} (2023), no.~4, 3240--3277; MR4612666

%\bibitem{C2023} C. Charlier, \emph{Hole probabilities and balayage of measures for planar Coulomb gases}, arXiv:2311.15285.

\bibitem{CL2023} C. Charlier and J. Lenells, \emph{Exponential moments for disk counting statistics of random normal matrices in the critical regime}, \textit{Nonlinearity} \textbf{36} (2023), no.3, 1593--1616.



%\bibitem{CGS2019} T. Claeys, M. Girotti and D. Stivigny, Large gap asymptotics at the hard edge for product random matrices and Muttalib-Borodin ensembles, \textit{Int. Math. Res. Not.} \textbf{2019} (2019), 2800--2847. 

\bibitem{CGMY2020} T. Claeys, G. Glesner, A. Minakov and M. Yang, \emph{Asymptotics for averages over classical orthogonal ensembles}, \textit{Int. Math. Res. Not. IMRN} \textbf{2022} (2022), 7922--7966.

\bibitem{CK2015} T. Claeys and I. Krasovsky, \emph{Toeplitz determinants with merging singularities}, \textit{Duke Math. J.} \textbf{164} (2015), no. 15, 2897--2987.


\bibitem{CMV2016} F.D. Cunden, F. Mezzadri and P. Vivo, \emph{Large deviations of radial statistics in the two-dimensional one-component plasma}, \textit{J. Stat. Phys.} \textbf{164} (2016), no. 5, 1062--1081.

\bibitem{DXZ2020} D. Dai, S.-X. Xu and L. Zhang, \emph{Asymptotics of Fredholm determinant associated with the Pearcey kernel}, \textit{Comm. Math. Phys.} \textbf{382} (2021), no. 3, 1769--1809.

\bibitem{DXZ2022}%new add
D. Dai, S.-X. Xu and L. Zhang, \emph{On the deformed Pearcey determinant}, Adv. Math. {\bf 400} (2022), Paper No. 108291, 64 pp.; MR4386547

\bibitem{DXZ2023}%new add
D. Dai, S.-X. Xu and L. Zhang, \emph{Gap probability for the hard edge Pearcey process}, Ann. Henri Poincar\'e{} {\bf 24} (2023), no.~6, 2067--2136; MR4586858

\bibitem{DeanoSimm} A. Dea\~{n}o and N. Simm, \emph{Characteristic polynomials of complex random matrices and Painlev\'{e} transcendents}, \textit{Int. Math. Res. Not. IMRN} \textbf{2022} (2022), 210--264.

\bibitem{DKMS25}
A. Dea\~{n}o, K. McLaughlin, L. Molag and N. Simm, \emph{Asymptotic analysis for a class of planar orthogonal polynomials on the unit disc}, arXiv:2505.12633.

\bibitem{DIZ1997} P. Deift, A. Its and X. Zhou, \emph{A Riemann-Hilbert approach to asymptotic problems arising in the theory of random matrix models, and also in the theory of integrable statistical mechanics}, \textit{Ann. of Math.} (2) \textbf{146} (1997), no. 1, 149--235. 


\bibitem{DIK} P. Deift, A. Its and I. Krasovky,
\emph{Asymptotics of Toeplitz, Hankel, and Toeplitz+Hankel determinants with Fisher-Hartwig singularities}, \textit{Ann. Math.} \textbf{174} (2011), 1243--1299.

\bibitem{DIKreview} P. Deift, A. Its, and I. Krasovsky, \emph{Toeplitz matrices and Toeplitz determinants under the impetus of the Ising model: some history and some recent results}, \textit{Comm. Pure Appl. Math.} \textbf{66} (2013), no. 9, 1360--1438. 

\bibitem{DIK14}
P. Deift, A. Its, I. Krasovsky, \emph{On the Asymptotics of a Toeplitz Determinant with Singularities}, MSRI Publications, vol. \textbf{65}, Cambridge University Press, 2014.


%\bibitem{Ehr2001} T. Ehrhardt, A status report on the asymptotic behavior of Toeplitz determinants with Fisher-Hartwig singularities. \textit{Oper. Theory Adv. Appl.} \textbf{124} (2001), 217--241.

%\bibitem{ElbauFedler} P. Elbau and G. Felder, Density of eigenvalues of random normal matrices, \textit{Comm. Math. Phys.}, \textbf{259} (2005), no. 2, 433--450.

%\bibitem{ES2020} B. Estienne and J.-M. St\'{e}phan, Entanglement spectroscopy of chiral edge modes in the quantum Hall effect, \textit{Physical Review B} \textbf{101} (2020), no. 11, 115136.

\bibitem{Fahs} B. Fahs, \emph{Uniform asymptotics of Toeplitz determinants with Fisher-Hartwig singularities}, \textit{Comm. Math. Phys.} \textbf{383} (2021), no. 2, 685--730.



\bibitem{FisherHartwig} M.E. Fisher and R.E. Hartwig, \emph{Toeplitz determinants: Some applications, theorems, and conjectures}, \textit{Advan. Chem. Phys.} \textbf{15} (1968), 333--353.


%\bibitem{ForKea} J. Forkel and J.P. Keating, The classical compact groups and Gaussian multiplicative chaos, \textit{Nonlinearity} \textbf{34} (2021), no. 9, 6050--6119.



%\bibitem{FyoKea} Y.V. Fyodorov and J.P. Keating, Freezing transitions and extreme values: random matrix theory, and disordered landscapes, \textit{Philos. Trans. R. Soc. Lond. Ser. A} \textbf{372} (2014), no. 2007, 20120503, 32 pp.

\bibitem{FenzlLambert} M. Fenzl and G. Lambert, \emph{Precise deviations for disk counting statistics of invariant determinantal processes}, \textit{Int. Math. Res. Not. IMRN} \textbf{2022} (2022), 7420--7494.


\bibitem{ForresterHoleProba} P.J. Forrester, \emph{Some statistical properties of the eigenvalues of complex random matrices}, \textit{Phys. Lett. A} \textbf{169} (1992), no. 1-2, 21--24. 

\bibitem{Forrester}
P. J. Forrester, \emph{Log-gases and Random Matrices} (LMS-34), Princeton University Press, Princeton, 2010.

\bibitem{Garoni} T.M. Garoni, \emph{On the asymptotics of some large Hankel determinants generated by Fisher-Hartwig symbols defined on the real line}, \textit{J. Math. Phys.} \textbf{46} (2005), 19 pp.

\bibitem{Ginibre} J. Ginibre, \emph{Statistical ensembles of complex, quaternion, and real matrices}, \textit{J. Mathematical Phys.} \textbf{6} (1965), 440--449. 


%\bibitem{Gould} H.W. Gould, Euler's formula for $n$th differences of powers, \textit{Amer. Math. Monthly} \textbf{85} (1978), no.6, 450--467.


\bibitem{GRtable} I.S. Gradshteyn and I.M. Ryzhik, \emph{Table of integrals, series, and products, Seventh edition}, Elsevier Academic Press, Amsterdam, 2007.


\bibitem{HM2013} H. Hedenmalm and N. Makarov, \textit{Coulomb gas ensembles and Laplacian growth}, \textit{Proc. Lond. Math. Soc.} (3) \textbf{106} (2013), no.4, 859--907.

\bibitem{ItsKrasovsky} A. Its and I. Krasovsky, \emph{Hankel determinant and orthogonal polynomials for the Gaussian weight with a jump}, \textit{Contemporary Mathematics} \textbf{458} (2008), 215--248.


\bibitem{Krasovsky} I. Krasovsky, \emph{Correlations of the characteristic polynomials in the Gaussian unitary ensemble or a singular Hankel determinant}, \textit{Duke Math J.} \textbf{139} (2007), 581--619. 


\bibitem{JLM1993} B. Jancovici, J. Lebowitz and G. Manificat, \emph{Large charge fluctuations in classical Coulomb systems}, \textit{J. Statist. Phys.} \textbf{72} (1993), no. 3-4, 773--787.



%\bibitem{KS2008} A. Karageorghis and Y.-S. Smyrlis, Conformal mapping for the efficient MFS solution of Dirichlet boundary value problems, \textit{Computing} \textbf{83} (2008), no.1, 1--24. % ref for ellipse




\bibitem{LMS2018} B. Lacroix-A-Chez-Toine, S.N. Majumdar and Gr\'{e}gory Schehr, \emph{Rotating trapped fermions in two dimensions and the complex Ginibre ensemble: Exact results for the entanglement entropy and number variance}, \textit{Phys. Rev. A} \textbf{99} (2019), 021602.

\bibitem{L et al 2019} B. Lacroix-A-Chez-Toine, J.A.M. Garz\'{o}n, C.S.H. Calva, I.P. Castillo, A. Kundu, S.N. Majumdar, and G. Schehr, \emph{Intermediate deviation regime for the full eigenvalue statistics in the complex Ginibre ensemble}, \textit{Phys. Rev. E} \textbf{100} (2019), 012137.



\bibitem{LeeRiser2016} S.-Y. Lee and R. Riser, \emph{Fine asymptotic behavior for eigenvalues of random normal matrices: ellipse case}, \textit{J. Math. Phys.} \textbf{57} (2016), no. 2, 023302, 29 pp. 

\bibitem{LeeYang} S.-Y. Lee and M. Yang, \emph{Discontinuity in the asymptotic behavior of planar orthogonal polynomials under a perturbation of the Gaussian weight}, \textit{Comm. Math. Phys.} \textbf{355} (2017), no. 1, 303--338. 

\bibitem{LeeYang2} S.-Y. Lee and M. Yang, Planar orthogonal polynomials as Type II multiple orthogonal polynomials, \textit{J. Phys. A} \textbf{52} (2019), no. 27, 275202, 14 pp. 

\bibitem{LY2023} S.-Y. Lee, M. Yang, \emph{Strong asymptotics of planar orthogonal polynomials: Gaussian weight perturbed by finite number of point charges}, \textit{Comm. Pure Appl. Math.} \textbf{76} (2023), no.10, 2888--2956.


\bibitem{LS2017} T. Lebl\'{e} and S. Serfaty, \emph{Large deviation principle for empirical fields of log and Riesz gases}, \textit{Invent. Math.} \textbf{210} (2017), no.3, 645--757.


\bibitem{Lenard} A. Lenard, \emph{Some remarks on large Toeplitz determinants}, \textit{Pacific J. Math.} \textbf{42} (1972), 137--145. 



\bibitem{LYZ24}%new add
J. Liu, L. Yao, and L. Zhang, \emph{Large gap asymptotics of the hard edge tacnode process}, arXiv:2412.12920

\bibitem{NIST} F.W.J. Olver, A.B. Olde Daalhuis, D.W. Lozier, B.I. Schneider, R.F. Boisvert, C.W. Clark, B.R. Miller and B.V. Saunders, NIST Digital Library of Mathematical Functions. http://dlmf.nist.gov/, Release 1.0.13 of 2016-09-16.

%\bibitem{RFbook} T. Ransford, \textit{Potential theory in the complex plane}, London Math. Soc. Stud. Texts \textbf{28}, Cambridge University Press, Cambridge, 1995. 

\bibitem{MMO25}
J. Marzo, L. D. Molag, and J. Ortega-Cerd\'{a}, \emph{Universality for fluctuations of counting statistics of random normal matrices}, arXiv:2508.04386


\bibitem{Mehta} M.L. Mehta, \emph{Random matrices}. \textit{Pure and Applied Mathematics} (Amsterdam), Vol. 142, 3rd ed., Elsevier/Academic Press, Amsterdam, 2004.

\bibitem{NAKP2020} T. Nagao, G. Akemann, M. Kieburg and I. Parra, \emph{Families of two-dimensional Coulomb gases on an ellipse: correlation functions and universality}, \textit{J. Phys. A} \textbf{53} (2020), no. 7, 075201, 36 pp.

\bibitem{SaTo} E. B. Saff and V. Totik, {\em Logarithmic Potentials with External Fields},  Grundlehren der Mathematischen Wissenschaften, Springer-Verlag, Berlin, 1997.

\bibitem{Seo} S.-M. Seo, \emph{Edge behavior of two-dimensional Coulomb gases near a hard wall}, \textit{Ann. Henri Poincar\'{e}} \textbf{23} (2022), no. 6, 2247--2275.



\bibitem{SDMS2020} N.R. Smith, P. Le Doussal, S.N. Majumdar and G. Schehr, \emph{Counting statistics for non-interacting fermions in a $d$-dimensional potential}, \textit{Phys. Rev. E} \textbf{103} (2021), L030105.

\bibitem{SDMS2021} N.R. Smith, P. Le Doussal, S.N. Majumdar and G. Schehr, \emph{Counting statistics for non-interacting fermions in a rotating trap}, \textit{Phys. Rev. A} \textbf{105} (2022), 043315

\bibitem{Temme} N.M. Temme, \textit{Special functions: An introduction to the classical functions of mathematical physics}, John Wiley \& Sons (1996).

\bibitem{Van Assche} W. Van Assche, \emph{Orthogonal polynomials, associated
polynomials and functions of the second kind}, \textit{J. Comput. Appl. Math.} \textbf{37} (1991), 237--249.

%\bibitem{Webb} C. Webb, The characteristic polynomial of a random unitary matrix and Gaussian multiplicative chaos--the $L^{2}$-phase, \textit{Electron. J. Probab.} \textbf{20} (2015), no. 104, 21 pp. 

\bibitem{WebbWong} C. Webb and M.D. Wong, \emph{On the moments of the characteristic polynomial of a Ginibre random matrix}, \textit{Proc. Lond. Math. Soc.} (3) \textbf{118} (2019), no. 5, 1017--1056. 



\bibitem{Widom2} H. Widom. \emph{Toeplitz determinants with singular generating functions}, {\em Amer. J. Math.} {\bf 95} (1973), 333--383.

\bibitem{Wunsche} A. W\"{u}nsche, \emph{Generalized Hermite polynomials associated with functions of parabolic cylinder}, \textit{Appl. Math. Comput.} \textbf{141} (2003) 197–213.

\bibitem{Xetal}%new add
J. Xia et al., \emph{Asymptotics of the deformed higher order Airy-kernel determinants and applications}, Nonlinearity {\bf 36} (2023), no.~8, 4384--4424; MR4619236

\bibitem{XZZ2024}%new add
S.-X. Xu, S.-Q. Zhao and Y.-Q. Zhao, \emph{On the Fredholm determinant of the confluent hypergeometric kernel with discontinuities}, Phys. D {\bf 461} (2024), Paper No. 134101, 23 pp.; MR4711936

\bibitem{YZ24}%new add
L. Yao and L. Zhang, \emph{On the gap probability of the tacnode process}, Adv. Math. {\bf 438} (2024), Paper No. 109474, 85 pp.; MR4687908

\bibitem{ZW2006}
A. Zabrodin and P. Wiegmannm, \emph{Large-N expansion for the 2D Dyson gas}, J. Phys. A, \textbf{39}(28):8933--8963, 2006.



\end{thebibliography}




\end{document}

